% T12 C/D, chapitre 4 : Nombres complexes. Source : le livre « Terminales scientifiques » (Complexe/Lesona, Complexe/Fiche, Complexe/Exo, Complexe/Valiny).
% Les exercices tirés d’autres ouvrages ne sont pas repris (décision du 2 novembre 2026).

\TCDLessonStyleOn
\newcommand{\Reel}[1]{\mathfrak{Re}(#1)}
\newcommand{\Imag}[1]{\mathfrak{Im}(#1)}

\begin{chapter}{Ensemble des nombres complexes}
\section{Nombres complexes}
Dans l'ensemble des nombres r\'eels, certaines \'equations polynomiales telles $x^2 + 1 = 0$ n'admettent pas de solutions. 
Pour rem\'edier \`a ce probl\`eme, on introduit un nouveau \emph{nombre imaginaire} $i$ v\'erifiant $i^2 = -1$.
On obtient ainsi un nouvel ensemble qu'on d\'enote usuellement par $\mathbb{C}$. 
Cet ensemble $\mathbb{C}$ est appel\'e le corps des nombres complexes. 
Il est constitu\'e des nombres de la forme $a + i b$ o\`u $a$ et $b$ sont deux nombres r\'eels et $i$ l'unit\'e imaginaire.
%
\subsection{Forme alg\'ebrique d'un nombre complexe}
%
\index{Nombres complexes!Forme algebrique@Forme alg\'ebrique}
%
Si $z$ est un \'el\'ement de $\mathbb{C}$, alors $z$ s'\'ecrit de mani\`ere unique sous la forme \defbox{$a + i b$ \label{eq:forme alg}}.
Cette forme est appel\'ee la \emph{forme alg\'ebrique} de $z$.
Le r\'eel $a$ est appel\'e \emph{partie r\'eelle de $z$} et est not\'e $\Reel{z}$.
Le r\'eel $b$ est appel\'e \emph{partie imaginaire de $z$} et est not\'e $\Imag{z}$. \\
Un nombre complexe dont la partie r\'eelle est nulle est qualifi\'e d'\emph{imaginaire pur}.

\subsection{Op\'erations dans $\mathbb{C}$}
%
\index{Nombres complexes!operations@op\'erations}
%
Soient $z_1 = a_1 + i b_1$ et $z_2 = a_2 + i b_2$ deux nombres complexes.
\begin{itemize}
\item $z_1 = z_2 $ si et seulement si $(a_1, b_1) = (a_2, b_2)$,
\item $z_1 + z_2 = (a_1 + i b_1) + (a_2 + i b_2) = (a_1 + a_2 ) + i (b_1 + b_2) $,
\item $z_1 \cdot z_2 = (a_1 + i b_1) \cdot (a_2 + i b_2) = a_1 a_2 + i a_1 b_2 + i b_1 a_2 + i^2 b_1 b_2$.
      Utilisant le fait que $i^2 = -1$ on a:
      $z_1 \cdot z_2 = (a_1 a_2 - b_1 b_2) + i (a_1 b_2 + b_1 a_2)$.
\end{itemize}
En pratique, toutes les op\'erations se font de la m\^eme fa\c{c}on que dans $\mathbb{R}$ (on peut m\^eme diviser par tout nombre complexe non nul).

\begin{exe}
Mettre les nombres complexes suivants sous la forme alg\'ebrique : \\
$z_1 = (1-i)^2 (3 + 2i)$, $z_2 = \frac{10-20i}{1-3i}$ et $z_3 = \frac{(10+15i)(10-20i)}{(3+i)^2} $. \\
Dans chaque cas, donner la partie r\'eelle et la partie imaginaire.
\label{exe:mettre sous forme algebrique}
\end{exe}

\subsection{Conjugu\'e d'un nombre complexe}
\index{Nombres complexes!conjugue@conjugu\'e}
%
\begin{defi}
Pour un nombre complexe $z = a + i b$ donn\'e, on d\'efinit \emph{le conjugu\'e} de $z$ comme \'etant le nombre complexe
\defbox{$\bar{z} = a - i b$}.
\end{defi}
%
\begin{prop}
      \label{prp:proprietes du conjugue}
      %
	Si $z$, $z_1$ et $z_2$ sont trois nombres complexes, on a : 
	%
	\renewcommand{\bar}{\overline}
	%
	\begin{enumerate}
		\item $\bar{\bar{z}} = z$,
		\item $\bar{z_1 + z_2} = \bar{z_1} + \bar{z_2}$,
		\item $\bar{z_1 \cdot z_2} = \bar{z_1} \cdot \bar{z_2}$,
		      \label{prp:proprietes du conjugue, multiplication}
		\item $\bar{\frac{1}{z}} = \frac{1}{\bar{z}}$,
		\item $\bar{z^n} = \bar{z}^n$.
	\end{enumerate}
\end{prop}
%
% La preuve est laiss\'e \`a l'\'etudiant \`a titre d'exercices.
%
\subsection{Module d'un nombre complexe}
%
\index{Nombres complexes!module}
%
\begin{defi}
Soit $z = a + i b$ un nombre complexe.
On appelle \emph{module de $z$} le nombre r\'eel not\'e $| z |$ et d\'efini par 
\defbox{$|z| = \sqrt{z \bar{z}} = \sqrt{a^2 + b^2}$}.
\end{defi}
%
\begin{prop}
	\label{prp:proprietes du module}
	Soient $z = a + i b$ et $z'$ deux nombres complexes. On a:
	\begin{enumerate}
		\item $|z| = |\bar{z}|$,
		\item $|z| = 0 \iff z = 0$,
		\item $|z z'| = |z| |z'|$,
		\item $|\frac{1}{z}| = \frac{1}{|z|}$,
		\item $|z+z'| \leq |z| + |z'|$.
		      \label{prp:proprietes du module, inegalite triangulaire}
	\end{enumerate}
\end{prop}
%
\begin{exe}
\label{exe:forme algebrique et modules}
Ecrire les nombres complexes suivants sous forme alg\'ebrique et donner leurs modules:
$ z_1 = \overline{- (1 - 2 i)^2} $ et
$ z_2 = \overline{\left( \frac{1}{3 + 4i} \right) }$.

\end{exe}
%
\section{Point de vue g\'eometrique des nombres complexes}
\index{Nombres complexes!Representation geometrique@Repr\'esentation g\'eometrique}
\subsection{Repr\'esentation dans le plan complexe}
%%%%%%%%%%%%%%%%%%%%%%%%%%%%%%%%%%%%%%%%%%%%%%%%%%%%%%%%%%%%%%%%%%%%%%%%%%%%%%%%%%%%%%%%%%%%%%%%%%%%%%%%%%%%%%%%%%%%%%%%%%%%%%
%\begin{wrapfigure}{l}{0.5\actualtextwidth}
\begin{figure}
      \centering
		\begin{tikzpicture}[line cap=round,line join=round,>=triangle 45,x=1.0cm,y=1.0cm]
			\clip (-1,-1) rectangle (5,5);
			\draw[->,color=black] (-1,0) -- (5,0);
			\foreach \x in {-1,1,2,3,4}
			\draw[shift={(\x,0)},color=black] (0pt,-2pt);
			\draw[color=black] (4.,0) node [anchor=north] { Axe r\'eel};
			\draw[->,color=black] (0,-1) -- (0,5);
			\foreach \y in {-1,-1,1,2,3,4}
			\draw[color=black] (0.06,4.77) node [anchor=west] { Axe imaginaire};
			\clip(-3,-2) rectangle (5,5);
			\draw [->] (0,0) -- (0,1);
			\draw [->] (0,0) -- (1,0);
			\draw [->] (0,0) -- (3.22,2.22);
			\draw [dashed,loko2] (3.22,2.22)-- (3.21,0);
			\draw [dashed,loko1] (3.22,2.22)-- (0,2.19);
			\draw [color=black] (0,0) circle (1.5pt);
			\draw[color=black] (-0,-0) node[below left] {$O$};
			\draw[color=black] (-0.0,0.5) node[left] {$\vec{v}$};
			\draw[color=black] (0.5,-0.0) node[below] {$\vec{u}$};
			\fill [color=black] (3.22,2.22) circle (1.5pt);
			\draw[color=black] (3.8,2.5) node[above] {\parbox{6em}{$M (a,b)$ \\ $z = a+ib$}};
			\fill [color=black] (3.21,0) circle (0.5pt);
			\draw[loko2, very thick] (0,0) -- (3.23,0);
			\draw[color=loko2] (3.35,0.) node[above right] {$a$};
			\fill [color=black] (0,2.19) circle (0.5pt);
			\draw[loko1, very thick] (0,0) -- (0,2.2);
			\draw[color=loko1] (0,2.3) node[above right] {$b$};
		\end{tikzpicture}
	\label{fig:conjugue geom}
	\caption{Repr\'esentation d'un nombre dans le plan complexe.}
%\end{wrapfigure}
\end{figure}
%%%%%%%%%%%%%%%%%%%%%%%%%%%%%%%%%%%%%%%%%%
L'ensemble $\mathbb{C}$ des nombres complexes peut ais\'ement \^etre identifi\'e au plan $\mathbb{R}^2$.
\`A chaque nombre complexe $z = a + i b$, on associe le point $M(a,b)$ de $\mathbb{R}^2$.
On pourra ainsi parler de \emph{plan complexe}, en r\'ef\'erence \`a cette repr\'esentation dans un rep\`ere orthonorm\'e $(O,\vec{u},\vec{v})$.
\paragraph{Terminologie}
Avec cette repr\'esentation, on dit que:
\begin{itemize}
\item $M$ est l'image de $z$,
\item $z$ est l'affixe du point $M$,
\item l'axe des abscisses est l'axe r\'eel,
\item l'axe des ordonn\'ees est l'axe imaginaire.
\end{itemize}
%
De m\^eme que pour un point, on d\'efinit l'affixe d'un vecteur $\vec{u} \left( \begin{matrix} a \\ b \end{matrix} \right)$
comme \'etant le nombre complexe $z = a + i b$.
%
%%%%%%%%%%%%%%%%%%%%%%%%%%%%%%%%%%%%%%%%%%
%\end{minipage}
%%%%%%%%%%%%%%%%%%%%%%%%%%%%%%%%%%%%%%%%%%%%%%%%%%%%%%%%%%%%%%%%%%%%%%%%%%%%%%%%%%%%%%%%%%%%%%%%%%%%%%%%%%%%%%%%%%%%%%%%%%%%%%
Le module du nombre complexe $z$ prend alors sens comme \'etant la distance Euclidienne entre l'origine et son image: le point $M$.
Autrement dit, il s'agit de la norme du vecteur $\overrightarrow{OM}$. 

\subsection{Interpr\'etation g\'eometrique de la somme} 
%
\index{Nombres complexes!Representation geometrique@Repr\'esentation g\'eometrique!somme}
%
L'addition de deux nombres complexes se traduit alors comme l'addition de deux vecteurs du plan. (Voir Fig \ref{fig:complex_sum}) \\

\begin{figure}
\centering
\begin{tikzpicture}[line cap=round,line join=round,>=triangle 45,scale=.75]
\draw[->,color=black] (-1.51,0) -- (9.07,0);
\foreach \x in {-1,1,2,3,4,5,6,7,8,9}
\draw[shift={(\x,0)},color=black] (0pt,-2pt);
\draw[color=black] (8.11,0.06) node [anchor=south] { {Axe r\'eel}};
\draw[->,color=black] (0,-1.45) -- (0,6.91);
\foreach \y in {-1,1,2,3,4,5,6}
\draw[color=black] (0.08,6.6) node [anchor=west] { {Axe imaginaire}};
\clip(-1.51,-1.45) rectangle (9.07,6.91);
\draw [->] (0,0) -- (0,1);
\draw [->] (0,0) -- (1,0);
\draw [->,loko2] (0,0) -- (4.56,0.9);
\draw [->,loko1] (0,0) -- (1.21,4.26);
\draw [->] (0,0) -- (5.77,5.16);
\draw [dashed,loko2] (1.21,4.26)-- (5.77,5.16);
\draw [dashed,loko1] (4.56,0.9)-- (5.77,5.16);
\draw[color=black] (0,0.5) node[left] {$\vec{v}$};
\draw[color=black] (.5,0) node[below] {$\vec{u}$};
\begin{scriptsize}
\draw [color=black] (0,0) circle (1.5pt);
\draw[color=black] (-0.25,-0.2) node {$O$};
\fill [color=black] (4.56,0.9) circle (1.5pt);
\draw[color=loko2] (5.3,1.2) node {$M (a,b)$};
\draw[color=loko2] (5.5,0.8) node {$z = a+ib$};
\fill [color=black] (1.21,4.26) circle (1.5pt);
\draw[color=loko1] (2.,4.1) node {$M' (a',b')$};
\draw[color=loko1] (2.,3.8) node {$z' = a' + i b'$};
\fill [color=black] (5.77,5.16) circle (1.5pt);
\draw[color=black] (7.1,5.5) node[above] {$M'' (a + a' , b + b')$};
\draw[color=black] (7.5,5.1) node[above] {$z'' = (a+a') + i (b+b')$};
\end{scriptsize}
\end{tikzpicture}
\caption{Repr\'esentation de la somme de deux nombres complexes en tant que somme de deux vecteurs du plan complexe}
\label{fig:complex_sum}
\end{figure}

Soient $z = a+ib$ et $z' = a'+ib'$ deux nombres complexes. 
Soient $M (a,b)$ et $M'(a',b')$ leurs images respectifs.
Consid\'erons les vecteurs $\overrightarrow{OM}$ et $\overrightarrow{OM'}$, soit $M''$ le point du plan complexe tel que 
$
\overrightarrow{OM''} = \overrightarrow{OM} + \overrightarrow{OM'}
$.
On sait que $\overrightarrow{OM''}$ a pour coordonn\'ees $\left( \begin{matrix} a+a' \\ b+b' \end{matrix} \right)$.
On v\'erifie qu'il s'agit de l'image de $z'' = z + z'$.
%%%%%%%%%%%%%%%%%%%%%%%%%%%%%%%%%%%%%%%%%%
%\end{minipage}
%%%%%%%%%%%%%%%%%%%%%%%%%%%%%%%%%%%%%%%%%%%%%%%%%%%%%%%%%%%%%%%%%%%%%%%%%%%%%%%%%%%%%%%%%%%%%%%%%%%%%%%%%%%%%%%%%%%%%%%%%%%%%%
\subsection{Interpr\'etation g\'eometrique du conjugu\'e}
%
\index{Nombres complexes!Representation geometrique@Repr\'esentation g\'eometrique!conjugue@conjugu\'e}
%
\begin{figure}
\centering
\begin{tikzpicture}[line cap=round,line join=round,>=triangle 45,scale=.9]
\draw[->,color=black] (-1,0) -- (4,0);
\foreach \x in {-1,1,2,4}
\draw[color=black] (4,0) node [anchor=south] {$x$};
\draw[->,color=black] (0,-3) -- (0,3);
\foreach \y in {-3,-2,-1,1,2}
\clip(-1,-3) rectangle (4,3);
\draw[] (1.85,0.3) -- (2.15,0.3) -- (2.15,0) -- (1.85,0) -- cycle; 
\draw [gray,->] (0,0) -- (0,1);
\draw [gray,->] (0,0) -- (1,0);
\draw (1.85,2.48)-- (1.85,-2.48);
% Itony ilay hahitana hoe mitovy halava ilay tsipika roa
\draw[color=black] (1.85,1.42) node {$ \smallsetminus $};
\draw[color=black] (1.85,1.5) node {$ \smallsetminus $};
\draw[color=black] (1.85,-1.4) node {$ \smallsetminus $};
\draw[color=black] (1.85,-1.5) node {$ \smallsetminus $};
% Izay ihany 
\draw [] (0,0) circle (1.5pt) node[below left] {$O$};
\draw[color=black] (0.25,0.5) node {$\vec{v}$};
\draw[color=black] (0.5,-0.25) node {$\vec{u}$};
\fill [color=loko2] (1.85,2.48) circle (1.5pt);
\draw[color=loko2] (2.75,2.75) node {$M(a,b)$};
\draw[color=loko2] (3,2.35) node {$z = a+ib$};
\fill [color=loko2] (1.85,-2.48) circle (1.5pt);
\draw[color=loko2] (3,-2.75) node {$\overline{M}(a,-b)$};
\draw[color=loko2] (3,-2.25) node {$\bar{z} = a - i b$};
\end{tikzpicture}
\caption{Repr\'esentation du conjugu\'e d'un nombre complexe}
\label{fig:conjugue}
\end{figure}

Si $z = a + i b$ alors $\bar{z} = a - i b$; ainsi, l'image de $z$ a pour coordonn\'ees $(a,b)$ et celle de $\bar{z}$ a pour coordonn\'ees $(a,-b)$.
On constate alors que l'image de $\bar{z}$ n'est autre que l'image du point $M$ par la sym\'etrie d'axe $O\vec{u}$ (Voir Fig \ref{fig:conjugue}).

\section{Forme trigonom\'etrique d'un nombre complexe}
\subsection{Argument d'un nombre complexe}
%
\index{Nombres complexes!argument}
%
Dans le plan complexe muni d'un rep\`ere orthonorm\'e $(O,\vec{u},\vec{v})$, consid\'erons le point $M(x,y)$ d'affixe $z$ non nul\footnote{$0$ n'a pas d'argument}.
On appelle \emph{argument de $z$}, et on note par $\arg(z)$, la mesure de l'angle form\'e par $\overrightarrow{OM}$ avec $\vec{u}$.

Pour tout nombre complexe $z$ diff\'erent de $0$, on peut d\'efinir le nombre $z' = \frac{z}{|z|}$ qui a pour module $1$ \footnote{En effet, $\left|\frac{z}{|z|}\right| = \frac{|z|}{|z|} = 1$.}.
Ainsi, la distance entre l'origine et l'image $M'$ de $z'$ est \'egale \`a $1$ ; autrement dit, $M'$ appartient au cercle unit\'e.
On sait alors qu'il existe un angle $\theta$ tel que $\Reel{z'} = \cos \theta $ et $\Imag{z'} = \sin \theta$.
L'angle $\theta$ ainsi d\'efini est un argument de $z'$.
Puisque $\overrightarrow{OM} = |z| \overrightarrow{OM'}$, les deux vecteurs sont colin\'eaires ; il s'ensuit que l'angle form\'e par $\overrightarrow{OM}$ avec $\vec{u}$ est le m\^eme que celui form\'e par $\overrightarrow{OM'}$ avec $\vec{u}$ (\`a un multiple de $2\pi$ pr\`es).

\begin{figure}[h]
\centering
\begin{tikzpicture}[line cap=round,line join=round,>=triangle 45,x=3.0cm,y=3.0cm]
\draw[->,color=black] (-1.25,0) -- (2.5,0);
\foreach \x in {-1,-0.5,0.5,1,1.5,2}
\draw[->,color=black] (0,-1.25) -- (0,1.5);
\foreach \y in {-1,-0.5,0.5,1}
\clip(-1.25,-1.25) rectangle (2.5,1.5);
\draw [shift={(0,0)},line width=1.2pt,color=loko3,fill=loko3,fill opacity=0.1] (0,0) -- (0:0.24) arc (0:39.73:0.24) -- cycle;
\draw(0,0) circle (3cm);
\draw (0,0)-- (1.26,1.05);
\draw [dashed, color=loko1] (0,0.64)-- (0.77,0.64);
\draw [dashed, color=loko2] (0.77,0.64)-- (0.77,0);
\fill [color=black] (0,0) circle (1.5pt);
\draw[color=black] (-0.1,-0.1) node {$O$};
\fill [color=black] (1.26,1.05) circle (1.5pt);
\draw[color=black] (1.7,1.1) node {$M \text{ d'affixe } z$};
\fill [color=black] (0.77,0.64) circle (1.5pt);
\draw[color=black] (1.25,0.6) node {$M' \text{ d'affixe } \frac{z}{|z|}$};
\draw[color=loko3] (0.3,0.1) node {$\theta$};
\fill [color=black] (0.77,0) circle (1.5pt);
\draw [color=loko2,very thick] (0,0) -- (0.75,0);
\draw[color=loko2] (0.75,-0.1) node {$\cos \theta$};
\fill [color=black] (0,0.64) circle (1.5pt);
\draw [color=loko1,very thick] (0,0) -- (0,0.64);
\draw[color=loko1] (-0.15,0.7) node {$\sin \theta$};
\end{tikzpicture}
\caption{Repr\'esentation d'un nombre complexe}
\label{fig:complex_trigo}
\end{figure}

On constate que le point $M'$ a pour affixe $\cos \theta + i \sin \theta$.
De l\`a, on aboutit \`a la conclusion que $M$ a pour affixe $|z|(\cos \theta + i\sin \theta)$.
Si on pose $r = |z|$ alors on a \defbox{$z = r (\cos \theta + i \sin \theta)$ \label{eq:forme trigo}}.
Cette expression est la \emph{forme trigonom\'etrique} de $z$.
%
\index{Nombres complexes!Forme trigonometrique@Forme trigonom\'etrique}
%
Le m\^eme nombre complexe $z$ peut aussi \^etre d\'esign\'e par le couple \defbox{$[r,\theta]$ \label{eq:forme pol}}.
Cette derni\`ere expression est appel\'ee la \emph{forme polaire} du nombre complexe $z$.

\begin{exe}
	\label{exe:trigo-alg}
	\begin{enumerate}
				\item D\'eterminer la forme trigonom\'etrique, le module et un argument
				      de chacun des nombres complexes suivants :
	            		\[
	            		\begin{gathered}
	            		z_1=1,\quad z_2=i,\\
	            		z_3=2+2i,\quad z_4=1+\sqrt{3}i.
	            		\end{gathered}
	            		\]
		\item Pour chacun des nombres complexes suivants, 
		      donner le module et l'argument puis mettre sous forme alg\'ebrique :
		      \[ z'_1 = 2 \left(\cos \frac{\pi}{3} - \sin \frac{\pi}{3} \ i \right) \ , \quad 
		      z'_2 = \sqrt{3} \left(\cos \frac{\pi}{6} + \sin \frac{\pi}{6} \ i \right) \;. \]
	\end{enumerate}
\end{exe}
%

\begin{rmk}
Pour passer de la forme trigonom\'etrique \`a la forme alg\'ebrique, il suffit de developper l'expression. 
En effet, si $z = r (\cos \theta + i \sin \theta)$ alors 
la partie r\'eelle de $z$ n'est autre que $r \cos \theta$ et la partie imaginaire est $r \sin \theta$. 
L'expression alg\'ebrique est alors $z = r \cos \theta + r \sin \theta i$. \\
Pour passer de la forme alg\'ebrique \`a la forme trigonom\'etrique, on peut d'abord commencer par normaliser le nombre. 
C'est-\`a-dire qu'on le divise par son module. 
On sait que le nombre $\frac{z}{|z|}$ appartient au cercle unit\'e.
Il existe alors un angle $\theta$ v\'erifiant $\frac{z}{|z|} = \cos \theta + i \sin \theta$.
Finalement, l'expression trigonom\'etrique sera $z = |z| (\cos \theta + i \sin \theta)$.
\end{rmk}

\subsection{Quelques propri\'et\'es}
\begin{prop}
	\label{prop:argument}
	Soient $z$ et $z'$ deux nombres complexes non nuls. On a alors :
	\begin{itemize}
		\item $\arg(z z') = \arg(z) + \arg(z')$ et $\arg(z^n) = n \, \arg(z)$,
		\item $\arg \left(\frac{1}{z} \right) = - \arg(z)$ 
		      et $\arg\left(\frac{z}{z'}\right) = \arg(z) - \arg(z')$,
		\item $\arg(\bar{z}) = - \arg(z)$.
		\item Si $z \in \mathbb{R}^*$ alors 
		      $\arg(z) = 
		            \begin{multiform}
		                  0 \text{ si } z > 0 \\
		                  \pi \text{ si } z < 0 \;.
		            \end{multiform}
		      $
		\item Si $z = \lambda i$ avec $\lambda \in \mathbb{R}^*$ alors 
		      $\arg(z) = 
		            \begin{multiform}
		                  \frac{\pi}{2} \text{ si } \lambda > 0 \\
		                  \frac{3\pi}{2} \text{ si } \lambda < 0 \;.
		            \end{multiform}
		      $
	\end{itemize}
	Les \'egalit\'es portant sur des arguments s'entendent \`a un multiple de $2\pi$ pr\`es.
\end{prop}

\begin{thm}[Formule de MOIVRE]
	%
	\index{Moivre!Formule de}
	%
	\label{thm:Moivre}
	Soit $n$ un nombre entier relatif. On a
	$$ (\cos \theta + i \sin \theta)^n = \cos n \theta + i \sin n \theta .$$
\end{thm}

\section{Notation exponentielle d'un nombre complexe}
%
\index{Nombres complexes!Forme exponentielle}
%
%%%%%%%%%%%%%%%%%%%%%%%%%%%%%%%%%%%%%%%%%%%%%%%%%%%%%%%%%%%%%%%%%%%%%%%%%%%%%%%%%%%%%%%%%%%%%%%%%%%%%%%%%%%%%%%%%%%%%%%%%%%%%%
\begin{thm}[Formules d'Euler]
\label{thm:Expo-trigo}
Pour tout nombre r\'eel $\theta$ on a
\thmbox{$$ \expo{i \, \theta} = \cos \theta + i \sin \theta .$$}

A partir de cette formule, on peut en obtenir d'autres: \\
%%%%%%%%%%%%%%%%%%%%%%%%%%%%%%%%%%%%%%%%%%%%%%%%%%%%%%%%%%%%%%%%%%%%%%%%%%%%%%%%%%%%%%%%%%%%%%%%%%%%%%%%%%%%%%%%%%%%%%%%%%%%%%
\begin{minipage}{0.45 \actualtextwidth}
%%%%%%%%%%%%%%%%%%%%%%%%%%%%%%%%%%%%%%%%%%
\begin{align*}
\expo{i \, \theta }   & = \cos \theta + i \sin \theta \\
\expo{- i \, \theta } & = \cos \theta - i \sin \theta \\
\intertext{et on additionne pour avoir}
\expo{i \, \theta } + \expo{ -i \, \theta } & = 2 \cos \theta \\
\end{align*}
On a alors: \\
\thmbox{$$ \cos \theta  = \frac{\expo{i \, \theta } + \expo{ -i \, \theta }}{2} .$$}
%%%%%%%%%%%%%%%%%%%%%%%%%%%%%%%%%%%%%%%%%%
\end{minipage}
%%%%%%%%%%%%%%%%%%%%%%%%%%%%%%%%%%%%%%%%%%%%%%%%%%%%%%%%%%%%%%%%%%%%%%%%%%%%%%%%%%%%
\begin{minipage}[c][][c]{.02\actualtextwidth}
      \rule{.4pt}{10\baselineskip}
\end{minipage}
%%%%%%%%%%%%%%%%%%%%%%%%%%%%%%%%%%%%%%%%%%%%%%%%%%%%%%%%%%%%%%%%%%%%%%%%%%%%%%%%%%%%
\begin{minipage}{0.45 \actualtextwidth}
\begin{align*}
\expo{i \, \theta }   & = \cos \theta + i \sin \theta \\
\expo{- i \, \theta } & = \cos \theta - i \sin \theta \\
\intertext{et on soustrait pour avoir}
\expo{i \, \theta } - \expo{ -i \, \theta } & = 2 \, i \sin \theta \\
\end{align*}
et \\
\thmbox{$$ \sin \theta = \frac{\expo{i \, \theta } - \expo{ -i \, \theta }}{2 \, i} .$$}
%%%%%%%%%%%%%%%%%%%%%%%%%%%%%%%%%%%%%%%%%%
\end{minipage}
%%%%%%%%%%%%%%%%%%%%%%%%%%%%%%%%%%%%%%%%%%%%%%%%%%%%%%%%%%%%%%%%%%%%%%%%%%%%%%%%%%%%%%%%%%%%%%%%%%%%%%%%%%%%%%%%%%%%%%%%%%%%%%
\end{thm}
%%%%%%%%%%%%%%%%%%%%%%%%%%%%%%%%%%%%%%%%%%%%%%%%%%%%%%%%%%%%%%%%%%%%%%%%%%%%%%%%%%%%%%%%%%%%%%%%%%%%%%%%%%%%%%%%%%%%%%%%%%%%%%
%%%%%%%%%%%%%%%%%%%%%%%%%%%%%%%%%%%%%%%%%%%%%%%%%%%%%%%%%%%%%%%%%%%%%%%%%%%%%%%%%%%%%%%%%%%%%%%%%%%%%%%%%%%%%%%%%%%%%%%%%%%%%%

\section{Lin\'earisation de fonctions trigonom\'etriques}
%
\index{Fonctions trigonometriques@Fonctions trigonom\'etriques!linearisation@lin\'earisation}
%
Les formules vues dans la section pr\'ec\'edente peuvent servir \`a lin\'eariser les fonctions trigonom\'etriques.
En effet, pour lin\'eariser les expressions de la forme $\cos ^n \theta$, on utilise la formule exponentielle de la fonction cosinus.
L'expression $\cos ^n \theta$ devient alors $\left( \frac{ \expo{i \theta} + \expo{-i \theta} }{ 2 } \right)^n$ et on peut utiliser la formule du bin\^ome pour la d\'evelopper.
Puis on regroupe les termes et on les identifie \`a des fonctions trigonom\'etriques connues. \\
Voyons cela par un exemple :
\begin{ex}
	      Lin\'earisons l'expression $\cos ^3 \theta$.
	      Pour cela, on utilise la formule d'Euler correspondante :
      \begin{align*}
            \cos ^3 \theta & = \left( \frac{\expo{i \theta } + \expo{- i \theta }}{2} \right)^3                            
                             = \frac{1}{2^3} \left( \expo{i \theta} + \expo{ - i \theta} \right)^3                         \\
                           & = \frac{1}{8} \left[ (\expo{i \theta})^3 + 3 (\expo{i \theta})^2 \expo{-i \theta} 
                                 + 3 \expo{ i \theta} (\expo{-i \theta})^2 + (\expo{-i \theta})^3 \right]                  \\
                           & = \frac{1}{8} (\expo{3i \theta} + 3 \expo{i \theta} + 3 \expo{-i \theta} + \expo{-3i \theta}) 
                             = \frac{1}{4} \left[ \frac{\expo{i 3\theta} + \expo{-i 3\theta}}{2} 
                                 + 3 \frac{ \expo{i \theta} + \expo{-i \theta} }{2} \right]                                \,.
      \end{align*}
      \mybox{$$ \cos ^3 \theta = \frac{1}{4} \left[ \cos 3\theta + 3 \cos \theta \right] .$$}
\end{ex}

La lin\'earisation de $\sin ^n \theta$ se fait de mani\'ere analogue.

\section{Equations du second degr\'e dans $\mathbb{C}$}
%
\index{Polynomes@Polyn\^omes!du second degre@du second degr\'e!dans C@dans $\mathbb{C}$}
%
Comme on l'a vu en introduction, certaines \'equations polynomiales dans $\mathbb{R}$ n'admettent pas de racine.
On sait d\'ej\`a comment r\'esoudre les \'equations du second degr\'e dans $\mathbb{R}$, cette m\^eme m\'ethode peut \^etre appliqu\'ee dans $\mathbb{C}$.
On va voir dans cette section comment \'etendre cette m\'ethode aux nombres complexes.
%
\subsection{Equations du second degr\'e \`a coefficients r\'eels}
%
\index{Polynomes@Polyn\^omes!du second degre@du second degr\'e!a coefficients reels@\`a coefficients r\'eels}
%
Consid\'erons l'\'equation $a z^2 + b z + c = 0$, avec $a\neq0$, o\`u $z$ est l'inconnue.
On se propose de r\'esoudre cette \'equation dans $\mathbb{C}$, c'est-\`a-dire qu'on cherche les nombres complexes qui la v\'erifient.
Posons $\Delta = b^2 - 4 a c$.
On distingue alors 3 cas:
\begin{itemize}
\item Si $\Delta>0$:
     L'\'equation admet deux racines r\'eelles distinctes: $z' = \frac{-b - \sqrt{\Delta}}{2a}$ et $z'' = \frac{-b + \sqrt{\Delta}}{2a}$.
\item Si $\Delta=0$:
     L'\'equation admet une racine r\'eelle double: $z' = z'' = \frac{-b}{2a}$.
\item Si $\Delta<0$:
     On sait que l'\'equation n'admet pas de racines r\'eelles.
\end{itemize}
Dans ce dernier cas, il nous faut alors chercher des solutions complexes.
Puisque $b^2 - 4 ac<0$, alors $4ac - b^2$ est strictement positif.
On peut alors r\'eecrire $\Delta$ sous la forme $b^2 - 4ac = - (4ac - b^2) = i^2 (4ac - b^2)$.
Puisque $4ac - b^2$ est positif, on peut prendre sa racine carr\'ee et ainsi \'ecrire $\Delta = (i \cdot \sqrt{4ac - b^2} )^2$.
Ainsi, on obtient deux nombres complexes $i \cdot \sqrt{4ac - b^2}$ et $-i \cdot \sqrt{4ac - b^2}$ dont les carr\'es sont \'egaux \`a $\Delta$.
En prenant \mybox{$z' = \frac{-b -i \cdot \sqrt{4ac - b^2} }{2a}$} et \mybox{$z'' = \frac{-b + i \cdot \sqrt{4ac - b^2} }{2a}$}, on v\'erifie bien que $z'$ et $z''$ sont racines de l'\'equation $a z^2 + b z + c = 0$.

\begin{exe}
\label{exe:santatra second degre}
	R\'esoudre dans $\mathbb{C}$ les \'equations:
	\begin{itemize}
	\item $ z^2 - 2 z + 5 = 0 $,
	\item $ z^2 - 4 z + 5 = 0 $,
	\item $ 4 z^2 - 8 z + 5 = 0 $.
	\end{itemize}
\end{exe}
%
\begin{rmk}
      Les racines non r\'eelles d'un polyn\^ome \`a coefficients r\'eels sont deux \`a deux conjugu\'ees.
      Autrement dit, si $P$ est un polyn\^ome \`a coefficients r\'eels et que $z_0$ est une racine de $P$,
      alors $\bar{z_0}$ est aussi une racine de $P$.
\end{rmk}
%
\subsection{Equations du second degr\'e \`a coefficients complexes}
%
\index{Polynomes@Polyn\^omes!du second degre@du second degr\'e!a coefficients complexes@\`a coefficients complexes}
%
Les \'equations du second degr\'e \`a coefficients complexes se r\'esolvent de la m\^eme mani\`ere que celles \`a coefficients r\'eels.
Consid\'erons l'\'equation $a z^2 + b z + c = 0$ o\`u, cette fois, $a$, $b$ et $c$ sont des nombres complexes avec $a\neq0$.
On prend toujours $\Delta = b^2 - 4ac$, et on distingue deux cas:
\begin{itemize}
\item Si $\Delta = 0$:
     L'\'equation admet une racine double \thmbox{$z' = z'' = \frac{-b}{2a}$}.
\item Si $\Delta \neq 0$:
	     Dans ce cas, il existe deux nombres complexes $\delta$ et $-\delta$ tels que $\delta ^2 = (- \delta)^2 = \Delta$.
	     L'\'equation admet deux racines complexes distinctes : \thmbox{$z'=\frac{-b-\delta}{2a}$} et \thmbox{$z''=\frac{-b+\delta}{2a}$}.
\end{itemize}
%
\begin{methode}[Racines carr\'es d'un nombre complexe]
      Soit $Z=X+iY$, on veut chercher un nombre $z=x+iy$ tel que $z^2 = Z$.
      Comme $z^2 = (x+iy)^2 = x^2 - y^2 + 2ixy = X + iY$, 
            $
                  \left\lbrace
                        \begin{array}{l}
                              x^2 - y^2 = X \\
                              2xy = Y       
                        \end{array}
                  \right. 
            $.
      De plus, $|z^2| = |Z|$, donc $x^2 + y^2 = \sqrt{X^2 + Y^2}$.
      
      On peut alors r\'esoudre le syst\`eme suivant pour trouver $z$ :
            \[
                  \left\lbrace
                        \begin{array}{l}
                              x^2 + y^2 = |Z|           \\
                              x^2 - y^2 = \Reel{Z}      \\
                              2 x y     = \Imag{Z}      \;.
                        \end{array}
                  \right. 
            \]
      Les solutions de ce syst\`eme sont :
			\[
				\begin{multiform}
					x = \pm \sqrt{\frac{|Z| + \Reel{Z}}{2}}		\\
					y = \pm \sqrt{\frac{|Z| - \Reel{Z}}{2}}		
				\end{multiform}
			\]
			o\`u $x$ et $y$ sont de m\^eme signes si $\Imag{Z} > 0$ et de signe contraire si $\Imag{Z} < 0$.
\end{methode}
%
\begin{exe}
\label{exe:second degre complexe}
      R\'esoudre dans $\mathbb{C}$ les \'equations suivantes: 
      \begin{itemize}
            \item $z^2 - (1 + 2i) z - 1 + i = 0$,
            \item $i z^2 + 2 z - 2 i = 0$,
            \item $z^2 - (1+3i) z - (4 - 3i) = 0$,
            \item $z^2 + (1-2i) z - 3 - i = 0$.
      \end{itemize}
\end{exe}

\subsection{\'Equations de degr\'e sup\'erieur}
\begin{prop}
      Soient $P$ un polyn\^ome \`a coefficients complexes de degr\'e $n \geq 1$ et $z_0$ un nombre complexe.
      Si $P(z_0) = 0$, alors il existe un polyn\^ome $Q$ de degr\'e $n - 1$ tel que, pour tout $z \in \mathbb{C}$,
      $P(z) = (z - z_0)\,Q(z)$.
\end{prop}
\begin{methode}[R\'esoudre une \'equation du troisi\`eme degr\'e]
      On cherche une racine particuli\`ere $z_0$ (souvent r\'eelle ou imaginaire pure, indiqu\'ee par l'\'enonc\'e) :
      pour $z_0 = x$ r\'eel, $P(x) = 0$ \'equivaut \`a $\Reel{P(x)} = 0$ et $\Imag{P(x)} = 0$ ; pour $z_0 = iy$, on remplace $z$ par $iy$ de m\^eme.
      On d\'etermine ensuite $Q(z) = a z^2 + b z + c$ par identification des coefficients (ou par division),
      puis on r\'esout l'\'equation du second degr\'e $Q(z) = 0$.
\end{methode}

\section{Racine $n^{\text{i\`eme}}$ d'un nombre complexe}
\index{Nombres complexes!racine nieme@racine $n^{\text{i\`eme}}$}
%
\begin{defi}[Racine $n^{\text{i\`eme}}$]
\index{Racine nieme@Racine $n^{\text{i\`eme}}$}
      Soient $Z$ un nombre complexe et $n$ un nombre entier naturel sup\'erieure \`a 1.
      On appelle \emph{racine $n^{\text{i\`eme}}$} de $Z$ tout nombre complexe $z$ v\'erifiant
            \defbox{$z^n = Z$}.
\end{defi}
%
\subsection{Racine $n^{\text{i\`eme}}$ de l'unit\'e}
\index{Nombres complexes!racine nieme@racine $n^{\text{i\`eme}}$!de l'unite@de l'unit\'e}
%
Les racines $n^{\text{i\`eme}}$ de l'unit\'e sont les nombres complexes $z$ v\'erifiant $z^n=1$.
Ce sont les nombres $\left(u_k = \left[ 1,\frac{2k\pi}{n} \right] \right)_{0\leq k <n}$.
\begin{ex}\vspace{-1\baselineskip}
      \begin{description}
            \item[$n=2$ :] $u_0 = 1$ et $u_1 = -1$.
            \item[$n=3$ :] $u_0 = 1$, $u_1 = j$ et $u_2 = \bar{j}$ o\`u $j = -\frac{1}{2}+i\frac{\sqrt{3}}{2}$.
            \item[$n=4$ :] $u_0 = 1$, $u_1 = i$, $u_2 = -1$ et $u_3 = -i$.
      \end{description}
\end{ex}

\subsection{Racine $n^{\text{i\`eme}}$ d'un nombre complexe}
%
Soient $[R,\Phi]$ et $[\rho,\varphi]$ les representations polaires de $Z$ et $z$.
Rappelons que l'argument d'un nombre complexe est d\'efini \`a un multiple de $2 \pi$ pr\`es.
Donc pour que $z^n$ soit \'egal \`a $Z$ il faut que $ [\rho^n,n \varphi] $ soit \'egal \`a $[R,\Phi]$ et donc que:
\begin{itemize}
\item $\rho^n = R$
\item $n \varphi = \Phi + 2k \pi $ pour un certain $k \in \mathbb{Z}$
\end{itemize}
Ainsi, $\rho = \sqrt[n]{R}$ et $\varphi = \frac{\Phi + 2 k \pi}{n}$.
Les possibilit\'es pour $z$ sont alors:
\[
      \left[ \sqrt[n]{R} , \frac{\Phi}{n} \right],
      \left[ \sqrt[n]{R} , \frac{\Phi + 2 \pi}{n} \right],
      \dots,
      \left[ \sqrt[n]{R} , \frac{\Phi + 2 k \pi}{n} \right],
      \dots,
      \left[ \sqrt[n]{R} , \frac{\Phi + 2(n-1) \pi}{n} \right].
      \]
Toute autre valeur de $k$ donnerait l'une de ces valeurs \`a un multiple de $2 \pi$ pr\`es.
Cela nous m\`ene au th\'eor\`eme suivant : 
\begin{thm}
Soit $n$ un entier naturel non nul et $z = [r,\theta]$ un nombre complexe non nul.
Alors $z$ admet $n$ racines $n^{\text{i\`eme}}$ qui sont :
$$
z_k = \sqrt[n]{r} \left[ \cos \left( \frac{\theta + 2k \pi}{n} \right) + i \sin \left( \frac{\theta + 2k \pi}{n} \right) \right] \text{ avec } 0 \leq k < n
$$
\end{thm}
%
\begin{rmk}
Repr\'esent\'ees g\'eom\'etriquement dans le plan complexe, ces racines $n^{\text{i\`eme}}$ forment les sommets d'un polyg\^one r\'egulier centr\'e en l'origine.
\end{rmk}
%
\begin{exe}
\label{exe:racines de l'unite}
%
\index{Nombres complexes!racine nieme@racine $n^{\text{i\`eme}}$!de l'unite@de l'unit\'e}
%
Donner les racines cubiques et racine quatri\`eme de l'unit\'e.

\end{exe}
%
\begin{exe}
\label{exe:racine nieme}
Donner les racines $5^{\text{i\`eme}}$ de $1+ \sqrt{3} i$. \\
\end{exe}
%
\section{Transformation de $a \cos x + b \sin x$}
%
\index{Fonctions trigonometriques@Fonctions trigonom\'etriques!Transformation de a cos x + b sin x@Transformation de $a \cos x + b \sin x$}
%
On se propose de r\'esoudre l'\'equation 
\begin{equation}
a \cos x + b \sin x = c
\label{eq:a cos x + b sin x}
\end{equation}
o\`u $a$, $b$ et $c$ sont des nombres r\'eels.
Posons $z = a+ i b $ et soient $r$ son module et $\theta$ son argument.
On sait que $ a = r \cos \theta$ et $b = r \sin \theta$ avec $r = \sqrt{a^2 + b^2}$.
Ainsi, l'expression $a \cos x + b \sin x $ devient
$$
r \cos \theta \, \cos x + r \sin \theta \, \sin x = r \left( \cos \theta \, \cos x + \sin \theta \, \sin x \right) \; .
$$
Faisant usage de formule trigonom\'etrique connue \footnote{$\cos( u - v ) = \cos u \, \cos v + \sin u \, \sin v$}, et de l'\'equation $r = \sqrt{a^2 + b^2}$, on obtient
$$
a \, \cos x + b \, \sin x = \sqrt{a^2 + b^2} \; \cos (\theta - x) \; .
$$
Ainsi, l'Equation \ref{eq:a cos x + b sin x} devient 
\begin{equation}
\cos (\theta - x) = \frac{c}{\sqrt{a^2 + b^2}} \; .
\label{eq:cos (theta - x) = c/r}
\end{equation}
On distingue alors $2$ cas : 
\begin{itemize}
\item \textbf{ Si $ \left| \frac{c}{\sqrt{a^2 + b^2}} \right| > 1$ : } \newline
      Alors l'\'equation n'admet pas de solution.
\item \textbf{ Si $ \left| \frac{c}{\sqrt{a^2 + b^2}} \right| \leq 1$ : } \newline
      Alors il existe $\alpha$ tel que $\cos \alpha =  \frac{c}{\sqrt{a^2 + b^2}} $.
      Les solutions sont alors les $x$ tels que 
      \begin{align*}
      \theta - x = \alpha + 2 k \pi 
      \quad \text{ ou } \quad
      x - \theta = \alpha + 2 k \pi
      \end{align*}
      Ainsi, l'ensemble des solutions est
      \[
      \left\lbrace \theta - \alpha + 2 k \pi \mid k \in \mathbb{Z} \right\rbrace 
      \cup 
      \left\lbrace \theta + \alpha + 2 k \pi \mid k \in \mathbb{Z} \right\rbrace
      .
      \]
\end{itemize}
%
\begin{ex}
On veut r\'esoudre l'\'equation 
\begin{equation}
\cos x + \sqrt{3} \sin x = \sqrt{2} \; .
\label{eq:ex: cos x + sin x}
\end{equation}
Posons $z = 1 + i \sqrt{3}$; on a alors $z = 2 e^{i \frac{\pi}{3}}$.
Ainsi, $1 = 2 \cos \frac{\pi}{3}$ et $\sqrt{3} = 2 \sin \frac{\pi}{3}$.
L'Equation \ref{eq:ex: cos x + sin x} devient alors 
\begin{align*}
2 \left[ \cos \frac{\pi}{3} \; \cos x + \sin \frac{\pi}{3} \; \sin x \right] & = \sqrt{2} \\
\cos \left( \frac{\pi}{3} - x \right) & = \frac{\sqrt{2}}{2} \\
\frac{\pi}{3} - x = \frac{\pi}{4} + 2 k \pi \quad \text{ ou } \quad x - \frac{\pi}{3} = \frac{\pi}{4} + 2 k \pi \\
\end{align*}
ce qui nous donne comme solutions : 
$
\left\lbrace \frac{\pi}{12} + 2 k \pi \mid k \in \mathbb{Z} \right\rbrace 
\cup 
\left\lbrace \frac{7 \pi }{12} + 2 k \pi \mid k \in \mathbb{Z} \right\rbrace
$.
\end{ex}
%
\section{Utilisation en g\'eom\'etrie}
%
\index{Nombres complexes!Utilisation en geometrie@Utilisation en g\'eom\'etrie}
%
On rapporte le plan complexe \`a un rep\`ere orthonorm\'e direct $(O,\vec{e_1},\vec{e_2})$.
Soit $f$ une transformation du plan qui au point $M$ d'affixe $z$ associe le point $M'$ d'affixe $z'$.
La relation qui exprime $z'$ en fonction de $z$ est appel\'ee \emph{\'ecriture complexe} de $f$.
%
\subsection{Translation}
\index{Nombres complexes!Utilisation en geometrie@Utilisation en g\'eom\'etrie!Translation}
%
\begin{figure}[h!]
\centering
%      \begin{wrapfigure}{r}{0.35\actualtextwidth}
            \begin{tikzpicture}[scale=1.2,>=stealth]
            \clip (-.5,-1.5) rectangle (2.5,2);
            \draw [color=black, xstep=.1,ystep=.1, opacity=.1, line width=.1] (-1.5,-1.5) grid (3.5,2.5);
            \draw [color=black, xstep=.5,ystep=.5, opacity=.2, line width=.2] (-1.5,-1.5) grid (3.5,2.5);
            \draw [color=black, xstep=1,ystep=1, opacity=.5, line width=.5] (-1.5,-1.5) grid (3.5,2.5);
            \draw[->,color=black,thick] (-1.5,0) -- (2.5,0);
            \foreach \x in {-1,1,2,3}
                  \draw[shift={(\x,0)},color=black] (0pt,2pt) -- (0pt,-2pt) node[below] {\tiny $\x$};
            \draw[->,color=black,thick] (0,-1.5) -- (0,2);
            \foreach \y in {-1,1}
                  \draw[shift={(0,\y)},color=black] (2pt,0pt) -- (-2pt,0pt) node[left] {\tiny $\y$};
            \draw[color=black] (0pt,0pt) node[above right] {\tiny $0$};
            \begin{scriptsize}
                  \fill[loko1] (0.5,-1) circle (1pt) node[below] {$M$} ;
                  \fill[loko1] (2,1.5) circle (1pt) node[above] {$M'$} ;
            \end{scriptsize}
            \draw[->,thick,color=loko2] (0.5,-1) -- node[left] {$\vec{v}$} (2,1.5);            
            \end{tikzpicture}
%      \end{wrapfigure}
	\caption{Le point $M'$ est le translat\'e du point $M$ par le vecteur $\vec{v}$.}
\end{figure}
%
      Soit $M (z_M)$ un point du plan complexe et $\overrightarrow{V}$ un vecteur d'affixe $z_{\overrightarrow{V}}$.
      Le translat\'e de $M$ par le vecteur $\overrightarrow{V}$ est le point $M'$ d'affixe $z_M + z_{\overrightarrow{V}}$.
      
      En effet, comme $z_M$ est l'affixe de $\overrightarrow{OM}$ et que $z_{\overrightarrow{V}}$ est l'affixe de $\overrightarrow{V}$ 
      alors l'affixe du vecteur $\overrightarrow{OM} + \overrightarrow{V}$ est $z_M + z_{\overrightarrow{V}}$ 
      ce qui doit correspondre \`a l'affixe de $\overrightarrow{OM'}$.
      En identifiant les points aux vecteurs, on r\'ealise que l'affixe du point $M'$ est $z_M + z_{\overrightarrow{V}}$.

\subsection{Homoth\'etie}
Soient $P$ un point du plan d'affixe $z_P$ et $k\in\mathbb{R}^*$.
Soient $h$ l'homoth\'etie de centre $P$ et de rapport $k$ et $M'$ l'image de $M$ par $h$.
On a alors $z' - z_P = k (z - z_P) $.

L'\'ecriture complexe de $h$ est donc \thmbox{$ z' = k \cdot z + (1-k) z_P $}.
%
\subsection{Angle entre deux vecteurs}
\index{Nombres complexes!Utilisation en geometrie@Utilisation en g\'eom\'etrie!Angle entre deux vecteurs}
%
      Soient $A$, $B$ et $C$ trois points du plan complexe d'affixes respectifs $z_A$, $z_B$ et $z_C$.
      L'affixe du vecteur $\overrightarrow{AB}$ est alors $z_B - z_A$ et l'affixe du vecteur $\overrightarrow{AC}$ est $z_C - z_A$.
      D\'esignons par $\vec{u}$ le vecteur d'affixe $1$.
%
		\begin{figure}[h!]
			\centering
			\begin{tikzpicture}[line cap=round,line join=round,>=stealth,x=0.15\actualtextwidth,y=0.15\actualtextwidth]
				\clip(-0.5,-1) rectangle (2.5,2.5);
					\draw [color=gray,line width=.5, xstep=1.,ystep=1., opacity=.5] (-0.5,-1) grid (2.5,2.5);
					\draw [color=gray,line width=.3, xstep=.5,ystep=.5, opacity=.3] (-0.5,-1) grid (2.5,2.5);
					\draw [color=gray,line width=.2, xstep=.1,ystep=.1, opacity=.2] (-0.5,-1) grid (2.5,2.5);
					\draw[->,color=black] (-0.5,0) -- (2.5,0);
					\draw[->,color=black] (0,-1) -- (0,2.5);
					\draw [shift={(0.5,0.5)},color=loko3!20!black,fill=loko3!20!black,fill opacity=0.1] 
						(0,0) -- (-63.43:0.26) arc (-63.43:0:0.26) -- cycle;
					\draw [shift={(0.5,0.5)},color=loko1,fill=loko1,fill opacity=0.1] 
						(0,0) -- (0:0.28) arc (0:56.31:0.28) -- cycle;
					\draw [shift={(0.5,0.5)},color=loko1] (0:0.28) arc (0:56.31:0.28);
					\draw [shift={(0.5,0.5)},color=loko1] (0:0.23) arc (0:56.31:0.23);
					\begin{footnotesize}
						\fill [color=loko2] (0.5,0.5) circle (1.5pt);
						\draw[color=loko2] (0.3,0.65) node {$A$};
						\fill [color=loko2] (1,-0.5) circle (1.5pt);
						\draw[color=loko2] (0.85,-0.65) node {$B$};
						\fill [color=loko2] (1.5,2) circle (1.5pt);
						\draw[color=loko2] (1.35,2.15) node {$C$};
						\draw[color=black] (1.55,0.5) node {$\vec{u}$};
						\draw[color=loko3!20!black] (1.25,0.25) node {$\left( \overrightarrow{AB} , \vec{u} \right)$};
						\draw[color=loko1] (1.25,0.75) node {$\left( \vec{u} , \overrightarrow{AC} \right)$};
					\end{footnotesize}
					\draw [->] (0.5,0.5) -- (1,-0.5);
					\draw [->] (0.5,0.5) -- (1.5,2);
					\draw [->] (0.5,0.5) -- (1.5,0.5);
			\end{tikzpicture}
		\end{figure}
%
      On sait que $\Mes{\vec{u}}{\overrightarrow{AB}} = \arg(z_B - z_A)$ et 
      $\Mes{\vec{u}}{\overrightarrow{AC}} = \arg(z_C - z_A)$.
      Il suit que 
      $\mes{AB}{AC} = \arg(z_C-z_A)-\arg(z_B-z_A)=\arg\frac{z_C - z_A}{z_B - z_A}$.
\begin{exe}
\label{exe:angle}
Soient $A (-1 - i)$, $B (2 + 2 i)$ et $C (-5 + i ) $ trois points du plan.
Calculer la mesure de l'angle $\hat{A}$ du triangle $ABC$.
\end{exe}
%
\subsection{Rotations}
%
\index{Nombres complexes!Utilisation en geometrie@Utilisation en g\'eom\'etrie!Rotations}
%
Soient $P$ un point du plan complexe et $\theta$ un nombre r\'eel.
On d\'esigne par $r$ la rotation de centre $P$ et d'angle $\theta$.
Soit $M$ un point du plan distinct de $P$ et soit $M'$ son image par $r$.
On sait alors que 
\begin{equation}
\norm{PM} = \norm{PM'} 
\label{eq:1}
\end{equation}
et 
\begin{equation}
\mathrm{mes}\left(\overrightarrow{PM} , \overrightarrow{PM'} \right) = \theta \; .
\label{eq:2}
\end{equation}
Notons par $z_P$, $z_M$ et $z_{M'}$ les affixes de $P$, $M$ et $M'$ respectivement.
En utilisant les affixes, on a alors 
$
\left\lbrace
\begin{matrix}
| z_{M'} - z_P | = | z_M - z_P | \\
\text{arg} \left( \frac{z_{M'} - z_P}{z_M - z_P } \right) = \theta
\end{matrix}
\right.
$
et donc 
$
\left\lbrace
\begin{matrix}
\frac{| z_{M'} - z_P |}{| z_M - z_P |} = 1 \\
\text{arg} \left( \frac{z_{M'} - z_P}{z_M - z_P } \right) = \theta
\end{matrix}
\right.
$
. 

En posant $Z = \frac{z_{M'} - z_P}{z_M - z_P}$, on a $|Z| = 1$ et $\text{arg} (Z) = \theta$, ce qui entra\^ine que 
$Z = \cos \theta + i \sin \theta$. 
Cela nous permet d'\'ecrire $( z_{M'} - z_P ) = ( \cos \theta + i \sin \theta) \; ( z_M - z_P )$ et donc 
\mybox{$ z_{M'}  = ( \cos \theta + i \sin \theta) \; ( z_M - z_P ) + z_P $}.
Ainsi, l'\'ecriture complexe de $r$ est : 
\begin{center}
\thmbox{$ z' = \expo{i\theta} (z - z_P) + z_P$}.
\end{center}
%
\begin{exe}
\label{exe:rotation}
Soient $A (2 + i ) $ et $B (3 - 2 i)$ deux points du plan. 
D\'eterminer l'image de $B$ par la rotation de centre $A$ et d'angle $\frac{\pi}{3}$.
D\'eterminer l'image de $B$ par la rotation de centre $A$ et d'angle $-\frac{\pi}{3}$.
\end{exe}
%
\subsection{Similitudes planes directes}
      \begin{defi}[Similitude plane directe]
      \index{Transformation du plan!Similitude plane directe}
      On appelle \emph{similitude plane directe} toute transformation du plan qui conserve les angles orient\'es
      et qui multiplie les distances par une constante non nulle.
      \end{defi}
      %
      \begin{thm}
      Une similitude plane directe appartient \`a l'une des deux cat\'egories suivantes :
            \begin{itemize}
                  \item translation,
                  \item composition d'une rotation et d'une homoth\'etie de m\^eme centre.
            \end{itemize}
      \end{thm}
      %
      \begin{prop}
      Une similitude plane directe qui n'est pas une translation est caract\'eris\'ee par : 
            \begin{itemize}
                  \item son centre,
                  \item son angle,
                  \item son rapport.
            \end{itemize}
      \end{prop}
      %
      \begin{prop}[\'Ecriture complexe]
            Soient $P$, d'affixe $z_P$, un point, $\theta$ un nombre r\'eel et $k$ un nombre r\'eel strictement positif.
            Notons par $S$ la similitude plane directe de centre $P$, d'angle $\theta$ et de rapport $k$.
            
            Pour tout point $M$ d'affixe $z$, on note par $M'$ son image par $S$ et par $z'$ l'affixe de $M'$.            
            On a alors :
                  \begin{displaymath}
                        \left\lbrace
                              \begin{array}{l}
                                    PM' = k\, PM \\
                                    \mathrm{mes} \left( \overrightarrow{PM} ; \overrightarrow{PM'} \right) = \theta
                              \end{array}
                        \right.
                  \end{displaymath}
            Avec ces notations, l'\'ecriture complexe de $S$ est \thmbox{$ z' = k \expo{i\theta} (z - z_P) + z_P $}.
      \end{prop}
      %
            \subsection*{Classification des similitudes planes directes}
            L'expression complexe d'une similitude plane directe est de la forme $z'=a\,z+b$, avec $a \neq 0$.
            \begin{figure}[h!]
            \centering
            \resizebox{\linewidth}{!}{%
            \begin{tikzpicture}[line cap=round,line join=round,>=triangle 45,x=.055\actualtextwidth,y=.075\textheight]
            \clip(-9,-1.5) rectangle (11,6.5);
            \draw (0,6) node[ ] {\fbox{$ z' = a\,z+b $}};
            \draw (-4,4.5) node[ ] {\fbox{$ a \in \mathbb{R}_+^* $}};
            \draw (4,4.5) node[ ] {\fbox{$ a \not \in \mathbb{R}_+ $, $a \neq 0$}};
            \draw (-6,3) node[ ] { \fbox{$a = 1 $}};
            \draw (-2,3) node[ ] { \fbox{$a \neq 1$} };
            \draw (2,3) node[ ] { \fbox{$|a| = 1$} };
            \draw (6,3) node[ ] { \fbox{$|a| \neq 1$} };
            %
            \draw [->] (0,5.75) to[out=-90,in=90] (-4,4.75);
            \draw [->] (0,5.75) to[out=-90,in=90] (4,4.75);
            \draw [->] (-4,4.25) to[out=-90,in=90] (-6,3.25);
            \draw [->] (-4,4.25) to[out=-90,in=90] (-2,3.25);
            \draw [->] (4,4.25) to[out=-90,in=90] (6,3.25);
            \draw [->] (4,4.25) to[out=-90,in=90] (2,3.25);
            %
            \draw (-6,2.5) node[below] {\fbox{\parbox{.2\actualtextwidth}{\raggedright %
                  Translation de vecteur $\vec{u}$ d'affixe $b$.}}};
            \draw (-2,1) node[below] %
                  {\fbox{\parbox{.2\actualtextwidth}{\raggedright %
                        Homoth\'etie de centre $I$ d'affixe $\frac{b}{1-a}$ de rapport $a$.}}};
            \draw (2,2.5) node[below] %
                  {\fbox{\parbox{.29\actualtextwidth}%
                  {\raggedright %
                   Rotation de centre $I$ d'affixe $\frac{b}{1-a}$ d'angle $\theta = \text{arg} (a) \qquad [2\;\pi]$.}}};
            \draw (6,1) node[below] %
                  {\fbox{\parbox{.35\actualtextwidth}%
                  {\raggedright Similitude plane directe de centre $I$ d'affixe $\frac{b}{1-a}$, %
                  de rapport $k =|a|$ et d'angle $\theta = \text{arg}(a) \qquad [2\,\pi]$.}}};
            %
            \draw [->] (-6,2.75) to[out=-90,in=90] (-6,2.5);
            \draw [->] (-2,2.75) to[out=-90,in=90] (-2,1);
            \draw [->] (2,2.75) to[out=-90,in=90] (2,2.5);
            \draw [->] (6,2.75) to[out=-90,in=90] (6,1);
            \end{tikzpicture}
            }
            \caption{Classification des similitudes planes directes.}
            \label{fig:classification similitudes planes directes}
            \end{figure}
\end{chapter}

\TCDLessonStyleOff
\providecommand{\Reel}[1]{\mathfrak{Re}\left(#1\right)}
\providecommand{\Imag}[1]{\mathfrak{Im}\left(#1\right)}
\providecommand{\argmt}[1]{\mathrm{arg}\left( #1 \right)}
\providecommand{\mes}[2]{\mathrm{mes}\left( \widehat{ \overrightarrow{#1} , \overrightarrow{#2} } \right)}
\providecommand{\AngleEgal}[3]{#1 = #2 \quad \left[ #3 \right] }
\begin{fiche}
      \subsection*{Diff\'erentes formes d'un nombre complexe}~\\[-2.5\baselineskip]
\index{Nombre complexe!formes}
            \[
                  \begin{array}{|c|c|c|}
                        \hline
                        \rowcolor[gray]{0.8}
                              \color{black}
                              \text{Forme algébrique} & \text{Forme trigonométrique}        & \text{Forme exponentielle} \\[2pt]
                        \hline
                              z = a + i b             & z = r (\cos \theta + i \sin \theta) & z = r \expo{i\theta}       \\
                        \hline
                              \displaystyle{
                                 \left\lbrace
                                    \begin{array}{l}
                                       a = r \cos \theta \\
                                       b = r \sin \theta 
                                    \end{array}
                                 \right.
                              }
                                                      & 
                                                  \multicolumn{2}{c|}{
                                                     r = \sqrt{a^2 + b^2}
                                                     \quad \text{ et } \quad 
                                                     \displaystyle{
                                                         \left\lbrace
                                                            \begin{array}{l}
                                                               \cos \theta = \frac{a}{r} \\
                                                               \sin \theta = \frac{b}{r} 
                                                            \end{array}
                                                         \right.
                                                      }
                                                                    }
                                                                                                                         \\
                        \hline
                  \end{array}
            \]
      %
      ~\\[-2\baselineskip]
      %
      \subsection*{Propri\'et\'es}
\index{Nombre complexe!Proprietes@Propri\'et\'es}
                  \subsubsection*{Sur les conjugu\'es} ~\\[-1.5\baselineskip]
\index{Nombre complexe!Proprietes@Propri\'et\'es!conjugues@conjug\'es}
                              \begin{itemize}
                                    \item $\overline{a+ib} = a-ib$,
                                    \item $\bar{\bar{z}} = z$,
                                    \item $z \cdot \bar{z} = |z|^2$,
                                    \item $\overline{z_1 + z_2} = \bar{z_1} + \bar{z_2}$,
                                    \item $\overline{z_1 \cdot z_2} = \bar{z_1} \cdot \bar{z_2}$,
                                    \item $\overline{\left(\frac{z_1}{z_2}\right)} = \displaystyle{\frac{\bar{z_1}}{\bar{z_2}}}$.
                              \end{itemize}
                  \subsubsection*{Sur les modules} ~\\[-1.5\baselineskip]
\index{Nombre complexe!Proprietes@Propri\'et\'es!modules}
%                  \begin{minipage}{.4\actualtextwidth}
                        \begin{itemize}
                              \item $|z_1 \cdot z_2| = |z_1| \cdot |z_2|$,
                              \item $\left| \frac{z_1}{z_2} \right| = \frac{|z_1|}{|z_2|}$,
                              \item $|z^n| = |z|^n$,
                              \item $|\bar{z}| = |z|$,
                              \item $|z_1 + z_2| \leq |z_1| + |z_2|$.
                        \end{itemize}
%                  \end{minipage}
                  \subsubsection*{Sur les arguments} ~\\[-1.5\baselineskip]
\index{Nombre complexe!Proprietes@Propri\'et\'es!arguments}
%                  \begin{minipage}{.4\actualtextwidth}
                        \begin{itemize}
                              \item $\argmt{z_1 \cdot z_2} = \argmt{z_1} + \argmt{z_2}$,
                              \item $\argmt{\frac{1}{z}} = -\argmt{z}$,
                              \item $\argmt{\frac{z_1}{z_2}} = \argmt{z_1} - \argmt{z_2}$,
                              \item $\argmt{z^n} = n \cdot \argmt{z}$,
                              \item $\argmt{\bar{z}} = -\argmt{z}$.
                        \end{itemize}
%                  \end{minipage}
                  \subsection*{Utilisant les affixes} ~\\[-1.5\baselineskip]
\index{Nombre complexe!Proprietes@Propri\'et\'es!angle entre deux vecteurs}
                        Soient $A$, $B$, $C$ et $D$ quatre points d'affixes respectifs $a$, $b$, $c$ et $d$.
                        On a alors : 
                              \[
                              \mes{AB}{CD} = \argmt{\frac{d - c}{b - a} }\;.
                              \]
                        
\index{Nombre complexe!Proprietes@Propri\'et\'es!identification de parallelogramme@identification de parall\'elogramme}
                        Le quadrilat\`ere $ABCD$ est un parall\'elogramme si et seulement si : 
                              $a-b = d-c$.
      \subsection*{Ensembles de points}
\index{Ensemble de points!sous condition sur les angles}
            \subsubsection*{D\'efini par l'argument}
                  \begin{itemize}
	                        \item L'ensemble des points $M(z)$ v\'erifiant
	                              \[
	                                    \AngleEgal{\argmt{\frac{z-z_A}{z-z_B}}}{0}{2\pi}
	                                          \iff
	                                    \AngleEgal{\mes{BM}{AM}}{0}{2\pi}
	                              \]
                              est
                              $(AB) \setminus [AB]$ ;
	                        \item l'ensemble des points $M(z)$ v\'erifiant
	                              \[
	                                    \AngleEgal{\argmt{\frac{z-z_A}{z-z_B}}}{0}{\pi}
	                                          \iff
	                                    \AngleEgal{\mes{BM}{AM}}{0}{\pi}
	                              \]
                              est
                              $(AB) \setminus \{A,B\}$ ;
	                        \item l'ensemble des points $M(z)$ v\'erifiant
	                              \[
	                                    \AngleEgal{\argmt{\frac{z-z_A}{z-z_B}}}{\pi}{2\pi}
	                                          \iff
	                                    \AngleEgal{\mes{BM}{AM}}{\pi}{2\pi}
	                              \]
                              est 
                              $]AB[$ ;
	                        \item l'ensemble des points $M(z)$ v\'erifiant
	                              \[
	                                    \AngleEgal{\argmt{\frac{z-z_A}{z-z_B}}}{\frac{\pi}{2}}{\pi}
	                                          \iff
	                                    \AngleEgal{\mes{BM}{AM}}{\frac{\pi}{2}}{\pi}
	                              \]
                              est
                              le cercle de diam\`etre $[AB]$ priv\'e des points $A$ et $B$ ;
	                        \item l'ensemble des points $M(z)$ v\'erifiant
	                              \[
	                                    \AngleEgal{\argmt{\frac{z-z_A}{z-z_B}}}{\frac{\pi}{2}}{2\pi}
	                                          \iff
	                                    \AngleEgal{\mes{BM}{AM}}{\frac{\pi}{2}}{2\pi}
	                              \]
                              est
                              l'ensemble $\mathscr{E}$ repr\'esent\'e ci-dessous : un demi-cercle de diam\`etre $[AB]$, priv\'e de $A$ et de $B$.
                              
                              \begin{tikzpicture}[scale=.75]
                              \draw [dotted] (-0.9,1.58) circle (2);
                              \draw [dash pattern=on 2pt off 2pt] (-2.86,1.19)-- (1.06,1.97);
                              \draw [shift={(-0.9,1.58)},line width=1.2pt,color=loko1,smooth]  
                                    plot[domain=-2.95:0.19,variable=\t]
                                          ({1*2*cos(\t r)+0*2*sin(\t r)},{0*2*cos(\t r)+1*2*sin(\t r)});
                              \fill [color=black] (-0.9,1.58) circle (1.0pt);
                              \draw [color=loko1,fill=white] (-2.86,1.19) circle (1.5pt) node[left] {$A$};
                              \draw [color=black,fill=white] (1.06,1.97) circle (1.5pt) node[right] {$B$};
                              \draw[color=loko1] (0.54,-0.12) node {$\mathscr{E}$};
                              \end{tikzpicture}
                  \end{itemize}
                  %
                  \subsubsection*{D\'efini par le module}
\index{Ensemble de points!sous condition sur les distances}
                        Soient $A(a)$ et $B(b)$ deux points.
                        L'ensemble des points $M(z)$ v\'erifiant
                        $|z-a| = |z-b|$ (c'est-\`a-dire v\'erifiant $AM = BM$)
                        est la \emph{m\'ediatrice de $[AB]$}.
                  
                  %
                  \subsubsection*{Sous condition sur la partie réelle ou imaginaire}
\index{Ensemble de points!sous condition sur la partie reelle@sous condition sous la partie r\'elle}
\index{Ensemble de points!sous condition sur la partie imaginaire}
                        Pour trouver l’ensemble des points $M(z)$ tels que $\Reel{f (z)} = 0$ ou $\Imag{f (z)} = 0$, 
                        il faut souvent remplacer $z$ par $x + iy$ dans $f (z)$.
                        \begin{ex}
                              Soit $f$ l'application d\'efinie sur $\mathbb{C}^*$ par
                              $\displaystyle{ f(z) = \frac{1}{iz} }$.
                              
                              On veut trouver l'ensemble des points $M(z)$ du plan tels que $f(z)$ est r\'eel.
                              On pose alors $z = x + i y$ et on a 
                              \begin{align*}
                                    f(z) & = \frac{1}{i(x+iy)}                        
                                           = \frac{1}{-y+ix}                          
                                           = \frac{-y-ix}{y^2 + x^2}                  \\
                                         & = -\frac{y}{x^2+y^2}-\frac{x}{x^2+y^2}i    \;.
                              \end{align*}
                              Pour que $f(z)$ soit r\'eel, il faut que $-\frac{x}{x^2+y^2}$ soit nul
                              et donc que $x$ soit \'egal \`a $0$.
                              
                              Ainsi, l'ensemble des points $M(z)$ tels que $f(z)$ soit r\'eel est 
                              la droite d'\'equation $x=0$ priv\'ee du point $O$ (car $z=0$ est exclu).
                        \end{ex}
      %
      \subsection*{Pour trouver une racine carr\'ee d'un nombre complexe}
\index{Nombre complexe!racines carree@racines carr\'ees}
            Soit $\Delta\in\mathbb{C}$, on cherche $\delta = x+iy$ v\'erifiant $\delta^2 = \Delta$.
            Les r\'eels $x$ et $y$ v\'erifient 
            \[
                  \left\lbrace
                        \begin{array}{l}
                              x^2 + y^2 = |\Delta|       \\
                              x^2 - y^2 = \Reel{\Delta}  \\
                              2 x y     = \Imag{\Delta}  
                        \end{array}
                  \right.
            \]
            ce qui donne
            \[
					\delta = \sqrt{\frac{|\Delta| + \Reel{\Delta}}{2}} \pm \sqrt{\frac{|\Delta| - \Reel{\Delta}}{2}} \; i
            \]
            o\`u on a <<~$+$~>> si $\Imag{\Delta} > 0$ et <<~$-$~>> si $\Imag{\Delta} < 0$ ;
            l'autre racine carr\'ee de $\Delta$ est $-\delta$.
            
      
      
      
      
      
      
      
      
      
      
      
      
      
\end{fiche}

\TCDLessonStyleOff
\clearpage
\providecommand{\Arg}[1]{\mathrm{Arg}\left( #1 \right)}
\section{Exercices}
%
\begin{exo}
      On consid\`ere les nombres complexes suivants :
            \[
                  z_1 = (1+i) (\sqrt{3}-i)
 ; \quad
                  z_2 = (1+i)^4 \cdot (\sqrt{3}-i)^3
 ; \quad
                  z_3 = \frac{1+i}{\sqrt{3}-i}
                  .
            \]
      \begin{enumerate}
            \item Ecrire chacun d'eux sous forme alg\'ebrique.
                  En d\'eduire leurs parties r\'eelles et imaginaires.
            \item D\'eterminer le module et un argument de chacun des nombres complexes $1+i$ et $\sqrt{3}-i$
                  puis donner la forme trigonom\'etrique de $z_1$, $z_2$ et $z_3$.
      \end{enumerate}
\end{exo}
%
\begin{exo}
      On pose $z=2\sqrt{3} + 2i$, $z' = 1 + \sqrt{3} + (1-\sqrt{3})i$ et $z'' = \frac{z}{z'}$.
      \begin{enumerate}
            \item Donner la forme alg\'ebrique de $z''$.
            \item Donner la forme trigonom\'etrique de $z''$ et de $z$.
            \item En d\'eduire $|z'|$, $\arg z'$, $\cos \frac{\pi}{12}$, $\sin \frac{\pi}{12}$ et $\tan \frac{\pi}{12}$.
      \end{enumerate}
\end{exo}
%
%%%%%%%%%%%%%%%%%%%%%%%%%%%%%%%%%%%%%%%%%%%%%%%%%%%%%%%%%%%%%%%%%%%%%%%%%%%%%%%%%%%%
\begin{exo}
	On consid\`ere les nombres complexes $p = \sqrt{3} + i$ et $q = -2 + 2i$.
	\begin{enumerate}
		\item D\'eterminer le module et un argument de $p$ et de $q$.
		\item En d\'eduire, sans d\'evelopper, le module et un argument de chacun des nombres
		      $p^2$, $\frac{q}{p}$, $\frac{1}{q}$, $p^6$ et $\frac{q^4}{p^3}$.
		\item Donner la forme alg\'ebrique de $p^6$ et de $\frac{q^4}{p^3}$.
		\item \'Ecrire $\frac{q}{p}$ sous forme alg\'ebrique et en d\'eduire les valeurs exactes de
		      $\cos \frac{7\pi}{12}$ et $\sin \frac{7\pi}{12}$.
	\end{enumerate}
\end{exo}
%
\begin{exo}
	Soient $z = a + ib$ et $z' = a' + ib'$ deux nombres complexes non nuls ($a$, $b$, $a'$, $b'$ r\'eels).
	\begin{enumerate}
		\item V\'erifier l'identit\'e $(aa' - bb')^2 + (ab' + a'b)^2 = (a^2 + b^2)(a'^2 + b'^2)$
		      et en d\'eduire que $|z z'| = |z| \, |z'|$.
		\item On \'ecrit $z = r(\cos\theta + i\sin\theta)$ et $z' = r'(\cos\theta' + i\sin\theta')$.
		      En utilisant les formules d'addition, montrer que $\theta + \theta'$ est un argument de $zz'$.
		\item Montrer par r\'ecurrence que, pour tout entier $n \geq 1$, $|z^n| = |z|^n$ et $n\theta$ est un argument de $z^n$.
		\item En remarquant que $z' \cdot \frac{z}{z'} = z$, en d\'eduire le module et un argument de $\frac{z}{z'}$.
		\item Application : calculer, sans d\'evelopper, le module de $\frac{(3+4i)^5}{(1-i)^4}$.
	\end{enumerate}
\end{exo}
%
\begin{exo}
	Soient $z$ et $z'$ deux nombres complexes.
	\begin{enumerate}
		\item Montrer que $|z + z'|^2 = |z|^2 + |z'|^2 + 2\,\mathfrak{Re}\left(z \overline{z'}\right)$.
		\item Justifier que, pour tout nombre complexe $w$, $\mathfrak{Re}(w) \leq |w|$.
		      En d\'eduire l'in\'egalit\'e triangulaire $|z + z'| \leq |z| + |z'|$.
		\item En \'ecrivant $z = (z - z') + z'$, montrer que $\left| |z| - |z'| \right| \leq |z - z'|$.
		\item Application : soit $z$ un nombre complexe de module $2$.
		      Montrer que $3 \leq |z + 3 - 4i| \leq 7$ et trouver un nombre $z$ de module $2$ pour lequel $|z + 3 - 4i| = 7$.
	\end{enumerate}
\end{exo}
%
\begin{exo}
	Soit $\omega$ un nombre complexe diff\'erent de $-1$. On pose $Z = \frac{\omega - 1}{\omega + 1}$.
	\begin{enumerate}
		\item On suppose $|\omega| = 1$.
		      Justifier que $\overline{\omega} = \frac{1}{\omega}$, puis montrer que $\overline{Z} = -Z$.
		      Que peut-on dire de $Z$\,?
		\item R\'eciproquement, on suppose que $Z = iy$ avec $y$ r\'eel.
		      Exprimer $\omega$ en fonction de $y$ et montrer que $|\omega| = 1$.
		\item On pose $\omega = \expo{i\theta}$ avec $\theta \in \left] -\pi , \pi \right[$.
		      Montrer que $Z = i \tan \frac{\theta}{2}$
		      (on pourra mettre $\expo{i\frac{\theta}{2}}$ en facteur au num\'erateur et au d\'enominateur).
		\item Interpr\'etation : on note $A$, $B$ et $M$ les points d'affixes $1$, $-1$ et $\omega$, avec $|\omega| = 1$ et $\omega \neq \pm 1$.
		      Que repr\'esente un argument de $Z$\,? Retrouver ainsi que le triangle $AMB$ est rectangle en $M$.
	\end{enumerate}
\end{exo}
%
\begin{exo}
	On consid\`ere l'application $g$ d\'efinie sur $\mathbb{C}^*$ par $g(z) = z + \frac{1}{z}$.
	\begin{enumerate}
		\item Montrer que, pour tous $z$ et $z'$ non nuls,
		      $g(z) - g(z') = \frac{(z - z')(z z' - 1)}{z z'}$.
		      En d\'eduire que $g(z) = g(z')$ si et seulement si $z = z'$ ou $z z' = 1$.
		\item Montrer que si $|z| > 1$, $|z'| > 1$ et $g(z) = g(z')$, alors $z = z'$.
		\item On pose $z = \expo{i\theta}$, $\theta \in \mathbb{R}$. Calculer $g(z)$.
		      Quel est l'ensemble des images par $g$ des points du cercle de centre $O$ et de rayon $1$\,?
		\item R\'esoudre dans $\mathbb{C}^*$ l'\'equation $g(z) = 1$ et v\'erifier le r\'esultat \`a l'aide de la question pr\'ec\'edente.
	\end{enumerate}
\end{exo}
%%%%%%%%%%%%%%%%%%%%%%%%%%%%%%%%%%%%%%%%%%%%%%%%%%%%%%%%%%%%%%%%%%%%%%%%%%%%%%%%%%%%
\begin{exo}
      On consid\`ere le nombre complexe $u=1+i$.
      \begin{enumerate}
            \item Ecrire $u$ et $\bar{u}$ sous forme trigonom\'etrique.
                  \label{item:1}
            \item Soit $n$ un entier naturel.
                  On pose $S_n = u^n + \bar{u}^n$.
                  D\'eduire de \ref{item:1} que $S_n = \lambda_n \cos \left( n \frac{\pi}{4} \right)$
                  o\`u $\lambda_n$ est un r\'eel \`a pr\'eciser en fonction de $n$.
            \item Pour quelles valeurs de $n$ a-t-on $S_n = 0$\,?
            \item Prouver que si $n$ est pair, $S_n$ est un entier relatif.
      \end{enumerate}
\end{exo}
%
\begin{exo}
      On pose $u=-1 + i \sqrt{3}$.
      \begin{enumerate}
            \item Donner la forme trigonom\'etrique des nombres :
                  $u$, $\bar{u}$, $-2u$, $iu$, $\frac{1}{u}$, $\frac{-i}{u^2}$
                  et construire leurs images dans le plan complexe.
            \item R\'esoudre dans $\mathbb{C}$ l'\'equation : $z^2 - 2i z - i\sqrt{3} = 0$.
      \end{enumerate}
\end{exo}
%
\begin{exo}
      R\'esoudre dans $\mathbb{C}$ les \'equations suivantes : 
      \begin{enumerate}
            \item $z^{2} - 2 \, z + 5 = 0$,
            \item $z^{2} - 4 \, z + 6 = 0$,
            \item $z^{2} - \left(1 + i\right) \, z + 5 i = 0$,
            \item $z^{2} + \left(1 + 3 i\right) \, z - 4 + 3 i = 0$,
            \item $z^{2} - ( 3 +  2 i \, \sqrt{2} )z + 3 i \, \sqrt{2} = 0$,
            \item $2 \, z^{2} - \left(6 - 5i\right) \, z + 7 - 4 i = 0$,
            \item $z^{2} - \left(3 - 4 i\right) \, z - 1 - 7 i = 0$.
      \end{enumerate}
\end{exo}
%
\begin{exo}
      En utilisant la formule de Moivre et la formule du bin\^ome, exprimer $\cos 3\theta$ en fonction de $\cos \theta$,
      puis $\sin 3\theta$ en fonction de $\sin \theta$.
      
      M\^eme questions pour $\cos 5\theta$ et $\sin 5\theta$.
\end{exo}
%
\begin{exo}
      On pose $z = \cos \theta + i \sin \theta$.
      \begin{enumerate}
            \item Calculer $z + \frac{1}{z}$, $z - \frac{1}{z}$, 
                  $z^n + \frac{1}{z^n}$ et $z^n - \frac{1}{z^n}$ ($n\in\mathbb{Z}$)
                  en fonction de $\theta$.
            \item D\'evelopper 
                  $\left( z + \frac{1}{z} \right)^3$,
                  $\left( z - \frac{1}{z} \right)^3$,
                  $\left( z + \frac{1}{z} \right)^4$ et
                  $\left( z - \frac{1}{z} \right)^4$ ; 
                  puis lin\'eariser $\cos^3 \theta$, $\sin^3 \theta$, $\cos^4 \theta$ et $\sin^4 \theta$.
      \end{enumerate}
\end{exo}
%
\begin{exo}
      \begin{enumerate}
            \item Lin\'eariser : $\cos ^2 x$, $\sin^3 x$, $\cos^4 x \cdot \sin x$ en utilisant les formules d'Euler.
            \item M\^eme question pour : $\cos^4 x$, $\sin^4 x$, $\cos^5 x$, $\sin^5 x$, $\cos^6 x$ et $\sin^6 x$.
      \end{enumerate}
\end{exo}
%
\begin{exo}
      R\'esoudre dans $\mathbb{C}^2$ les syst\`emes : 
            \begin{enumerate}						
                  \item $\displaystyle{ \left\lbrace \begin{array}{l} 
                              (1-i) x + (2-i) y = 6 + i \\
                              (2+i) x - (1+i) y = 5 - 8 i \;.
                         \end{array} \right. }$
                  \item $\displaystyle{ \left\lbrace \begin{array}{l} 
                              2ix + (1-3i) y = 7 - 3i \\
                              (1-i) \bar{x} - (5+2i) \bar{y} = -10 + i \;.
                         \end{array} \right. }$
                  \item $\displaystyle{ \left\lbrace \begin{array}{l} 
                              (1+i) x + 2 i y = -2 - 2i \\
                              (3-i) \bar{x} + (1+3i) \bar{y} = 3 - 11 i \;.
                         \end{array} \right. }$
						\item $\displaystyle{
									\begin{multiform}
										\left(2 + i\right) \, \bar{x} + \left(3 - 2 i\right) \, \bar{y} = 9 + 8 i			\\
										\left(4 - 3 i\right) \, x - \left(2 - 3 i\right) \, y = 7 - 12 i		.
									\end{multiform}	 
													}$
                  \item $\displaystyle{ \left\lbrace \begin{array}{l} 
                              (2+i) x + (2-i) \bar{y} = -12-4i \\
                              (3+2i) x - (3-2i) \bar{y} = 8-4i  \;.
                         \end{array} \right. }$
            \end{enumerate}
\end{exo}
%%%%%%%%%%%%%%%%%%%%%%%%%%%%%%%%%%%%%%%%%%%%%%%%%%%%%%%%%%%%%%%%%%%%%%%%%%%%%%%%%%%%
\begin{exo}
	R\'esoudre dans $\mathbb{C}$ les \'equations suivantes :
	\begin{enumerate}
		\item $z^2 + 4z + 13 = 0$,
		\item $z^2 - 2iz - 5 = 0$,
		\item $z^2 - (3 - 2i) z + 5 - 5i = 0$,
		\item $z^2 - (2 - 2i) z + 3 + 2i = 0$,
		\item $(1+i) z^2 - 2z + 4 = 0$.
	\end{enumerate}
	Dans chaque cas, v\'erifier que la somme et le produit des solutions valent $-\frac{b}{a}$ et $\frac{c}{a}$ (avec les notations $az^2 + bz + c = 0$).
\end{exo}
%
\begin{exo}
	\begin{enumerate}
		\item D\'eterminer les racines cubiques de $8i$ ; les donner sous forme exponentielle puis sous forme alg\'ebrique.
		\item D\'eterminer les racines sixi\`emes de $-64$ et placer leurs images dans le plan complexe.
		      Quelle figure forment-elles\,?
		\item R\'esoudre dans $\mathbb{C}$ les \'equations :
			\begin{enumerate}
				\item $z^3 + 27 = 0$,
				\item $z^4 = -8 + 8\sqrt{3}\, i$.
			\end{enumerate}
		\item Parmi les solutions de $z^4 = -8 + 8\sqrt{3}\, i$, lesquelles sont aussi des racines sixi\`emes de $-64$\,?
	\end{enumerate}
\end{exo}
%%%%%%%%%%%%%%%%%%%%%%%%%%%%%%%%%%%%%%%%%%%%%%%%%%%%%%%%%%%%%%%%%%%%%%%%%%%%%%%%%%%%
\begin{exo}
      \begin{enumerate}
            \item R\'esoudre dans $\mathbb{C}$ l'\'equation : $z^2 - (5+4i\sqrt{3}) z + 9 = 0$.
            \item On d\'esigne par $z'$ la racine ayant le plus grand module.
                  Calculer les modules et un argument de chacune des racines $z'$ et $z''$ 
                  ainsi que de $\frac{z'}{z''}$ et $z' z''$.
            \item Trouver les nombres complexes $Z$ v\'erifiant $Z^2 = z'$.
                  Exprimer les solutions sous forme alg\'ebrique, trigonom\'etrique puis polaire.
      \end{enumerate}
\end{exo}
%
\begin{exo}
      \begin{enumerate}
            \item Calculer $(1-2i)^3$ et $(2-i)^4$.
            \item En d\'eduire les solutions des \'equations : 
                  \begin{enumerate}
                        \item $z^3 = - 11 + 2i$
                        \item $z^4 + 7 + 24 i = 0$.
                  \end{enumerate}
      \end{enumerate}
\end{exo}
%
\begin{exo}
      R\'esoudre dans $\mathbb{C}$ les \'equations suivantes : 
      \begin{enumerate}
            \item $z^3 + z^2 + z + 1 = 0$,
            \item $\left(\frac{z-i}{z+i}\right)^3+\left(\frac{z-i}{z+i}\right)^2+\left(\frac{z-i}{z+i}\right)+1 = 0$,
            \item $1-z+z^2-z^3 = 0$,
            \item $z^4 = 2 - 2 i \sqrt{3}$.
      \end{enumerate}
\end{exo}
%
\begin{exo}
      On consid\`ere le polyn\^ome $P(z) = 2\, z^3 - (5+6i)\, z^2 + 9i\, z + 1 - 3i$.
      \begin{enumerate}
            \item Trouver le r\'eel $\alpha$ tel que $P(\alpha) = 0$.
            \item R\'esoudre dans $\mathbb{C}$ l'\'equation $P(z) = 0$.
      \end{enumerate}
      M\^eme questions pour le polyn\^ome $Q(z) = z^3 - (4+5i) z^2 - (1-19i) z + 12 - 12 i$.
      En d\'eduire les solutions de l'\'equation $z^3 - (4-5i) z^2 - (1+19i) z + 12 + 12 i = 0 $.
\end{exo}
%
\begin{exo}
      On consid\`ere le polyn\^ome $P(z) = z^3 - (3+2i)z^2 + (3+5i) z - (2+6i)$.
      \begin{enumerate}
            \item Trouver le r\'eel $\alpha$ tel que $P(i\alpha) = 0$.
            \item R\'esoudre dans $\mathbb{C}$ l'\'equation $P(z) = 0$.
      \end{enumerate}
      M\^eme questions pour les polyn\^omes
      $Q(z) = 2 \, z^{3} - 3 i \, z^{2} + \left(10 + 24 i\right) \, z + 36 - 15 i$
      et 
      $R(z) = 3 \, z^{3} - \left(3 + 5 i\right) \, z^{2} - \left(2 - 17 i\right) \, z + 10$.
\end{exo}
%
\begin{exo}
      On consid\`ere le polyn\^ome 
      $P(z) = \left(1-i\right) \, z^{3} + \left(1 + 3 i\right) \, z^{2} - \left(4 - 2 i\right) \, z + 12 - 4 i$.
      \begin{enumerate}
            \item Trouver le r\'eel $\alpha$ tel que $P(\alpha - i\alpha) = 0$.
            \item R\'esoudre dans $\mathbb{C}$ l'\'equation $P(z) = 0$.
      \end{enumerate}
      M\^eme questions pour le polyn\^ome
      $Q(z) = \left(1-i\right) \, z^{3} - \left(3-5 i\right) \, z^{2} + 18 \, z - 44 + 12 i$.
\end{exo}
%
\begin{exo}
      On pose $P(z) = z^4 - i z^3 - 3 (1-i) z^2 + 4 i z - 12 i - 4 $.
      \begin{enumerate}
            \item D\'emontrer que $P$ admet deux racines r\'eelles que l'on d\'eterminera.
            \item R\'esoudre dans $\mathbb{C}$ l'\'equation $P(z) = 0$.
      \end{enumerate}
\end{exo}
%
\begin{exo}
      R\'esoudre dans $\mathbb{C}$ les \'equations 
      \begin{enumerate}
            \item $ (2z^2-3z+2)^2 + (z^2-3z+2)^2 = 0 $,
            \item $ (z^2 + 4z + 1)^2 + (3z+5)^2 = 0 $,
            \item $ z^2 + 4z + 1 + i (3z + 5) = 0 $.
      \end{enumerate}
      (Remarquer que $a^2 + b^2 = (a+ib)(a-ib)$).
\end{exo}
%
\begin{exo}
      On pose $P(z) = z^4 - z^2 + 1$.
      \begin{enumerate}
            \item R\'esoudre dans $\mathbb{C}$ l'\'equation $P(z) = 0$.
            \item Montrer que $P(z)$ peut s'\'ecrire sous forme d'un produit de 
                  deux polyn\^omes du second degr\'e \`a coefficients r\'eels.
      \end{enumerate}
\end{exo}
%
\begin{exo}
      R\'esoudre dans $\mathbb{C}$ l'\'equation : $z^6 - (1-i) z^3 - i = 0$.
      Construire les images des solutions dans le plan complexe.
\end{exo}
%
\begin{exo}
      \begin{enumerate}
            \item D\'eterminer les modules et les arguments des solutions dans $\mathbb{C}$ de l'\'equation $x^5+1=0$.
            \item D\'emontrer que l'on peut \'ecrire : 
                  \[
                  x^5+1= (x+1) \left( x^2-2x\cos \frac{\pi}{5} +1 \right) \left( x^2 - 2x\cos \frac{3\pi}{5} +1 \right) .
                  \]
            \item En d\'eduire les valeurs de $\cos \frac{\pi}{5}$ et $\cos \frac{3\pi}{5}$.
      \end{enumerate}
\end{exo}
%
\begin{exo}
      On consid\`ere dans $\mathbb{C}$ la suite de terme g\'en\'eral $z_n$, 
      d\'efinie par son premier terme $z_0 = 1$ et par la relation de r\'ecurrence $2z_{n+1} = z_n + i$.
      \begin{enumerate}
            \item D\'emontrer que : $\forall n \in\mathbb{N}^*, |z_n|<1$.
            \item On pose $z_n = x_n + i y_n$ o\`u $x_n$ et $y_n$ sont r\'eels et $u_n = z_n - i$.
                  \begin{enumerate}
                        \item Trouver une relation entre $u_{n+1}$ et $u_n$.
                        \item En d\'eduire que la suite de terme g\'en\'eral 
                              $x_n$ est une suite g\'eom\'etrique qui converge vers $0$ 
                              et que les suite de termes g\'en\'eraux $y_n$ et $|z_n|$ convergent vers $1$.                  
                  \end{enumerate}
            \item Calculer le plus petit entier $n_0$ tel que : $\forall n \geq n_0,\ |z_n - i| < 10^{-5}$.
      \end{enumerate}
\end{exo}
%
\begin{exo}
      On consid\`ere la suite de nombre complexe d\'efinie par : 
            \[
                  z_0 = 0 ; z_1 = i ; z_n - z_{n-1} = i (z_{n-1} - z_{n-2} ) \quad (n\geq 2)
            \]
      \begin{enumerate}
            \item Exprimer $z_n - z_{n-1}$ en fonction de $i$ et de $n$.
            \item Etablir : $z_n = \frac{1-i}{2} (i^n - 1)$.
            \item D\'emontrer que cette suite est p\'eriodique.
            \item Repr\'esenter dans le plan complexe les points $M_{2017}$, $M_{2018}$, $M_{2019}$ et $M_{2020}$ 
                  d'affixes respectifs  $z_{2017}$, $z_{2018}$, $z_{2019}$ et $z_{2020}$.
      \end{enumerate}
\end{exo}
%
\begin{exo}
      On consid\`ere la suite de nombres complexes d\'efinie par son premier terme $z_0=8$ et la relation de r\'ecurrence :
            \[
                  z_{n+1} = \frac{\sqrt{2}+i\sqrt{6}}{4} \cdot z_n      .
            \]
      Dans le plan complexe muni d'un rep\`ere orthonormal direct $(O,\vec{\imath},\vec{\jmath})$,
      on d\'esigne par $M_n$ le point d'affixe $z_n$.
      \begin{enumerate}
            \item Donner la forme trigonom\'etrique et la forme alg\'ebrique des nombres       
                  $z_0$, $z_1$, $z_2$, $z_3$, $z_4$ et $z_5$
                  puis construire les points 
                  $M_0$, $M_1$, $M_2$, $M_3$, $M_4$ et $M_5$.
            \item Exprimer $z_n$ en fonction de $n$.
                  Montrer que les modules des nombres complexes $z_n$ forment une suite g\'eom\'etrique.
                  Que peut-on dire des arguments de $z_n$\,?
            \item D\'eterminer le plus petit entier $n$ tel que $|z_n| < 10^{-2}$.
                  Quel est le nombre $z_n$ correspondant\,?
                  Quelle est la limite de $z_n$ quand $n$ tend vers l'infini\,?
                  Que peut-on dire des points $M_n$\,?
            \item Comparer $OM_n$ et $OM_{n+1}$ puis calculer la mesure de l'angle 
                  $ \left( \overrightarrow{OM}_n,\overrightarrow{OM}_{n+1} \right)$.
                  En d\'eduire une construction de $M_{n+1}$ connaissant $M_n$.
      \end{enumerate}
\end{exo}
%
\begin{exo}
      Dans le plan complexe muni d'un rep\`ere orthonorm\'e direct $(O,\vec{u},\vec{v})$ d'unit\'e 2~cm,
      on consid\`ere les points $A$, $B$, $C$ et $D$ d'affixes respectifs 
            $z_A = -i$, $z_B = 3$, $z_C = 2+3i$ et $z_D = -1+2i$.
      \begin{enumerate}
            \item Placer sur une figure les points $A$, $B$, $C$ et $D$.
            \item \begin{enumerate}
                        \item Interpr\'eter g\'eom\'etriquement le module et l'argument du nombre complexe 
                              $\displaystyle{ \frac{z_C - z_A}{z_D - z_B} }$.
                        \item Calculer ce nombre complexe.
                        \item Que peut-on conclure concernant les segments $[AC]$ et $[BD]$\,?
                  \end{enumerate}
            \item \begin{enumerate}
                        \item Quelle est la nature du quadrilat\`ere $ABCD$\,? Justifier.
                        \item Calculer l'aire $s_0$ du quadrilat\`ere $ABCD$.
                  \end{enumerate}
      \end{enumerate}
\end{exo}
%
\begin{exo}
      Soit $P(z) = z^3 - (6+9i) z^2 + (-9+45i) z + 54 - 54 i$.
      \begin{enumerate}
            \item Montrer que l'\'equation $P(z) = 0$ admet une racine r\'eelle $z_0$ que l'on d\'eterminera.
                  En d\'eduire une factorisation de $P(z)$.
            \item D\'eterminer toutes les solutions de l'\'equation $P(z) = 0$.
                  On appellera $z_2$ la solution imaginaire pure, $z_0$ la solution r\'eelle et $z_1$ l'autre.
                  Pour chacune d'elles, pr\'eciser le module et l'argument.
            \item Montrer que $z_0$, $z_1$ et $z_2$ sont les trois premiers termes d'une suite g\'eom\'etrique
                  $(z_n)_{n\in\mathbb{N}}$ dont on pr\'ecisera la raison.
                  Calculer $z_n$ en fonction de $n$.
            \item Montrer que les modules des $z_n$ forment une suite g\'eom\'etrique.
                  Que peut-on dire de leurs arguments\,?
            \item Pour quelles valeurs de $n$, $z_n$ est-il r\'eel\,?
                  R\'eel n\'egatif\,?
                  Imaginaire pur\,?
      \end{enumerate}
\end{exo}
%
\begin{exo}
      On consid\`ere le polyn\^ome $P(z)$ suivant : $P(z) = z^3 + 9 i z^2 + 2(6i - 11)z - 3(4i + 12)$.
      \begin{enumerate}
            \item D\'emontrer que l'\'equation $P ( z ) = 0$ admet une solution r\'eelle $z_1$.
            \item D\'eterminer un polyn\^ome $Q ( z )$ tel que $P ( z ) = ( z - z_1 ) Q ( z )$.
            \item D\'emontrer que l'\'equation $Q ( z ) = 0$ admet une solution imaginaire pure $z_2$.
            \item R\'esoudre dans $\mathbb{C}$ l'\'equation $P ( z ) = 0$.
            \item On note $z_3$ la 3~i\`eme solution de l'\'equation $P ( z ) = 0$. 
                  D\'emontrer que les points du plan complexe $A$, $B$ et $C$ d'affixes respectives $z_1$, $z_2$ et $z_3$, 
                  sont align\'es.
      \end{enumerate}
\end{exo}
%%%%%%%%%%%%%%%%%%%%%%%%%%%%%%%%%%%%%%%%%%%%%%%%%%%%%%%%%%%%%%%%%%%%%%%%%%%%%%%%%%%
\begin{exo}
      On rapporte le plan \`a un rep\'ere orthonorm\'e $(O ; \vec{u} , \vec{v})$ d'unit\'e 2~cm.
      On consid\`ere le polyn\^ome $P$ d\'efinie par 
            $
                  P(z) = 4 \, z^{3} - \left(4 - 2 i\right) \, z^{2} + \left(5 + 2 i\right) \, z + 4 + \frac{1}{2} i \;.
            $
      \begin{enumerate}
            \item Trouver le r\'eel $\alpha$ v\'erifiant $P(\alpha) = 0$.
            \item R\'esoudre dans $\mathbb{C}$ l'\'equation $P(z) = 0$. \newline
                  On d\'esigne par $z_1$ et $z_2$ les deux racines non-r\'eels de $P$, o\`u $\mathfrak{Im}(z_1) < 0$.
                  On d\'esigne respectivement par $A$, $B$ et $C$ les points d'affixes $\alpha$, $z_1$ et $z_2$.
            \item Trouver les deux nombres complexes $a$ et $b$ v\'erifiant :
                  $
                        \left\lbrace
                              \begin{array}{l}
                                    \alpha = a z_2    + b \\
                                    z_1    = a \alpha + b \;.
                              \end{array}
                        \right.
                  $
                  
                  On pose $\omega = \frac{b}{1-a}$, et on note par $\Omega$ le point d'affixe $\omega$.
            \item D\'eterminer la forme alg\'ebrique de $\omega$.
            \item Placer les points $A$, $B$, $C$ et $\Omega$ dans le rep\`ere $(O ; \vec{u} , \vec{v})$.
      \end{enumerate}
\end{exo}
%%%%%%%%%%%%%%%%%%%%%%%%%%%%%%%%%%%%%%%%%%%%%%%%%%%%%%%%%%%%%%%%%%%%%%%%%%%%%%%%%%%
\begin{exo}
      Le plan complexe est rapporté au repère orthonormal direct $( A ; \overrightarrow{AB} , \overrightarrow{AD} )$. 
      On consid\`ere le point $C$ tel que $ABCD$ est un carr\'e.
      Soit E le milieu du segment [ AD ], on considère le carré EDGF tel que
      $( \overrightarrow{ED} ; \overrightarrow{EF} ) = \frac{\pi}{2} \quad [ \mathrm{mod} 2 \pi ] $.
      \begin{enumerate}
            \item \begin{enumerate}
                        \item Faire une figure en plaçant les points 
                              $A$, $B$, $C$, $D$, $E$, $F$, $G$. 
                              On complètera la figure au cours de l’exercice.
                        \item Préciser les nombres complexes $a$, $b$, $c$, $d$, $e$, $f$, $g$, 
                              affixes respectives des points $A$, $B$, $C$, $D$, $E$, $F$ et $G$. 
                  \end{enumerate}
                  Soit la similitude directe $s$ du plan telle que $s ( D ) = F$ et $s ( B ) = D$.
            \item On se propose de préciser les éléments caractéristiques de la similitude directe $s$.
                  \begin{enumerate}
                        \item Déterminer son rapport $k$ et son angle $\theta$.
                        \item Donner l’écriture complexe de cette similitude.
                        \item Déterminer le centre $\Omega$ de la similitude directe $s$.
                  \end{enumerate}
      \end{enumerate}
\end{exo}
%
\begin{exo}
      Le plan complexe est muni d'un repère orthonormal direct 
      $( O ; \vec{u} , \vec{v} )$, unité graphique 2~cm. 
      On note $A$, $B$, $C$, $D$ et $E$ les points d’affixes respectives 
      $z_A = 2 i$, $z_B = 2$, $z_C = 4 + 6i$, $z_D = -1 + i$ et $z_E = -3 + 3i$.
      \begin{enumerate}
            \item Placer les points sur une figure qui sera complétée au fur et à mesure des questions.
            \item Déterminer la nature du triangle $ABC$.
            \item Soit $f$ la similitude plane directe telle que $f(A) = D$ et $f(B) = A$.
                  \begin{enumerate}
                        \item Donner l’écriture complexe de $f$.
                        \item Déterminer l’angle, le rapport et le centre $\Omega$ de cette similitude.
                        \item Montrer que le triangle $DAE$ est l’image du triangle $ABC$ par la similitude $f$.
                        \item En déduire la nature du triangle $DAE$.
                  \end{enumerate}
            \item On désigne par $(\mathscr{C}_1)$ le cercle de diamètre $[AB]$ et
                  par $(\mathscr{C}_2)$ le cercle de diamètre $[AD]$.
                  On note $M$ le second point d’intersection du cercle $(\mathscr{C}_1)$ et de la droite $(BC)$, 
                  et $N$ le second point d’intersection du cercle $(\mathscr{C}_2)$ et de la droite $(AE)$.
                  \begin{enumerate}
                        \item Déterminer l’image de $M$ par la similitude $f$.
                        \item En déduire la nature du triangle $\Omega MN$.
                  \end{enumerate}
      \end{enumerate}
\end{exo}
%
\begin{exo}
      On consid\`ere le polyn\^ome $P$ d\'efini par 
      $P(z) = 2 \, z^{3} - \left(3 + i\right) \, z^{2} + \left( 3 - 3 i \right) \, z + 10 i$.
      \begin{enumerate}
            \item Trouver le r\'eel $\lambda$ tel que $P(\lambda + \lambda i) = 0$.
            \item Soient $A$, $B$ et $C$ les points d'affixes respectifs 
                  $z_A = -1-i$, $z_B = 1+2i$ et $z_C = \frac{3}{2}-\frac{1}{2}i$.
                  Soit $s$ la similitude plane directe qui laisse $A$ invariant et transforme $B$ en $C$.
                  \begin{enumerate}
                        \item Donner l'ecriture complexe de $s$.
                        \item Donner les \'el\'ements caract\'eristiques de $s$.
                  \end{enumerate}
            \item Soit $z\in\mathbb{C}$, on consid\`ere le nombre complexe $Z = \frac{z - z_B}{z - z_A}$.
                  \begin{enumerate}
                        \item Interpr\'eter $|Z|$ et $\Arg{Z}$.
                        \item D\'eterminer l'ensemble $(\mathscr{D})$ des points $M$ du plan d'affixe $z$ tels que 
                              $|\frac{z - z_B}{z - z_A}| = 1$.
                        \item D\'eterminer l'image de $(\mathscr{D})$ par $s$.
                  \end{enumerate}
      \end{enumerate}
\end{exo}
%%%%%%%%%%%%%%%%%%%%%%%%%%%%%%%%%%%%%%%%%%%%%%%%%%%%%%%%%%%%%%%%%%%%%%%%%%%%%%%%
\begin{exo}
	On consid\`ere le polyn\^ome $P(z) = z^3 - (4+4i) z^2 + (2+8i) z - 12 i$.
	\begin{enumerate}
		\item Calculer $P(2i)$.
		\item D\'eterminer les nombres complexes $a$, $b$ et $c$ tels que $P(z) = (z - 2i)(a z^2 + b z + c)$.
		\item R\'esoudre dans $\mathbb{C}$ l'\'equation $P(z) = 0$.
		\item On note $A$, $B$ et $C$ les points d'affixes $2i$, $1 - i$ et $3 + 3i$.
		      Calculer $\frac{z_C - z_A}{z_B - z_A}$ et en d\'eduire la nature du triangle $ABC$.
	\end{enumerate}
\end{exo}
%
\begin{exo}
	On consid\`ere le polyn\^ome $P(z) = z^3 - (2+i) z^2 + (-11 + 2i) z + 120 + 35i$.
	\begin{enumerate}
		\item Soit $x$ un r\'eel. \'Ecrire la partie r\'eelle et la partie imaginaire de $P(x)$.
		      En d\'eduire que l'\'equation $P(z) = 0$ admet une unique solution r\'eelle $z_0$, que l'on d\'eterminera.
		\item D\'eterminer les nombres complexes $a$ et $b$ tels que $P(z) = (z - z_0)(z^2 + a z + b)$.
		\item R\'esoudre dans $\mathbb{C}$ l'\'equation $P(z) = 0$.
		\item Montrer que les images des trois solutions sont sur un m\^eme cercle de centre $O$, dont on donnera le rayon.
		      On note $B$ et $C$ les images des deux solutions non r\'eelles : quelle est la nature du triangle $OBC$\,?
	\end{enumerate}
\end{exo}
%
\begin{exo}
	Soit le polyn\^ome $P(z) = z^3 - z^2 + 5z - 9 - 12i$.
	\begin{enumerate}
		\item V\'erifier que $z_A = 2 + i$ est une racine de $P$.
		\item Trouver les nombres complexes $a$ et $b$ tels que $P(z) = (z - z_A)(z^2 + a z + b)$.
		\item R\'esoudre dans $\mathbb{C}$ l'\'equation $P(z) = 0$.
		      On note $z_B$ la solution de partie r\'eelle n\'egative et $z_C$ la troisi\`eme solution.
		\item V\'erifier que $z_B = i z_A$. Interpr\'eter g\'eom\'etriquement.
		\item D\'eterminer l'affixe du point $D$ tel que $ABCD$ soit un parall\'elogramme, ainsi que celle de son centre.
	\end{enumerate}
\end{exo}
%
\begin{exo}
	On consid\`ere les \'equations
	\[
		(E) : \ z^4 - (2+3i) z^3 + (3+6i) z^2 + (4-15i) z - 10 = 0
		\qquad \text{et} \qquad
		(E') : \ z^2 - 2z + 5 = 0 .
	\]
	\begin{enumerate}
		\item R\'esoudre dans $\mathbb{C}$ l'\'equation $(E')$.
		\item D\'eterminer les nombres complexes $a$ et $b$ tels que, pour tout $z$,
		      $z^4 - (2+3i) z^3 + (3+6i) z^2 + (4-15i) z - 10 = (z^2 - 2z + 5)(z^2 + a z + b)$.
		\item R\'esoudre dans $\mathbb{C}$ l'\'equation $(E)$.
		\item Calculer la somme des quatre solutions de $(E)$ et la comparer au coefficient de $z^3$.
	\end{enumerate}
\end{exo}
%
\begin{exo}
	On consid\`ere l'\'equation $(E) : \ z^4 + (2 - 16i) z^2 - 63 + 16i = 0$.
	\begin{enumerate}
		\item On pose $Z = z^2$. \'Ecrire l'\'equation $(E')$ v\'erifi\'ee par $Z$ et la r\'esoudre
		      (on v\'erifiera que $(8 - 8i)^2 = -128\,i$).
		\item D\'eterminer les racines carr\'ees de $3 + 4i$ et de $-5 + 12i$.
		\item En d\'eduire les solutions de $(E)$.
		\item Montrer que les images des solutions de $(E)$ sont les sommets d'un parall\'elogramme de centre $O$.
	\end{enumerate}
\end{exo}
%%%%%%%%%%%%%%%%%%%%%%%%%%%%%%%%%%%%%%%%%%%%%%%%%%%%%%%%%%%%%%%%%%%%%%%%%%%%%%%%
\begin{exo}
      Le plan complexe est rapporté à un repère orthonormal direct $(O ; \vec{u} , \vec{v})$.
      Soit $f$ la fonction de $\mathbb{C}$ dans $\mathbb{C}$ définie par 
            $f(z) = \frac{2 - i z}{1-z}$.
      \begin{enumerate}
            \item Pr\'eciser l'ensemble de d\'efinition de $f$.
            \item \begin{enumerate}
                        \item Calculer $a = f(2)$ et donner la forme trigonom\'etrique de $a$.
                        \item Pour quelles valeurs de l'entier naturel $k$ le nombre $a^k$ est-il r\'eel\,?
                  \end{enumerate}
            \item On pose $z' = f(z)$.
                  \begin{enumerate}
                        \item Exprimer $z$ en fonction de $z'$.
                        \item Soient le point $M$ d'affixe $z$ et $M'$ d'affixe $z'$.
                              Montrer que $OM = \frac{M'A}{M'B}$
                              o\`u $A$ et $B$ sont des points dont on d\'eterminera les affixes.
                        \item Soit $(\mathscr{C})$ l'ensemble des points $M$ du plan pour lesquels on a 
                              $|z| = 1$ et $z \neq 1$.
                              Montrer que si $M \in (\mathscr{C})$ alors $M'$ appartient \`a une droite $(\mathscr{D})$ que l'on
                              d\'efinira g\'eom\'etriquement.
                        \item Tracer $(\mathscr{C})$ et $(\mathscr{D})$.
                  \end{enumerate}
            \item Soit $S$ la transformation du plan qui, \`a tout point $M(z)$, associe le point $M'(z')$
                  v\'erifiant la relation $z' = \frac{\sqrt{3} + i}{2} \, z + 1 + 4 i $.
                  \begin{enumerate}
                        \item D\'eterminer la nature de $S$.
                        \item D\'eterminer l'ensemble des points invariants par $S$.
                        \item En d\'eduire les \'el\'ements caract\'eristiques de $S$.
                  \end{enumerate}
      \end{enumerate}
\end{exo}





\begin{corrige}

\section{Solution des exercices et preuve des propri\'et\'es}
\begin{sol}
Solution de l'Exercice \ref{exe:mettre sous forme algebrique}. 
\begin{itemize}
	\item Pour le cas de $z_1$, il suffit de developper et de r\'eunir les termes r\'eels et les termes imaginaires: 
		\begin{align*}
		      z_1 & = (1-i)^2 (3+2i) = (1-2i+i^2) (3+2i)  \\
		          & = -2i (3+2i) = -6i - 4i^2 = 4 - 6i \;.
		\end{align*} 
		Ainsi on a $\mathfrak{Re}(z_1) = 4$ et $\mathfrak{Im}(z_1) = -6$.
	\item Pour le cas de $z_2$, une technique assez commune est de multiplier la fraction par le conjugu\'e du d\'enominateur: 
		\begin{align*}
		      z_2 & = \frac{10-20i}{1-3i} = \frac{(10 - 20 i) (1 + 3i)}{(1-3i) (1+3i)} 
		            = \frac{10 + 30i - 20i - 60 i^2}{1 + 3^2} \\
		          & = \frac{10 + 10 i + 60}{10} = 7 + i \;.
		\end{align*}
		Ainsi on a $\mathfrak{Re}(z_2) = 7$ et $\mathfrak{Im}(z_2) = 1$.
	\item Pour $z_3$, il suffit de combiner ce qu'on a fait dans le premier et dans le second cas:
		\begin{align*}
		      z_3 & = \frac{(10+15i)(10-20i)}{(3+i)^2} = \frac{(400-50i)(3-i)^2}{((3+i)(3-i))^2} \\
		          & = \frac{(400-50i)(8-6i)}{10^2} = \frac{2900 - 2800 i}{100} 
		            = 29 - 28 i \; . 
		\end{align*}
		Ainsi on a $\mathfrak{Re}(z_3) = 29$ et $\mathfrak{Im}(z_3) = -28$.
\end{itemize}
\label{sol:mettre sous forme algebrique}
\end{sol}
%%%%%%%%%%%%%%%%%%%%%%%%%%%%%%%%%%%%%%%%%%%%%%%%%%%%%%%%%%%%%%%%%%%%%%%%%%%%%%%%%%%%%%%%%%%%%%%%
\begin{dem}
      \label{dem:proprietes du conjugue}
      Preuve de la Proposition~\ref{prp:proprietes du conjugue}. 
      \begin{enumerate}
            \item Si $z = a + i b$ alors $\bar{z} = a - i b$ et $\bar{\bar{z}} = a + i b = z$.
            \item En posant $z_1 = a_1 + i b_1$ et $z_2 = a_2 + i b_2$, on a 
                  \begin{align*}
                        \overline{z_1 + z_2} = \overline{(a_1 + i b_1) + (a_2 + i b_2)} 
                              & = \overline{(a_1 + a_2) + i (b_1 + b_2)}                          \\
                              & = (a_1 + a_2) - i (b_1 + b_2)                                     \\
                              & = a_1 + a_2 - i b_1 - i b_2                                       \\
                              & = a_1 - i b_1 + a_2 - i b_2                                       \\
                              & = \bar{z_1} + \bar{z_2}                                           \;.
                  \end{align*}
            \item Avec les m\^eme notations que pr\'ec\'edemment :
                  \begin{align*}
                        \overline{z_1 \cdot z_2} = \overline{(a_1 + i b_1) \cdot (a_2 + i b_2)}
                              & = \overline{ a_1 a_2 + i a_1 b_2 + i a_2 b_1 + i^2 b_1 b_2 }       \\
                              & = \overline{ a_1 a_2 - b_1 b_2 + i (a_1 b_2 + a_2 b_1) }           \\
                              & = (a_1 a_2 - b_1 b_2) - i (a_1 b_2 + a_2 b_1)                      \;.
                  \end{align*}
                  %
                  \begin{align*}
                        \overline{z_1} \cdot \overline{z_2} = \overline{ a_1 + i b_1 } \cdot \overline{ a_2 + i b_2 }
                              & = ( a_1 - i b_1 ) \cdot ( a_2 - i b_2 )                            \\
                              & = (a_1 a_2 - i a_1 b_2 - i a_2 b_1 + i^2 b_1 b_2)                  \\
                              & = (a_1 a_2 - b_1 b_2) - i (a_1 b_2 + a_2 b_1)                      \;.
                  \end{align*}
                  %
                  Ce qui montre que $\overline{z_1 \cdot z_2} = \overline{z_1} \cdot \overline{z_2}$.
            \item Comme
                  \begin{align*}
                        \frac{1}{z}
                              & =  {\frac{1}{a + i b}}                                          
                                =  {\frac{a - i b}{(a+ib)(a-ib)}}                               
                                =  {\frac{a}{a^2 + b^2} - i \frac{b}{a^2 + b^2}}                        \\
                        \intertext{ on a  }
                        \overline{\frac{1}{z}}
                              & = \frac{a}{a^2 + b^2} + i \frac{b}{a^2 + b^2}                           
                                = \frac{a + i b}{a^2 + b^2}                                             
                                = \frac{a+ib}{(a-ib)(a+ib)}                                             \\
                              & = \frac{1}{a-ib}                                                        
                                = \frac{1}{\bar{z}}                                                     \;.
                  \end{align*}
            \item En utilisant la propri\'et\'e \ref{prp:proprietes du conjugue, multiplication}, on a :
                  $\overline{z^n} = \overline{ \underbrace{z \cdot z \cdots z}_{n \text{ fois }} }
                        = \underbrace{ \bar{z} \cdot \bar{z} \cdots \bar{z} }_{n \text{ fois }} 
                        = {\bar{z}}^n$.
      \end{enumerate}
\end{dem}
%%%%%%%%%%%%%%%%%%%%%%%%%%%%%%%%%%%%%%%%%%%%%%%%%%%%%%%%%%%%%%%%%%%%%%%%%%%%%%%%%%%%%%%%%%%%%%%%
\begin{dem}
	\label{dem:proprietes du module}
	Preuve de la Proposition~\ref{prp:proprietes du module}. 
	\begin{enumerate}
	      \item $|z| = \sqrt{a^2 + b^2} = \sqrt{a^2 + (-b)^2} = |\bar{z}|$,
	      \item $|z| = 0 \iff a^2 + b^2 = 0 \iff a = b = 0 \iff z = 0$,
	      \item $|z z'| = \sqrt{(zz') (\overline{zz'})} = \sqrt{(z \bar{z}) (z' \overline{z'})} 
	            = \sqrt{z \bar{z}} \sqrt{z' \overline{z'}} = |z| |z'| $,
	      \item Si $z \neq 0$ on a: $|z \cdot \frac{1}{z}| = 1$ ce qui implique que 
	            $|z| \cdot |\frac{1}{z}| = 1$ et donc que $|\frac{1}{z}| = \frac{1}{|z|}$.
	      \item D'abord, montrons l'in\'egalit\'e suivante :
		            \begin{lem}
		                  \label{lem: inegalite x**2 + y**2 >= 2xy}
	                        Pour tout r\'eels $X$ et $Y$, $2 X Y \leq X^2 + Y^2$.
	      	            \end{lem}
	      	            %
	      	            \begin{proof}
	      	                  Comme, pour tout r\'eel $X$ et $Y$, on a $(X - Y)^2 \geq 0$, on a 
	      	                  $X^2 - 2 X Y + Y^2 \geq 0$ et donc $X^2 + Y^2 \geq 2 X Y $.
	      	            \end{proof}
	      	      %
	      	      Posons $z = a + i b$ et $z' = \alpha + i \beta$, o\`u $a$, $b$, $\alpha$ et $\beta$ sont des nombres r\'eels.
	      	      Les deux membres \'etant positifs, l'in\'egalit\'e \`a d\'emontrer \'equivaut \`a chacune des suivantes :
	            \begin{align*}
	                  |z + z'|                   
	                        & \leq 
	                              |z| + |z'|                                                                       \\
	                  |z + z'|^2                 
	                        & \leq 
	                              (|z| + |z'|)^2                                                                   \\
	                  |z + z'|^2                 
	                        & \leq 
	                              |z|^2 + |z'|^2 + 2 |z|\cdot |z'|                                                 \;.
	            \end{align*}
	            D'une part, le premier membre de l'in\'egalit\'e est : 
	                  \[
	                  (a+\alpha)^2 + (b+\beta)^2 
	                  \]
	            qui est \'egale \`a 
	                  \[
	                  a^2+\alpha^2+b^2+\beta^2+2(a\alpha+b\beta)                                                   \;.
	                  \]
	            D'autre part, le second membre de l'in\'egalit\'e est :
	                  \[
	                        a^2 +b^2 + \alpha^2 + \beta^2 + 2 \sqrt{(a^2+b^2)\cdot(\alpha^2+\beta^2)}
	                  \]
      	            qui est \'egale \`a 
	                  \[
	                        a^2+b^2+\alpha^2+\beta^2+2\sqrt{(a\alpha)^2+(b\beta)^2+(a\beta)^2+(b\alpha)^2}         \;.
	                  \]
                  Apr\`es simplification par $a^2+\alpha^2+b^2+\beta^2$ et division par $2$, il s'agit de montrer que $a\alpha + b\beta \leq \sqrt{(a\alpha)^2+(b\beta)^2+(a\beta)^2+(b\alpha)^2}$. Si $a\alpha + b\beta < 0$, c'est vrai car une racine carr\'ee est positive. Sinon, les deux membres sont positifs et l'in\'egalit\'e \'equivaut \`a celle des carr\'es :
	            \begin{align*}
	                  (a\alpha + b\beta)^2
	                        & \leq 
	                              (a\alpha)^2 + (b\beta)^2 + (a\beta)^2 + (b\alpha)^2                              \\
	                  (a\alpha)^2 + (b\beta)^2 + 2 a\alpha b\beta 
	                        & \leq 
	                              (a\alpha)^2 + (b\beta)^2 + (a\beta)^2 + (b\alpha)^2                             \\
	                  2 a\beta \; b\alpha
	                        & \leq 
	                              (a\beta)^2 + (b\alpha)^2                                                        \;.
            	      \end{align*}
            	      Cette derni\`ere in\'egalit\'e est vraie d'apr\`es le Lemme~\ref{lem: inegalite x**2 + y**2 >= 2xy},
            	      en prenant $X = a\beta$ et $Y = b \alpha$.
	\end{enumerate} 
\end{dem}
%%%%%%%%%%%%%%%%%%%%%%%%%%%%%%%%%%%%%%%%%%%%%%%%%%%%%%%%%%%%%%%%%%%%%%%%%%%%%%%%%%%%%%%%%%%%%%%%
\begin{sol}
\label{sol:forme algebrique et modules}
Solution de l'exercice \ref{exe:forme algebrique et modules} 
\begin{itemize}
	\item Pour $z_1$:
		\begin{align*}
		      z_1 & = \overline{- (1 - 2i)^2} = \overline{(-1)} \cdot \overline{(1-2i)}^2 \\
		          & = -1 \cdot (1+2i)^2 = - (1 + 4i + 2^2 i^2) = - (4i - 3) = 3 - 4i \;.
		\end{align*}
		Ce qui nous permet de calculer $|z_1| = \sqrt{3^2 + (-4)^2} = \sqrt{25} = 5$.
	\item Pour $z_2$:
		\begin{align*}
		      z_2 & = \overline{\left( \frac{1}{3 + 4i} \right) } = \frac{1}{\overline{3 + 4i}}  = \frac{1}{3 - 4i}
		              \left( = \frac{1}{z_1} \right) \\
		          & = \frac{3+4i}{(3-4i)(3+4i)} = \frac{3+4i}{3^2+4^2} = \frac{3}{25} + \frac{4}{25}i \;.
		\end{align*}
		Pour calculer le module de $|z_2|$ on peut utiliser le fait que $z_2$ soit \'egal \`a $\frac{1}{z_1}$. 
		Nous avons alors $|z_2| = |\frac{1}{z_1}| = \frac{1}{|z_1|} = \frac{1}{5}$. \\
		On peut aussi faire le calcul direct (qui est plus compliqu\'e)
		\[
			|z_2| = \sqrt{\frac{3}{25}^2 + \frac{4}{25}^2} = \sqrt{\frac{3^2 + 4^2}{25^2} } 
			      = \sqrt{\frac{25}{25^2} } = \sqrt{\frac{1}{25}} = \frac{1}{5}      \;.
		\]
\end{itemize}
\end{sol}
%%%%%%%%%%%%%%%%%%%%%%%%%%%%%%%%%%%%%%%%%%%%%%%%%%%%%%%%%%%%%%%%%%%%%%%%%%%%%%%%%%%%%%%%%%%%%%%%
\begin{sol}
\label{sol:trigo-alg}
Solution de l'exercice \ref{exe:trigo-alg}. 
{\renewcommand{\arraystretch}{1.5}
	\begin{table}[h]
		\centering
		\begin{tabular}{IlIl|l|l|l|lI} \whline 
			$x$ & $\frac{\pi}{2} $ & $\frac{\pi}{3}$ & $\frac{\pi}{4}$ & $\frac{\pi}{6}$ & $0$ \\ \whline
			$\sin(x)$ & $1$ & $\frac{ \sqrt{3}}{2} $ & $\frac{ \sqrt{2}}{2} $ & $\frac{1}{2}$ & $0$ \\ \hline
			$\cos(x)$ & $0$ & $\frac{1}{2}$ & $\frac{ \sqrt{2}}{2}$ & $\frac{ \sqrt{3}}{2} $ & $1$ \\ \whline 
		\end{tabular}
		\label{tab:angles_remarquables}
		\caption{Valeurs des fonctions $\sin$ et $\cos$ pour certains angles remarquables.}
	\end{table}
}
\begin{enumerate}
	\item \begin{itemize}
			\item On commence par calculer $|z_1| = \sqrt{1^2 + 0^2} = 1 $ puis on le met en facteur :
			      $z_1 = 1 (1 + 0 i) = (\cos \theta + i \sin \theta)$.
			      Ainsi, 
			      $
			            \begin{multiform}
			                  \cos \theta = 1 \\
			                  \sin \theta = 0
			            \end{multiform}
			      $
			      ce qui correspond \`a $\theta = 0$.
			      Ainsi, $z_1  = 1 (\cos 0 + i \sin 0)$,
			      donc le module est \'egal \`a $1$ et l'argument est \'egal \`a $0$.
			\item On commence par calculer $|z_2| = \sqrt{0^2 + 1^2} = 1 $ puis on le met en facteur :
			      $z_2 = 1 (0 + 1\; i) = (\cos \theta + i \sin \theta)$.
			      Ainsi, 
			      $
			            \begin{multiform}
			                  \cos \theta = 0 \\
			                  \sin \theta = 1
			            \end{multiform}
			      $
			      ce qui correspond \`a $\theta = \frac{\pi}{2}$.
			      Ainsi, $z_2  = 1 (\cos \frac{\pi}{2} + i \sin \frac{\pi}{2})$, 
			      donc le module est \'egal \`a $1$ et l'argument est \'egal \`a $\frac{\pi}{2}$.
			\item On commence par calculer $|z_3| = \sqrt{2^2 + 2^2} = \sqrt{2 \cdot 2^2} = 2 \sqrt{2}$.
			      On met ensuite $|z_3|$ en facteur.
			      $z_3 = 2 \sqrt{2} \left(\frac{\sqrt{2}}{2} + i \frac{\sqrt{2}}{2} \right)$
			      qui doit \^etre \'egal \`a 
			      $ 2 \sqrt{2} \left( \cos \theta + i \sin \theta \right) $.
			      Il vient alors que 
			      $
			            \begin{multiform}
			                  \cos \theta = \frac{\sqrt{2}}{2} \\
			                  \sin \theta = \frac{\sqrt{2}}{2}
			            \end{multiform}
			      $
			      ce qui correspond \`a $\theta = \frac{\pi}{4}$.
			      Ainsi, $z_3 = 2 \sqrt{2} \left( \cos \frac{\pi}{4} + i \sin \frac{\pi}{4} \right) $, 
			           donc le module est \'egal \`a $2 \sqrt{2}$ et l'argument est \'egal \`a $\frac{\pi}{4}$.
			\item De m\^eme, on commence par calculer le module de $z_4$,
			      $|z_4| = \sqrt{1^2 + \sqrt{3}^2} = \sqrt{4} = 2$,
			      puis on le met en facteur.
			      $z_4 = 2 \left( \frac{1}{2} + i \frac{\sqrt{3}}{2} \right) = 2 \left( \cos \theta + i \sin \theta \right)$.
			      Donc 
			      $
			            \begin{multiform}
			                  \cos \theta = \frac{1}{2}  \\
			                  \sin \theta = \frac{\sqrt{3}}{2}
			            \end{multiform}
			      $
			      ce qui correspond \`a $\theta = \frac{\pi}{3}$.
			      Ainsi, $z_4 = 2 \left( \cos \frac{\pi}{3} + i \sin \frac{\pi}{3} \right) $, 
			      donc le module est \'egal \`a $2$ et l'argument est \'egal \`a $\frac{\pi}{3}$.
		\end{itemize}
	\item \begin{itemize} 
			\item On a $z'_1 = 2 \left( \cos \left(-\frac{\pi}{3}\right) + i \sin \left(-\frac{\pi}{3}\right) \right)$ :
			      le module est \'egal \`a $2$ et un argument est $-\frac{\pi}{3}$ (et non $\frac{\pi}{3}$, \`a cause du signe $-$).
			      On a 
			      $z'_1 = 2 \left( \cos \frac{\pi}{3} - i \sin \frac{\pi}{3} \right) 
			            = 2 \left( \frac{1}{2} - \frac{\sqrt{3}}{2} i \right) = 1 - \sqrt{3} i$.
			\item Le module est \'egal \`a $\sqrt{3}$ et l'argument est \'egal \`a $\frac{\pi}{6}$. 
			      On a 
			      $z'_2 = \sqrt{3} \left( \cos \frac{\pi}{6} + i \sin \frac{\pi}{6} \right)
			            = \sqrt{3} \left( \frac{\sqrt{3}}{2} + i \cdot \frac{1}{2} \right)
			            = \frac{3}{2} + \frac{\sqrt{3}}{2} i$.
		\end{itemize}
\end{enumerate}
\end{sol}
%%%%%%%%%%%%%%%%%%%%%%%%%%%%%%%%%%%%%%%%%%%%%%%%%%%%%%%%%%%%%%%%%%%%%%%%%%%%%%%%%%%%%%%%%%%%%%%%
\begin{dem}
\label{dem:argument}
Preuve de la proposition \ref{prop:argument}. 
\begin{enumerate}
	\item Soient $z = r (\cos \theta + i \sin \theta )$ et $z' = r' (\cos \theta' + i \sin \theta' )$.
	      On a:
	      \begin{align*}
		      z z' & = r r' (\cos \theta + i \sin \theta ) (\cos \theta' + i \sin \theta' )              \\
		           & = r r' \left[ (\cos \theta \cos \theta' - \sin \theta \sin \theta') 
		                 + i (\cos \theta \sin \theta' + \sin \theta \cos \theta') \right]    .
		      \intertext{Rappelons que \newline 
		                 $\cos (a+b) = \cos a \, \cos b - \sin a \, \sin b$ et que 
		                 $\sin (a+b) = \sin a \, \cos b + \sin b \, \cos a$. \newline On a alors: }
		      z z' & = r r' \left[ \cos (\theta + \theta') + i \sin (\theta + \theta') \right] .
	      \end{align*}
	      On en d\'eduit que $\arg(z z') = \arg(z) + \arg(z')$.
%	      \footnote{et on retrouve aussi la propri\'et\'e: $|z z'| = |z| |z'|$}.	
	      \begin{align*}
	            \arg(z^n) & = \arg( \underbrace{z \cdot z \cdot \cdots \cdot z}_{n\text{ fois}} )       \\
	                      & = \underbrace{ \arg{z} + \arg{z} + \cdots + \arg{z} }_{n\text{ fois}}
	                        = n\ \arg{z} \;.
	      \end{align*}
	\item Remarquons d'abord que $z \cdot \frac{1}{z} = 1$ et que $\arg(1) = 0$.
	      Puis on applique la propri\'et\'e pr\'ec\'edente \`a $z \cdot \frac{1}{z}$ 
	      et on a $0 = \arg(1) = \arg(z \cdot \frac{1}{z}) = \arg(z) + \arg(\frac{1}{z})$. 
	      Il s'en suit que $\arg(z) + \arg(\frac{1}{z}) = 0$ ce qui entraine que $\arg(\frac{1}{z}) = - \arg(z)$.
	      
	      \begin{align*}
	            \arg\left( \frac{z}{z'} \right) & = \arg\left( z \cdot \frac{1}{z'} \right)      \\
	                                            & = \arg z + \arg\left( \frac{1}{z'} \right)
	                                              = \arg z - \arg z'                              .
	      \end{align*}
	\item Soit $r (\cos \theta + i \sin \theta)$ l'expression trigonom\'etrique de $z$. 
	      Son expression alg\'ebrique est alors $r \cos \theta + i \, r \sin \theta$ et 
	      son conjugu\'e est $r \cos \theta - i \, r \sin \theta$.
	      Ainsi, le conjugu\'e de $z$ est $ r ( \cos \theta - i \, \sin \theta) $ 
	      qui est \'egal \`a $ r \left[ \cos (- \theta ) + i \, \sin (- \theta) \right]$.
	      On d\'eduit que $\arg(\bar{z}) = - \arg(z)$.
\end{enumerate}
Les deux derni\`eres propri\'et\'es sont faciles \`a v\'erifier (il suffit de d\'essiner).
\end{dem}
%%%%%%%%%%%%%%%%%%%%%%%%%%%%%%%%%%%%%%%%%%%%%%%%%%%%%%%%%%%%%%%%%%%%%%%%%%%%%%%%%%%%%%%%%%%%%%%%
\begin{dem}
	\label{dem:Moivre}
	Pour prouver la formule de Moivre (Th\'eor\`eme \ref{thm:Moivre}), nous allons utiliser le Th\'eoreme \ref{thm:Expo-trigo}.
	Puisque 
	\begin{align*}
	(\cos \theta + i \sin \theta)^n & = \left( \expo{i \theta} \right)^n         
	                                  = \expo{i n \theta}                        
	                                  = \cos (n \theta) + i \sin (n \theta)     \;.
	\end{align*}
\end{dem}
%%%%%%%%%%%%%%%%%%%%%%%%%%%%%%%%%%%%%%%%%%%%%%%%%%%%%%%%%%%%%%%%%%%%%%%%%%%%%%%%%%%%%%%%%%%%%%%%
\begin{dem}
      \label{dem:Expo-trigo}
      La formule d'Euler (Th\'eor\`eme~\ref{thm:Expo-trigo}) sera admis.
\end{dem}
%%%%%%%%%%%%%%%%%%%%%%%%%%%%%%%%%%%%%%%%%%%%%%%%%%%%%%%%%%%%%%%%%%%%%%%%%%%%%%%%%%%%%%%%%%%%%%%%
\begin{sol}
\label{sol:santatra second degre}
Solution de l'\'exercice \ref{exe:santatra second degre} \\
	\begin{itemize}
		\item On part de l'\'equation $z^2 - 2 z + 5 = 0$. 
			On a $\Delta = 2^2 - 4 \cdot 1 \cdot 5 = 4 - 20 = -16 = (4i)^2$.
			Ainsi, les deux racines complexes sont $z' = \frac{2-4i}{2} = 1 - 2i$ et $z'' = \frac{2+4i}{2} = 1 + 2i$. 
		\item On part de l'\'equation $z^2 - 4 z + 5 = 0$. 
			On a $\Delta = 4^2 - 4 \cdot 1 \cdot 5 = 16 - 20 = -4 = (2i)^2$.
			Ainsi, les deux racines complexes sont $z' = \frac{4-2i}{2} = 2 - i$ et $z'' = \frac{4+2i}{2} = 2 + i$. 
		\item On part de l'\'equation $4 z^2 - 8 z + 5 = 0$. 
			On a $\Delta = 8^2 - 4 \cdot 4 \cdot 5 = 64 - 80 = -16 = (4i)^2$.
			Ainsi, les deux racines complexes sont 
			$z' = \frac{8-4i}{2 \cdot 4} = 1 - \frac{1}{2} i$ et $z'' = \frac{8+4i}{2 \cdot 4} = 1 + \frac{1}{2} i$. 
	\end{itemize}
\end{sol}
%%%%%%%%%%%%%%%%%%%%%%%%%%%%%%%%%%%%%%%%%%%%%%%%%%%%%%%%%%%%%%%%%%%%%%%%%%%%%%%%%%%%%%%%%%%%%%%%
\begin{sol}
\label{sol:second degre complexe}
Solution de l'exercice \ref{exe:second degre complexe} \\
\begin{itemize} 
\item Comme $\Delta = (- (1 + 2 i) )^2 - 4 \cdot 1 \cdot (-1+i) = (-3+4i)-(4i-4) = 1 = 1^2$
      les solutions sont donc $z' = \frac{1+2i-1}{2} = i$ et $z'' = \frac{1+2i+1}{2} = 1+i$.
\item Comme $\Delta = 2^2 - 4 \cdot i \cdot (-2i) = 4-8 = -4 = (2i)^2$
      les solutions sont donc: \\
      $z' = \frac{-2-2i}{2i} = \frac{-1-i}{i} = \frac{-i+1}{-1} = $ \fbox{$-1+i$}
      \\ et \\
      $z'' = \frac{-2+2i}{2i} = \frac{-1+i}{i} = \frac{-i-1}{-1} = $ \fbox{$1+i$}.
\item Comme $\Delta = (1+3i)^2 + 4 (4 - 3i) = 1-9 + 6i + 16 - 12i = 8 - 6i$. 
      Soit $\delta = x + i y \in \mathbb{C}$ tel que $\delta^2 = \Delta$.
      Le module de $\Delta$ est \'egal \`a $\sqrt{8^2 + 6^2} = \sqrt{100} = 10$.
      
      On a alors 
      $\displaystyle{
	      \begin{multiform} 
		      x^2 - y^2 = 8       \\
		      2 x y     = -6      \\
		      x^2 + y^2 = 10      \;.
	      \end{multiform}
      }$
      
      En consid\'eant
      $\displaystyle{
	      \begin{multiform}
		      x^2 - y^2 = 8       \\
		      x^2 + y^2 = 10      
	      \end{multiform}
      }$
      on a $ 2 x^2 = 18 $ ce qui donne $x = \pm 3$.
      Si $x = 3$, on a $2 \cdot 3 \cdot y = -6$ ce qui donne $y = -1$.
      Ainsi, $\delta = 3 - i$ v\'erifie $\delta^2 = \Delta$.
      Les racines sont donc 
      \[
      z' = \frac{(1+3i) - (3-i)}{2} = -1 + 2i
      \qquad
      \text{ et }
      \qquad
      z'' = \frac{(1+3i) + (3-i)}{2} = 2 + i
      \; .
      \]
      L'ensemble des solutions est 
      \fbox{
		      $\displaystyle{
            		      S = \left\lbrace -1 + 2i \ ; \ 2 + i  \right\rbrace
		      }$
	      }
      .
\item Pour $z^2 + (1-2i) z - 3 - i = 0$ : $\Delta = (1-2i)^2 + 4(3+i) = (-3-4i) + 12 + 4i = 9 = 3^2$.
      Les solutions sont $z' = \frac{-(1-2i) - 3}{2} = -2 + i$ et $z'' = \frac{-(1-2i) + 3}{2} = 1 + i$.
\end{itemize}
\end{sol}

\begin{sol}
\label{sol:racines de l'unite}
Solution de l'exercice \ref{exe:racines de l'unite}. \\
Le module de $1$ est \'egal \`a $1$ et $\sqrt[n]{1} = 1$ pour tout entiers naturels positifs $n$.
Comme $1$ est un r\'eel, son argument est \'egal \`a $0$. 
Ainsi, les racines cubiques de l'unit\'e sont: 
\begin{align*}
      z_0 & = \sqrt[3]{1} \left(\cos \frac{0 + 0 \, \pi}{3} + i \sin \frac{0 + 0 \, \pi}{3} \right)        \\
      z_0 & = \cos 0 + i \sin 0                                                                            \\
          & \fbox{$ z_0 = 1 $}                                                                             \\
      z_1 & = \sqrt[3]{1} \left(\cos \frac{0 + 2 \, \pi}{3} + i \sin \frac{0 + 2 \, \pi}{3} \right)        \\
      z_1 & = \cos \frac{2\pi}{3} + i \sin \frac{2\pi}{3}                                                  \\
          & \fbox{$ z_1 = -\frac{1}{2} + i \frac{\sqrt{3}}{2} $}                                           \\
      z_2 & = \sqrt[3]{1} \left(\cos \frac{0 + 4 \, \pi}{3} + i \sin \frac{0 + 4 \, \pi}{3} \right)        \\
      z_2 & = \cos \frac{4 \pi}{3} + i \sin \frac{4 \pi}{3}                                                \\
          & \fbox{$ z_2 = -\frac{1}{2} - i \frac{\sqrt{3}}{2} $}                                           \;.
\end{align*}
De m\^eme, on peut trouver que les racines quatri\`emes de l'unit\'e sont $1$, $i$, $-1$ et $-i$.
%%%%%%%%%%%%%%%%%%%%%%%%%%%%%%%%%%%%%%%%%%%%%%%%%%%%%%%%%%%%%%%%%%%%%%%%%%%%%%%%%%%%%%%
\begin{figure}[H]
\centering
%%%%%%%%%%%%%%%%%%%%%%%%%%%%%%%%%%%%%%%%%%%%%%%%%%%%%%%%%%
\begin{subfigure}[b]{0.4 \actualtextwidth}
	%%%%%%%%%%%%%%%%%%%%%%%%%%%%%
	\begin{tikzpicture}[line cap=round,line join=round,>=triangle 45,x=1.5cm,y=1.5cm]
		\draw [->,color=black] (-1.5,0) -- (1.5,0);
		\draw [->,color=black] (0,-1.5) -- (0,1.5);
		\clip (-1.5,-1.5) rectangle (1.5,1.5);
		\draw [shift={(0,0)},color=brown,fill=brown,fill opacity=0.1] 
		      (0,0) -- (0:0.18) arc (0:120:0.18) -- cycle;
		\draw [shift={(0,0)},color=green!50!black,fill=green!50!black,fill opacity=0.1] 
		      (0,0) -- (-120:0.16) arc (-120:0:0.16) -- cycle;
		\draw [shift={(0,0)},color=green!50!black,fill=green!50!black,fill opacity=0.1] 
		      (0,0) -- (-120:0.2) arc (-120:0:0.2) -- cycle;
		\draw (0,0) circle (1);
		\draw (0,0)-- (-0.5,0.85);
		\draw (0,0)-- (-0.5,-0.85);
		\fill [color=gray] (0,0) circle (1.5pt);
		\fill [color=black] (1,0) circle (1.5pt);
		\draw [color=black] (1.2,0.2) node {$1$};
		\fill [color=black] (-0.5,0.85) circle (1.5pt);
		\draw [color=black] (-1,1) node {$-\frac{1}{2} + i \frac{ \sqrt{3} }{2}$};
		\fill [color=black] (-0.5,-0.85) circle (1.5pt);
		\draw [color=black] (-1,-1) node {$ - \frac{1}{2} - i\frac{\sqrt{3}}{2}$};
		\draw [color=brown] (0.35,0.25) node {$\frac{2 \pi}{3}$};
		\draw [color=green!50!black] (0.35,-0.25) node {$\frac{- 2 \pi}{3}$};
	\end{tikzpicture}
	\label{fig:racine 3eme de l'unite}
	\caption{Les racines cubiques de l'unit\'e forment un polygone r\'egulier \`a $3$ c\^ot\'e.}
	%%%%%%%%%%%%%%%%%%%%%%%%%%%%%
\end{subfigure}
%%%%%%%%%%%%%%%%%%%%%%%%%%%%%%%%%%%%%%%%%%%%%%%%%%%%%%%%%%
\hfill
%%%%%%%%%%%%%%%%%%%%%%%%%%%%%%%%%%%%%%%%%%%%%%%%%%%%%%%%%%
\begin{subfigure}[b]{0.4 \actualtextwidth}
%%%%%%%%%%%%%%%%%%%%%%%%%%%%%
	\definecolor{zzttqq}{rgb}{0.6,0.2,0}
	\definecolor{uququq}{rgb}{0.25,0.25,0.25}
	\begin{tikzpicture}[line cap=round,line join=round,>=triangle 45,x=1.5cm,y=1.5cm]
		\draw [->,color=black] (-1.5,0) -- (1.5,0);
		\draw [->,color=black] (0,-1.5) -- (0,1.5);
		\clip (-1.5,-1.5) rectangle (1.5,1.5);
		\draw (0,0) circle (1);
		\draw [color=brown] (0,1)-- (-1,0);
		\draw [color=brown] (-1,0)-- (0,-1);
		\draw [color=brown] (0,-1)-- (1,0);
		\draw [color=brown] (1,0)-- (0,1);
		\fill [ ] (0,0) circle (1.5pt);
		\draw [ ] (-0.15,-0.15) node {$O$};
		\fill [ ] (1,0) circle (1.5pt);
		\draw [ ] (1.1,0.1) node {$1$};
		\fill [ ] (0,1) circle (1.5pt);
		\draw [ ] (0.1,1.2) node {$i$};
		\fill [ ] (-1,0) circle (1.5pt);
		\draw [ ] (-1.2,-0.1) node {$-1$};
		\fill [ ] (0,-1) circle (1.5pt);
		\draw [ ] (-0.17,-1.15) node {$-i$};
	\end{tikzpicture}
	\label{fig:racine 4eme de l'unite}
	\caption{Les racines $4^{\text{i\`eme}}$ de l'unit\'e forment un polygone r\'egulier \`a $4$ c\^ot\'e.}
%%%%%%%%%%%%%%%%%%%%%%%%%%%%%
\end{subfigure}
%%%%%%%%%%%%%%%%%%%%%%%%%%%%%%%%%%%%%%%%%%%%%%%%%%%%%%%%%%
\caption{Racines $n^{\text{i\`eme}}$ de l'unit\'e.}
\label{fig:racine nieme de l'unite}
\end{figure}
%%%%%%%%%%%%%%%%%%%%%%%%%%%%%%%%%%%%%%%%%%%%%%%%%%%%%%%%%%%%%%%%%%%%%%%%%%%%%%%%%%%%%%%
\end{sol}


\begin{sol}
\label{sol:racine nieme}
Solution de l'exercice \ref{exe:racine nieme} \\
%%%%%%%%%%%%%%%%%%%%%%%%%%%%%%%%%%%%%%%%%%%%%%%%%%%%%%%%%%%%%%%%%%%%%%%%%%%%%%%%%%%%%%%
Posons $z = 1 + \sqrt{3} i = 2 \left( \frac{1}{2} + \frac{\sqrt{3}}{2} i \right)$; on a alors $|z| = 2$ et $\arg(z) = \frac{\pi}{3}$.
Les racines $5^{\text{i\`emes}}$ de $z$ sont les nombres 
$$
\sqrt[5]{2}  \left[ \cos \left( \frac{ \frac{\pi }{3} + 2k\pi}{5} \right) + i \sin \left( \frac{ \frac{\pi }{3} + 2k\pi}{5} \right) \right] 
$$ 
pour $k$ allant de $0$ \`a $4$.
Les racines $5^{\text{i\`eme}}$ de $z$ sont alors: 
\begin{enumerate}
\item $ z_0 = \sqrt[5]{2} \left( \cos \frac{\pi}{15} + i \sin \frac{\pi}{15} \right)         $
\item $ z_1 = \sqrt[5]{2} \left( \cos \frac{7 \pi}{15} + i \sin \frac{7 \pi}{15} \right)     $
\item $ z_2 = \sqrt[5]{2} \left( \cos \frac{13 \pi}{15} + i \sin \frac{13 \pi}{15} \right)   $
\item $ z_3 = \sqrt[5]{2} \left( \cos \frac{19 \pi}{15} + i \sin \frac{19 \pi}{15} \right)   $
\item $ z_4 = \sqrt[5]{2} \left( \cos \frac{5 \pi}{3} + i \sin \frac{5 \pi}{3} \right)       $
\end{enumerate}
%%%%%%%%%%%%%%%%%%%%%%%%%%%%%%%%%%%%%%%%%%%%%%%%%%%%%%%%%%
%\begin{figure}
%\includegraphics[scale=2]{./Complexe/Sary/5throot.png}
%\label{fig:racine 5ieme}
%\caption{Les racines $5^{\text{i\`eme}}$ forment un polygone r\'egulier centr\'e \`a l'origine.}
%\end{figure}
%%%%%%%%%%%%%%%%%%%%%%%%%%%%%%%%%%%%%%%%%%%%%%%%%%%%%%%%%%%%%%%%%%%%%%%%%%%%%%%%%%%%%%%
\end{sol}
%%%%%%%%%%%%%%%%%%%%%%%%%%%%%%%%%%%%%%%%%%%%%%%%%%%%%%%%%%%%%%%%%%%%%%%%%%%%%%%%%%%%%%%%%%%%%%%%
\begin{sol}
Solution de l'exercice \ref{exe:angle}.
$z_B - z_A = 3 + 3i$ et $z_C - z_A = -4 + 2i$, donc
\[
\frac{z_C - z_A}{z_B - z_A} = \frac{(-4+2i)(3-3i)}{(3+3i)(3-3i)} = \frac{-6 + 18i}{18} = -\frac{1}{3} + i .
\]
Un argument $\varphi$ de ce nombre est une mesure de $\left(\overrightarrow{AB},\overrightarrow{AC}\right)$ ; il v\'erifie
$\cos\varphi = \frac{-1/3}{\sqrt{10}/3} = -\frac{1}{\sqrt{10}}$ et $\sin\varphi = \frac{3}{\sqrt{10}} > 0$.
Donc $\hat{A} = \arccos\left(-\frac{1}{\sqrt{10}}\right) \approx 1{,}89$~rad, soit environ $108{,}4^\circ$.
\end{sol}
%%%%%%%%%%%%%%%%%%%%%%%%%%%%%%%%%%%%%%%%%%%%%%%%%%%%%%%%%%%%%%%%%%%%%%%%%%%%%%%%%%%%%%%%%%%%%%%%
\begin{sol}
Solution de l'exercice \ref{exe:rotation}.
L'image $B'$ de $B$ par la rotation de centre $A$ et d'angle $\theta$ a pour affixe $z_{B'} = \expo{i\theta}(z_B - z_A) + z_A$, avec $z_B - z_A = 1 - 3i$.
\begin{itemize}
\item $\theta = \frac{\pi}{3}$ : $z_{B'} = \left(\frac{1}{2} + i\frac{\sqrt{3}}{2}\right)(1 - 3i) + 2 + i
      = \boxed{\frac{5 + 3\sqrt{3}}{2} + i\,\frac{\sqrt{3} - 1}{2}}$.
\item $\theta = -\frac{\pi}{3}$ : $z_{B''} = \left(\frac{1}{2} - i\frac{\sqrt{3}}{2}\right)(1 - 3i) + 2 + i
      = \boxed{\frac{5 - 3\sqrt{3}}{2} - i\,\frac{1 + \sqrt{3}}{2}}$.
\end{itemize}
\end{sol}
%%%%%%%%%%%%%%%%%%%%%%%%%%%%%%%%%%%%%%%%%%%%%%%%%%%%%%%%%%%%%%%%%%%%%%%%%%%%%%%%%%%%%%%%%%%%%%%%
\subsection*{Solutions des exercices de fin de chapitre}
Dans ce qui suit, les arguments sont donn\'es \`a $2\pi$ pr\`es et $k$ d\'esigne un entier relatif.

\paragraph{Exercice 1.}
\begin{enumerate}
\item $z_1 = \sqrt{3} - i + i\sqrt{3} + 1 = \boxed{(\sqrt{3}+1) + i(\sqrt{3}-1)}$.
      $(1+i)^2 = 2i$, donc $(1+i)^4 = -4$ ; $(\sqrt{3}-i)^2 = 2 - 2\sqrt{3}i$ et $(\sqrt{3}-i)^3 = (2-2\sqrt{3}i)(\sqrt{3}-i) = -8i$ ;
      donc $z_2 = (-4)(-8i) = \boxed{32i}$.
      $z_3 = \frac{(1+i)(\sqrt{3}+i)}{(\sqrt{3}-i)(\sqrt{3}+i)} = \boxed{\frac{\sqrt{3}-1}{4} + i\,\frac{\sqrt{3}+1}{4}}$.
      Les parties r\'eelles et imaginaires se lisent sur ces \'ecritures.
\item $1 + i = \sqrt{2}\,\expo{i\frac{\pi}{4}}$ et $\sqrt{3} - i = 2\,\expo{-i\frac{\pi}{6}}$. Donc
      $z_1 = 2\sqrt{2}\left(\cos\frac{\pi}{12} + i\sin\frac{\pi}{12}\right)$,
      $z_2 = (\sqrt{2})^4 2^3\,\expo{i\left(\pi - \frac{\pi}{2}\right)} = 32\left(\cos\frac{\pi}{2} + i\sin\frac{\pi}{2}\right)$,
      $z_3 = \frac{\sqrt{2}}{2}\left(\cos\frac{5\pi}{12} + i\sin\frac{5\pi}{12}\right)$.
\end{enumerate}

\paragraph{Exercice 2.}
\begin{enumerate}
\item $|z'|^2 = (1+\sqrt{3})^2 + (1-\sqrt{3})^2 = 8$ et
      $z'' = \frac{z\,\overline{z'}}{|z'|^2} = \frac{(2\sqrt{3}+2i)\left(1+\sqrt{3} - (1-\sqrt{3})i\right)}{8} = \frac{8 + 8i}{8} = \boxed{1 + i}$.
\item $z'' = \sqrt{2}\left(\cos\frac{\pi}{4} + i\sin\frac{\pi}{4}\right)$ et $z = 4\left(\frac{\sqrt{3}}{2} + \frac{1}{2}i\right) = 4\left(\cos\frac{\pi}{6} + i\sin\frac{\pi}{6}\right)$.
\item $z' = \frac{z}{z''}$, donc $|z'| = \frac{4}{\sqrt{2}} = 2\sqrt{2}$ et $\arg z' = \frac{\pi}{6} - \frac{\pi}{4} = -\frac{\pi}{12}$.
      En identifiant $z' = 2\sqrt{2}\left(\cos\frac{\pi}{12} - i\sin\frac{\pi}{12}\right)$ avec sa forme alg\'ebrique :
      $\cos\frac{\pi}{12} = \frac{1+\sqrt{3}}{2\sqrt{2}} = \frac{\sqrt{6}+\sqrt{2}}{4}$,
      $\sin\frac{\pi}{12} = \frac{\sqrt{3}-1}{2\sqrt{2}} = \frac{\sqrt{6}-\sqrt{2}}{4}$ et
      $\tan\frac{\pi}{12} = \frac{\sqrt{3}-1}{\sqrt{3}+1} = \frac{(\sqrt{3}-1)^2}{2} = 2 - \sqrt{3}$.
\end{enumerate}

\paragraph{Exercice 3.}
\begin{enumerate}
\item $|p| = \sqrt{3+1} = 2$, $p = 2\left(\frac{\sqrt{3}}{2} + \frac{1}{2}i\right)$, donc $\arg p = \frac{\pi}{6}$.
      $|q| = \sqrt{8} = 2\sqrt{2}$, $q = 2\sqrt{2}\left(-\frac{\sqrt{2}}{2} + i\frac{\sqrt{2}}{2}\right)$, donc $\arg q = \frac{3\pi}{4}$.
\item $p^2$ : module $4$, argument $\frac{\pi}{3}$ ;
      $\frac{q}{p}$ : module $\sqrt{2}$, argument $\frac{3\pi}{4} - \frac{\pi}{6} = \frac{7\pi}{12}$ ;
      $\frac{1}{q}$ : module $\frac{1}{2\sqrt{2}} = \frac{\sqrt{2}}{4}$, argument $-\frac{3\pi}{4}$ ;
      $p^6$ : module $64$, argument $\pi$ ;
      $\frac{q^4}{p^3}$ : module $\frac{64}{8} = 8$, argument $3\pi - \frac{\pi}{2} \equiv \frac{\pi}{2}$.
\item $p^6 = 64(\cos\pi + i\sin\pi) = \boxed{-64}$ et $\frac{q^4}{p^3} = 8\left(\cos\frac{\pi}{2} + i\sin\frac{\pi}{2}\right) = \boxed{8i}$.
\item $\frac{q}{p} = \frac{(-2+2i)(\sqrt{3}-i)}{4} = \frac{(1-\sqrt{3}) + i(1+\sqrt{3})}{2}$.
      Comme $\frac{q}{p} = \sqrt{2}\left(\cos\frac{7\pi}{12} + i\sin\frac{7\pi}{12}\right)$, on obtient
      $\cos\frac{7\pi}{12} = \frac{1-\sqrt{3}}{2\sqrt{2}} = \frac{\sqrt{2}-\sqrt{6}}{4}$ et
      $\sin\frac{7\pi}{12} = \frac{1+\sqrt{3}}{2\sqrt{2}} = \frac{\sqrt{2}+\sqrt{6}}{4}$.
\end{enumerate}

\paragraph{Exercice 4.}
\begin{enumerate}
\item En d\'eveloppant, le premier membre vaut
      $a^2a'^2 - 2aa'bb' + b^2b'^2 + a^2b'^2 + 2ab'a'b + a'^2b^2 = a^2a'^2 + b^2b'^2 + a^2b'^2 + a'^2b^2$,
      qui est aussi le d\'eveloppement de $(a^2+b^2)(a'^2+b'^2)$.
      Or $zz' = (aa' - bb') + i(ab' + a'b)$, donc $|zz'|^2 = |z|^2|z'|^2$ et, les modules \'etant positifs, $|zz'| = |z|\,|z'|$.
\item \begin{align*}
      zz' & = rr'\left[(\cos\theta\cos\theta' - \sin\theta\sin\theta') + i(\sin\theta\cos\theta' + \cos\theta\sin\theta')\right] \\
          & = rr'\left[\cos(\theta+\theta') + i\sin(\theta+\theta')\right]
      \end{align*}
      avec $rr' > 0$ :
      c'est une forme trigonom\'etrique de $zz'$, donc $\theta+\theta'$ en est un argument.
\item Pour $n = 1$ c'est vrai. Si $|z^n| = |z|^n$ et si $n\theta$ est un argument de $z^n$, alors, d'apr\`es 1 et 2 appliqu\'ees \`a $z^n$ et $z$,
      $|z^{n+1}| = |z^n|\,|z| = |z|^{n+1}$ et $n\theta + \theta = (n+1)\theta$ est un argument de $z^{n+1}$.
\item Notons $w = \frac{z}{z'}$. De $z' w = z$ on tire $|z'|\,|w| = |z|$, soit $|w| = \frac{|z|}{|z'|}$,
      et $\arg z' + \arg w = \arg z$, soit $\arg w = \arg z - \arg z'$ (\`a $2\pi$ pr\`es).
\item $|3 + 4i| = 5$ et $|1 - i| = \sqrt{2}$, donc
      $\left|\frac{(3+4i)^5}{(1-i)^4}\right| = \frac{5^5}{(\sqrt{2})^4} = \boxed{\frac{3125}{4}}$.
\end{enumerate}

\paragraph{Exercice 5.}
\begin{enumerate}
\item $|z+z'|^2 = (z+z')\overline{(z+z')} = z\bar{z} + z'\overline{z'} + z\overline{z'} + \overline{z\overline{z'}}
      = |z|^2 + |z'|^2 + 2\,\mathfrak{Re}\left(z\overline{z'}\right)$,
      car $w + \bar{w} = 2\,\mathfrak{Re}(w)$.
\item Si $w = u + iv$, $\mathfrak{Re}(w) = u \leq |u| = \sqrt{u^2} \leq \sqrt{u^2+v^2} = |w|$.
      Donc $\mathfrak{Re}\left(z\overline{z'}\right) \leq |z\overline{z'}| = |z|\,|z'|$ et
      $|z+z'|^2 \leq |z|^2 + |z'|^2 + 2|z|\,|z'| = (|z| + |z'|)^2$ ; les deux nombres \'etant positifs, $|z+z'| \leq |z| + |z'|$.
\item $|z| = |(z-z') + z'| \leq |z - z'| + |z'|$, donc $|z| - |z'| \leq |z - z'|$.
      En \'echangeant les r\^oles, $|z'| - |z| \leq |z' - z| = |z - z'|$. D'o\`u $\left||z| - |z'|\right| \leq |z - z'|$.
\item Avec $w = 3 - 4i$, $|w| = 5$ : $|z + w| \leq |z| + |w| = 7$ et $|z + w| = |z - (-w)| \geq \left||z| - |w|\right| = 3$.
      Pour $z = \frac{2}{5}(3 - 4i) = \frac{6}{5} - \frac{8}{5}i$, on a $|z| = 2$ et $z + w = \frac{7}{5}(3-4i)$, de module $7$.
\end{enumerate}

\paragraph{Exercice 6.}
\begin{enumerate}
\item $\omega\bar{\omega} = |\omega|^2 = 1$, donc $\bar{\omega} = \frac{1}{\omega}$ ($\omega \neq 0$). Alors
      $\overline{Z} = \frac{\bar{\omega} - 1}{\bar{\omega} + 1} = \frac{\frac{1}{\omega} - 1}{\frac{1}{\omega} + 1} = \frac{1 - \omega}{1 + \omega} = -Z$.
      Si $Z = x + iy$, $\overline{Z} = -Z$ donne $x = -x$, donc $x = 0$ : $Z$ est imaginaire pur (\'eventuellement nul).
\item $\omega - 1 = iy(\omega + 1)$ donne $\omega(1 - iy) = 1 + iy$, et $1 - iy \neq 0$, donc
      $\omega = \frac{1 + iy}{1 - iy}$ et $|\omega| = \frac{\sqrt{1+y^2}}{\sqrt{1+y^2}} = 1$.
\item $Z = \frac{\expo{i\frac{\theta}{2}}\left(\expo{i\frac{\theta}{2}} - \expo{-i\frac{\theta}{2}}\right)}{\expo{i\frac{\theta}{2}}\left(\expo{i\frac{\theta}{2}} + \expo{-i\frac{\theta}{2}}\right)}
      = \frac{2i\sin\frac{\theta}{2}}{2\cos\frac{\theta}{2}} = i\tan\frac{\theta}{2}$
      (formules d'Euler ; $\cos\frac{\theta}{2} \neq 0$ car $\frac{\theta}{2} \in \left]-\frac{\pi}{2},\frac{\pi}{2}\right[$).
\item $Z = \frac{z_M - z_A}{z_M - z_B}$, donc un argument de $Z$ est une mesure de l'angle $\left(\overrightarrow{MB},\overrightarrow{MA}\right)$.
      Pour $\omega \neq \pm 1$, $Z$ est imaginaire pur non nul, donc cet angle mesure $\pm\frac{\pi}{2}$ : le triangle $AMB$ est rectangle en $M$
      (on retrouve que $M$ est sur le cercle de diam\`etre $[AB]$).
\end{enumerate}

\paragraph{Exercice 7.}
\begin{enumerate}
\item $g(z) - g(z') = (z - z') + \frac{z' - z}{zz'} = (z - z')\left(1 - \frac{1}{zz'}\right) = \frac{(z - z')(zz' - 1)}{zz'}$.
      Un produit est nul si et seulement si l'un des facteurs l'est : $g(z) = g(z') \iff z = z'$ ou $zz' = 1$.
\item Si $|z| > 1$ et $|z'| > 1$, alors $|zz'| = |z|\,|z'| > 1$, donc $zz' \neq 1$ ; d'apr\`es 1, $g(z) = g(z')$ impose $z = z'$.
\item $g\left(\expo{i\theta}\right) = \expo{i\theta} + \expo{-i\theta} = 2\cos\theta$. Quand $\theta$ d\'ecrit $\mathbb{R}$, $2\cos\theta$ d\'ecrit $[-2,2]$ :
      l'image du cercle unit\'e est le segment de l'axe r\'eel d'extr\'emit\'es les points d'affixes $-2$ et $2$.
\item $g(z) = 1 \iff z^2 - z + 1 = 0$ ($z \neq 0$) ; $\Delta = -3 = (i\sqrt{3})^2$, d'o\`u $z = \frac{1 \pm i\sqrt{3}}{2} = \expo{\pm i\frac{\pi}{3}}$.
      Ces nombres sont de module $1$ et $2\cos\left(\pm\frac{\pi}{3}\right) = 1$ : c'est coh\'erent avec 3.
\end{enumerate}

\paragraph{Exercice 8.}
\begin{enumerate}
\item $u = \sqrt{2}\left(\cos\frac{\pi}{4} + i\sin\frac{\pi}{4}\right)$ et $\bar{u} = \sqrt{2}\left(\cos\frac{\pi}{4} - i\sin\frac{\pi}{4}\right)$.
\item Par la formule de Moivre, $u^n + \bar{u}^n = (\sqrt{2})^n\left[2\cos\frac{n\pi}{4}\right]$, donc $\lambda_n = 2(\sqrt{2})^n = 2^{\frac{n}{2}+1}$.
\item $S_n = 0 \iff \cos\frac{n\pi}{4} = 0 \iff \frac{n\pi}{4} = \frac{\pi}{2} + k\pi \iff n = 2 + 4k$ :
      $S_n = 0$ pour les entiers $n$ de la forme $4k + 2$ ($k \in \mathbb{N}$).
\item Si $n = 2m$, $S_n = 2\cdot 2^m\cos\frac{m\pi}{2}$ et $\cos\frac{m\pi}{2} \in \{-1, 0, 1\}$ : $S_n \in \{0, 2^{m+1}, -2^{m+1}\} \subset \mathbb{Z}$.
\end{enumerate}

\paragraph{Exercice 9.}
\begin{enumerate}
\item $u = 2\expo{i\frac{2\pi}{3}}$ ; $\bar{u} = 2\expo{-i\frac{2\pi}{3}}$ ; $-2u = 4\expo{i\left(\frac{2\pi}{3}+\pi\right)} = 4\expo{-i\frac{\pi}{3}}$ ;
      $iu = 2\expo{i\left(\frac{2\pi}{3}+\frac{\pi}{2}\right)} = 2\expo{-i\frac{5\pi}{6}}$ ; $\frac{1}{u} = \frac{1}{2}\expo{-i\frac{2\pi}{3}}$ ;
      $\frac{-i}{u^2} = \frac{\expo{-i\frac{\pi}{2}}}{4\expo{i\frac{4\pi}{3}}} = \frac{1}{4}\expo{i\frac{\pi}{6}}$.
      Les images se construisent \`a partir de leur module (cercle de centre $O$) et de leur argument.
\item $z^2 - 2iz - i\sqrt{3} = (z - i)^2 + 1 - i\sqrt{3} = (z - i)^2 - u$. L'\'equation \'equivaut \`a $(z - i)^2 = u = \left(\sqrt{2}\,\expo{i\frac{\pi}{3}}\right)^2$,
      soit $z - i = \pm\left(\frac{\sqrt{2}}{2} + i\frac{\sqrt{6}}{2}\right)$.
      Solutions : $\boxed{\frac{\sqrt{2}}{2} + i\left(1 + \frac{\sqrt{6}}{2}\right)}$ et $\boxed{-\frac{\sqrt{2}}{2} + i\left(1 - \frac{\sqrt{6}}{2}\right)}$.
\end{enumerate}

\paragraph{Exercice 10.}
On calcule $\Delta$, une racine carr\'ee $\delta$ de $\Delta$, puis $z = \frac{-b \pm \delta}{2a}$.
\begin{enumerate}
\item $\Delta = -16 = (4i)^2$ : $S = \{1 - 2i ; 1 + 2i\}$.
\item $\Delta = 16 - 24 = -8 = (2\sqrt{2}i)^2$ : $S = \{2 - \sqrt{2}i ; 2 + \sqrt{2}i\}$.
\item $\Delta = (1+i)^2 - 20i = -18i = (3 - 3i)^2$ : $S = \left\{\frac{1 + i \pm (3 - 3i)}{2}\right\} = \{2 - i ; -1 + 2i\}$.
\item $\Delta = (1+3i)^2 - 4(-4+3i) = 8 - 6i = (3 - i)^2$ : $S = \left\{\frac{-1 - 3i \pm (3 - i)}{2}\right\} = \{1 - 2i ; -2 - i\}$.
\item $\Delta = (3 + 2\sqrt{2}i)^2 - 12\sqrt{2}i = 1 = 1^2$ : $S = \{1 + \sqrt{2}i ; 2 + \sqrt{2}i\}$.
\item $\Delta = (6 - 5i)^2 - 8(7 - 4i) = -45 - 28i = (2 - 7i)^2$ : $S = \left\{\frac{6 - 5i \pm (2 - 7i)}{4}\right\} = \left\{2 - 3i ; 1 + \frac{1}{2}i\right\}$.
\item $\Delta = (3 - 4i)^2 + 4(1 + 7i) = -3 + 4i = (1 + 2i)^2$ : $S = \left\{\frac{3 - 4i \pm (1 + 2i)}{2}\right\} = \{2 - i ; 1 - 3i\}$.
\end{enumerate}

\paragraph{Exercice 11.}
Notons $c = \cos\theta$ et $s = \sin\theta$. Par Moivre et la formule du bin\^ome,
$\cos 3\theta + i\sin 3\theta = (c + is)^3 = c^3 - 3cs^2 + i(3c^2s - s^3)$. En identifiant et avec $s^2 = 1 - c^2$ (resp. $c^2 = 1 - s^2$) :
\[
\cos 3\theta = 4\cos^3\theta - 3\cos\theta , \qquad \sin 3\theta = 3\sin\theta - 4\sin^3\theta .
\]
De m\^eme, $(c + is)^5 = c^5 - 10c^3s^2 + 5cs^4 + i(5c^4s - 10c^2s^3 + s^5)$, d'o\`u
\[
\cos 5\theta = 16\cos^5\theta - 20\cos^3\theta + 5\cos\theta , \qquad \sin 5\theta = 16\sin^5\theta - 20\sin^3\theta + 5\sin\theta .
\]

\paragraph{Exercice 12.}
\begin{enumerate}
\item $\frac{1}{z} = \cos\theta - i\sin\theta$, donc $z + \frac{1}{z} = 2\cos\theta$ et $z - \frac{1}{z} = 2i\sin\theta$ ;
      de m\^eme (Moivre) $z^n + \frac{1}{z^n} = 2\cos n\theta$ et $z^n - \frac{1}{z^n} = 2i\sin n\theta$.
\item $\left(z + \frac{1}{z}\right)^3 = \left(z^3 + \frac{1}{z^3}\right) + 3\left(z + \frac{1}{z}\right)$, soit $8\cos^3\theta = 2\cos 3\theta + 6\cos\theta$ :
      $\cos^3\theta = \frac{1}{4}(\cos 3\theta + 3\cos\theta)$.

      $\left(z - \frac{1}{z}\right)^3 = \left(z^3 - \frac{1}{z^3}\right) - 3\left(z - \frac{1}{z}\right)$, soit $-8i\sin^3\theta = 2i\sin 3\theta - 6i\sin\theta$ :
      $\sin^3\theta = \frac{1}{4}(3\sin\theta - \sin 3\theta)$.

      $\left(z + \frac{1}{z}\right)^4 = \left(z^4 + \frac{1}{z^4}\right) + 4\left(z^2 + \frac{1}{z^2}\right) + 6$, soit
      $\cos^4\theta = \frac{1}{8}(\cos 4\theta + 4\cos 2\theta + 3)$.

      $\left(z - \frac{1}{z}\right)^4 = \left(z^4 + \frac{1}{z^4}\right) - 4\left(z^2 + \frac{1}{z^2}\right) + 6$, et $(2i\sin\theta)^4 = 16\sin^4\theta$ :
      $\sin^4\theta = \frac{1}{8}(\cos 4\theta - 4\cos 2\theta + 3)$.
\end{enumerate}

\paragraph{Exercice 13.}
On \'ecrit $\cos x = \frac{\expo{ix} + \expo{-ix}}{2}$, $\sin x = \frac{\expo{ix} - \expo{-ix}}{2i}$, on d\'eveloppe par la formule du bin\^ome et on regroupe les termes conjugu\'es :
\begin{enumerate}
\item $\cos^2 x = \frac{1 + \cos 2x}{2}$ ; $\sin^3 x = \frac{3\sin x - \sin 3x}{4}$ ;
      $\cos^4 x\,\sin x = \frac{2\sin x + 3\sin 3x + \sin 5x}{16}$.
\item $\cos^4 x = \frac{\cos 4x + 4\cos 2x + 3}{8}$ ; $\sin^4 x = \frac{\cos 4x - 4\cos 2x + 3}{8}$ ;
      $\cos^5 x = \frac{\cos 5x + 5\cos 3x + 10\cos x}{16}$ ; $\sin^5 x = \frac{\sin 5x - 5\sin 3x + 10\sin x}{16}$ ;
      $\cos^6 x = \frac{\cos 6x + 6\cos 4x + 15\cos 2x + 10}{32}$ ; $\sin^6 x = \frac{-\cos 6x + 6\cos 4x - 15\cos 2x + 10}{32}$.
\end{enumerate}
Par exemple, $\cos^4 x\sin x = \frac{(\expo{ix} + \expo{-ix})^4(\expo{ix} - \expo{-ix})}{32i}$ et
$(\expo{ix} + \expo{-ix})^4(\expo{ix} - \expo{-ix}) = (\expo{5ix} - \expo{-5ix}) + 3(\expo{3ix} - \expo{-3ix}) + 2(\expo{ix} - \expo{-ix})$.

\paragraph{Exercice 14.}
Pour un syst\`eme o\`u figurent $\bar{x}$ et $\bar{y}$, on prend le conjugu\'e de l'\'equation concern\'ee
pour se ramener \`a un syst\`eme lin\'eaire en $x$ et $y$, que l'on r\'esout par combinaison ou par les d\'eterminants
($D = a d - b c$, $x = \frac{e d - b f}{D}$, $y = \frac{a f - e c}{D}$ pour le syst\`eme $ax + by = e$, $cx + dy = f$).
\begin{enumerate}
\item $D = -(1-i)(1+i) - (2-i)(2+i) = -7$, $x = \frac{-7 + 14i}{-7}$, $y = \frac{-14 - 21i}{-7}$ : $\boxed{x = 1 - 2i,\ y = 2 + 3i}$.
\item La seconde \'equation conjugu\'ee s'\'ecrit $(1+i)x - (5-2i)y = -10 - i$. Alors $D = -8 - 8i$ et $\boxed{x = 1 - i,\ y = 2 + i}$.
\item La seconde \'equation conjugu\'ee : $(3+i)x + (1-3i)y = 3 + 11i$. Alors $D = 6 - 8i$ et $\boxed{x = 1 + i,\ y = -2 + i}$.
\item La premi\`ere \'equation conjugu\'ee : $(2-i)x + (3+2i)y = 9 - 8i$ ; en \'ecrivant d'abord la seconde \'equation puis la premi\`ere conjugu\'ee, $D = (4-3i)(3+2i) + (2-3i)(2-i) = 19 - 9i$ et $\boxed{x = 3 - 2i,\ y = 1 - i}$.
\item Les inconnues sont $x$ et $Y = \bar{y}$ : $D = -16$, $x = \frac{32 + 4i}{-16} = -2 - \frac{1}{4}i$, $Y = \frac{48 + 36i}{-16} = -3 - \frac{9}{4}i$,
      donc $\boxed{x = -2 - \frac{1}{4}i,\ y = -3 + \frac{9}{4}i}$.
\end{enumerate}
On v\'erifie chaque r\'esultat en le reportant dans le syst\`eme.

\paragraph{Exercice 15.}
\begin{enumerate}
\item $\Delta = 16 - 52 = -36 = (6i)^2$ : $S = \{-2 - 3i ; -2 + 3i\}$.
\item $\Delta = (2i)^2 + 20 = 16$ : $S = \left\{\frac{2i \pm 4}{2}\right\} = \{2 + i ; -2 + i\}$.
\item $\Delta = (3 - 2i)^2 - 4(5 - 5i) = -15 + 8i$. On cherche $\delta = x + iy$ : $x^2 + y^2 = 17$, $x^2 - y^2 = -15$, $xy = 4 > 0$,
      d'o\`u $\delta = 1 + 4i$ et $S = \left\{\frac{3 - 2i \pm (1 + 4i)}{2}\right\} = \{2 + i ; 1 - 3i\}$.
\item $\Delta = (2 - 2i)^2 - 4(3 + 2i) = -12 - 16i$ ; $x^2 + y^2 = 20$, $x^2 - y^2 = -12$, $xy = -8 < 0$ : $\delta = 2 - 4i$,
      $S = \left\{\frac{2 - 2i \pm (2 - 4i)}{2}\right\} = \{2 - 3i ; i\}$.
\item $\Delta = 4 - 16(1+i) = -12 - 16i = (2 - 4i)^2$ :
      $z = \frac{2 \pm (2 - 4i)}{2(1+i)}$, soit $\frac{4 - 4i}{2 + 2i} = -2i$ et $\frac{4i}{2 + 2i} = 1 + i$ : $S = \{-2i ; 1 + i\}$.
\end{enumerate}
Dans chaque cas, on v\'erifie que la somme des solutions vaut $-\frac{b}{a}$ et leur produit $\frac{c}{a}$ ;
par exemple pour 5 : $-2i + (1+i) = 1 - i = \frac{2}{1+i}$ et $(-2i)(1+i) = 2 - 2i = \frac{4}{1+i}$.

\paragraph{Exercice 16.}
\begin{enumerate}
\item $8i = 8\expo{i\frac{\pi}{2}}$ ; ses racines cubiques sont $2\expo{i\left(\frac{\pi}{6} + \frac{2k\pi}{3}\right)}$, $k \in \{0,1,2\}$ :
      $2\expo{i\frac{\pi}{6}} = \sqrt{3} + i$, $2\expo{i\frac{5\pi}{6}} = -\sqrt{3} + i$, $2\expo{i\frac{3\pi}{2}} = -2i$.
\item $-64 = 64\expo{i\pi}$ ; ses racines sixi\`emes sont $2\expo{i\left(\frac{\pi}{6} + \frac{k\pi}{3}\right)}$, $0 \leq k \leq 5$ :
      $\sqrt{3} + i$, $2i$, $-\sqrt{3} + i$, $-\sqrt{3} - i$, $-2i$, $\sqrt{3} - i$.
      Leurs images sont les sommets d'un hexagone r\'egulier inscrit dans le cercle de centre $O$ et de rayon $2$.
\item (a) $z^3 = -27 = 27\expo{i\pi}$ : $z = 3\expo{i\left(\frac{\pi}{3} + \frac{2k\pi}{3}\right)}$, soit
      $S = \left\{\frac{3}{2} + \frac{3\sqrt{3}}{2}i ; -3 ; \frac{3}{2} - \frac{3\sqrt{3}}{2}i\right\}$.

      (b) $-8 + 8\sqrt{3}i = 16\expo{i\frac{2\pi}{3}}$ : $z = 2\expo{i\left(\frac{\pi}{6} + \frac{k\pi}{2}\right)}$, $0 \leq k \leq 3$, soit
      $S = \{\sqrt{3} + i ; -1 + \sqrt{3}i ; -\sqrt{3} - i ; 1 - \sqrt{3}i\}$.
\item Une solution $z = 2\expo{i\left(\frac{\pi}{6} + \frac{k\pi}{2}\right)}$ est racine sixi\`eme de $-64$ si et seulement si $\frac{k\pi}{2}$ est un multiple de $\frac{\pi}{3}$ \`a $2\pi$ pr\`es,
      soit $k = 0$ ou $k = 2$ : ce sont $\sqrt{3} + i$ et $-\sqrt{3} - i$ (on v\'erifie : $z^6 = 64\expo{i\pi} = -64$).
\end{enumerate}

\paragraph{Exercice 17.}
\begin{enumerate}
\item $\Delta = (5 + 4i\sqrt{3})^2 - 36 = 25 + 40\sqrt{3}i - 48 - 36 = -59 + 40\sqrt{3}i$.
      On cherche $\delta = x + iy$ : $x^2 + y^2 = |\Delta| = \sqrt{3481 + 4800} = 91$, $x^2 - y^2 = -59$, $xy = 20\sqrt{3} > 0$ ;
      d'o\`u $x^2 = 16$, $y^2 = 75$ et $\delta = 4 + 5\sqrt{3}i$.
      Les solutions sont $\frac{5 + 4\sqrt{3}i \pm (4 + 5\sqrt{3}i)}{2}$ : $\boxed{\frac{9}{2} + \frac{9\sqrt{3}}{2}i}$ et $\boxed{\frac{1}{2} - \frac{\sqrt{3}}{2}i}$.
\item $z' = \frac{9}{2} + \frac{9\sqrt{3}}{2}i = 9\expo{i\frac{\pi}{3}}$ (module $9$, argument $\frac{\pi}{3}$) et
      $z'' = \expo{-i\frac{\pi}{3}}$ (module $1$, argument $-\frac{\pi}{3}$).
      $\frac{z'}{z''} = 9\expo{i\frac{2\pi}{3}}$ (module $9$, argument $\frac{2\pi}{3}$) et $z'z'' = 9$ (module $9$, argument $0$ ; c'est le produit des racines $\frac{c}{a}$).
\item $Z^2 = 9\expo{i\frac{\pi}{3}}$ donne $Z = 3\expo{i\frac{\pi}{6}}$ ou $Z = 3\expo{i\frac{7\pi}{6}}$, soit
      $Z = \pm\left(\frac{3\sqrt{3}}{2} + \frac{3}{2}i\right) = \pm 3\left(\cos\frac{\pi}{6} + i\sin\frac{\pi}{6}\right)$ ;
      en forme polaire : $\left[3, \frac{\pi}{6}\right]$ et $\left[3, \frac{7\pi}{6}\right]$.
\end{enumerate}

\paragraph{Exercice 18.}
\begin{enumerate}
\item $(1-2i)^2 = -3 - 4i$, $(1-2i)^3 = (-3-4i)(1-2i) = \boxed{-11 + 2i}$ ;
      $(2-i)^2 = 3 - 4i$, $(2-i)^4 = (3-4i)^2 = \boxed{-7 - 24i}$.
\item (a) $z^3 = (1-2i)^3 \iff \left(\frac{z}{1-2i}\right)^3 = 1 \iff z = (1-2i)u$ avec $u \in \{1, j, \bar{j}\}$, $j = -\frac{1}{2} + \frac{\sqrt{3}}{2}i$.
      Solutions : $1 - 2i$, $(1-2i)j = \left(\sqrt{3} - \frac{1}{2}\right) + \left(1 + \frac{\sqrt{3}}{2}\right)i$ et
      $(1-2i)\bar{j} = \left(-\sqrt{3} - \frac{1}{2}\right) + \left(1 - \frac{\sqrt{3}}{2}\right)i$.

      (b) $z^4 = -7 - 24i = (2-i)^4 \iff \left(\frac{z}{2-i}\right)^4 = 1 \iff z = (2-i)u$, $u \in \{1, i, -1, -i\}$ :
      $S = \{2 - i ; 1 + 2i ; -2 + i ; -1 - 2i\}$.
\end{enumerate}

\paragraph{Exercice 19.}
\begin{enumerate}
\item $z^3 + z^2 + z + 1 = z^2(z+1) + (z+1) = (z+1)(z^2+1)$ : $S = \{-1 ; i ; -i\}$.
\item On pose $w = \frac{z-i}{z+i}$ ($z \neq -i$) : $w$ est solution de l'\'equation 1, donc $w \in \{-1, i, -i\}$.
      De $w(z + i) = z - i$ on tire $z = \frac{i(1+w)}{1-w}$ ($w \neq 1$).
      $w = -1$ donne $z = 0$ ; $w = i$ donne $z = \frac{i(1+i)}{1-i} = i \cdot i = -1$ ; $w = -i$ donne $z = \frac{i(1-i)}{1+i} = 1$.
      $S = \{0 ; -1 ; 1\}$.
\item $1 - z + z^2 - z^3 = (1 - z) + z^2(1 - z) = (1-z)(1+z^2)$ : $S = \{1 ; i ; -i\}$.
\item $2 - 2\sqrt{3}i = 4\expo{-i\frac{\pi}{3}}$ : $z = \sqrt{2}\,\expo{i\left(-\frac{\pi}{12} + \frac{k\pi}{2}\right)}$, $k \in \{0,1,2,3\}$,
      c'est-\`a-dire $\sqrt{2}\expo{-i\frac{\pi}{12}}$, $\sqrt{2}\expo{i\frac{5\pi}{12}}$, $\sqrt{2}\expo{i\frac{11\pi}{12}}$, $\sqrt{2}\expo{i\frac{17\pi}{12}}$.
      Avec $\cos\frac{\pi}{12} = \frac{\sqrt{6}+\sqrt{2}}{4}$ et $\sin\frac{\pi}{12} = \frac{\sqrt{6}-\sqrt{2}}{4}$ :
      $\sqrt{2}\expo{-i\frac{\pi}{12}} = \frac{\sqrt{3}+1}{2} - \frac{\sqrt{3}-1}{2}i$, et les autres s'en d\'eduisent en multipliant par $i$, $-1$, $-i$.
\end{enumerate}

\paragraph{Exercice 20.}
Pour $\alpha$ r\'eel, $P(\alpha) = 0$ \'equivaut \`a : partie r\'eelle et partie imaginaire de $P(\alpha)$ nulles.
\begin{enumerate}
\item $P(\alpha) = (2\alpha^3 - 5\alpha^2 + 1) + i(-6\alpha^2 + 9\alpha - 3)$ et $-6\alpha^2 + 9\alpha - 3 = -3(2\alpha - 1)(\alpha - 1)$.
      Seul $\alpha = \frac{1}{2}$ annule aussi la partie r\'eelle ($2\cdot\frac{1}{8} - \frac{5}{4} + 1 = 0$, alors que $2 - 5 + 1 \neq 0$) : $\boxed{\alpha = \frac{1}{2}}$.
\item $P(z) = \left(z - \frac{1}{2}\right)\left(2z^2 - (4+6i)z - 2 + 6i\right)$ (identification ou division).
      Pour $z^2 - (2+3i)z - 1 + 3i = 0$ : $\Delta = (2+3i)^2 + 4 - 12i = -1 = i^2$, d'o\`u $z = \frac{2 + 3i \pm i}{2}$.
      $S = \left\{\frac{1}{2} ; 1 + i ; 1 + 2i\right\}$.
\end{enumerate}
Pour $Q$ : la partie imaginaire de $Q(\alpha)$ est $-(\alpha - 3)(5\alpha - 4)$ et la partie r\'eelle $(\alpha - 3)(\alpha^2 - \alpha - 4)$ ; seule valeur commune : $\alpha = 3$.
$Q(z) = (z - 3)\left(z^2 - (1+5i)z - 4 + 4i\right)$, $\Delta = (1+5i)^2 + 16 - 16i = -8 - 6i = (1 - 3i)^2$, d'o\`u les racines $\frac{1 + 5i \pm (1 - 3i)}{2}$ :
$S_Q = \{3 ; 1 + i ; 4i\}$.

Le dernier polyn\^ome est $\overline{Q(\bar{z})}$ (on a conjugu\'e tous les coefficients) : $z$ en est racine si et seulement si $\bar{z}$ est racine de $Q$.
Ses racines sont donc les conjugu\'ees : $\{3 ; 1 - i ; -4i\}$.

\paragraph{Exercice 21.}
\begin{enumerate}
\item $P(i\alpha) = (\alpha - 2)(3\alpha + 1) - i(\alpha - 2)(\alpha^2 - 3)$ (calcul direct en rempla\c{c}ant $z$ par $i\alpha$) : la seule valeur qui annule
      les deux parties est $\boxed{\alpha = 2}$.
\item $P(z) = (z - 2i)\left(z^2 - 3z + 3 - i\right)$ ; $\Delta = 9 - 12 + 4i = -3 + 4i = (1 + 2i)^2$ ; racines $\frac{3 \pm (1 + 2i)}{2}$.
      $S = \{2i ; 2 + i ; 1 - i\}$.
\end{enumerate}
Pour $Q$ : $Q(i\alpha) = -12(2\alpha - 3) - i(2\alpha - 3)(\alpha^2 - 5)$, donc $\alpha = \frac{3}{2}$ ;
$Q(z) = \left(z - \frac{3}{2}i\right)\left(2z^2 + 10 + 24i\right)$, donc $z^2 = -5 - 12i = (2 - 3i)^2$ et $S_Q = \left\{\frac{3}{2}i ; 2 - 3i ; -2 + 3i\right\}$.

Pour $R$ : $R(i\alpha) = (\alpha - 5)(3\alpha - 2) - i\,\alpha(\alpha - 1)(3\alpha - 2)$, donc $\alpha = \frac{2}{3}$ ;
$R(z) = \left(z - \frac{2}{3}i\right)\left(3z^2 - (3+3i)z + 15i\right)$, soit $z^2 - (1+i)z + 5i = 0$, $\Delta = 2i - 20i = -18i = (3 - 3i)^2$ :
$S_R = \left\{\frac{2}{3}i ; 2 - i ; -1 + 2i\right\}$.

\paragraph{Exercice 22.}
\begin{enumerate}
\item On remplace $z$ par $\alpha(1 - i)$ : la partie imaginaire de $P(\alpha - i\alpha)$ vaut $-2(\alpha - 2)(\alpha - 1)$ et sa partie r\'eelle
      $-2(\alpha - 2)(2\alpha^2 + \alpha + 3)$, o\`u $2\alpha^2 + \alpha + 3$ n'a pas de racine r\'eelle : $\boxed{\alpha = 2}$, racine $2 - 2i$.
\item $P(z) = (z - 2 + 2i)\left((1-i)z^2 + (1-i)z - 4 - 2i\right)$.
      En divisant par $1 - i$ : $z^2 + z - (1 + 3i) = 0$, $\Delta = 1 + 4 + 12i = 5 + 12i = (3 + 2i)^2$, racines $\frac{-1 \pm (3 + 2i)}{2}$.
      $S = \{2 - 2i ; 1 + i ; -2 - i\}$.
\end{enumerate}
Pour $Q$ : la partie imaginaire de $Q(\alpha - i\alpha)$ vaut $6(\alpha - 2)(\alpha - 1)$ et la partie r\'eelle $-2(\alpha - 2)(2\alpha^2 - \alpha - 11)$, qui ne s'annule pas en $1$ : $\alpha = 2$.
$Q(z) = (z - 2 + 2i)\left((1-i)z^2 + (-3+i)z + 14 + 8i\right)$ ; en divisant par $1 - i$ : $z^2 - (2 + i)z + 3 + 11i = 0$,
$\Delta = (2+i)^2 - 12 - 44i = -9 - 40i = (4 - 5i)^2$, racines $\frac{2 + i \pm (4 - 5i)}{2}$ :
$S_Q = \{2 - 2i ; 3 - 2i ; -1 + 3i\}$.

\paragraph{Exercice 23.}
\begin{enumerate}
\item Pour $x$ r\'eel, $P(x) = (x^4 - 3x^2 - 4) + i(-x^3 + 3x^2 + 4x - 12)$, avec
      $x^4 - 3x^2 - 4 = (x^2 - 4)(x^2 + 1)$ et $-x^3 + 3x^2 + 4x - 12 = -(x - 3)(x^2 - 4)$.
      Les deux parties s'annulent pour $x = 2$ et $x = -2$ : ce sont les racines r\'eelles de $P$.
\item $P(z) = (z^2 - 4)\left(z^2 - iz + 1 + 3i\right)$ ; $\Delta = -1 - 4 - 12i = -5 - 12i = (2 - 3i)^2$, racines $\frac{i \pm (2 - 3i)}{2}$.
      $S = \{2 ; -2 ; 1 - i ; -1 + 2i\}$.
\end{enumerate}

\paragraph{Exercice 24.}
On utilise $a^2 + b^2 = (a + ib)(a - ib)$.
\begin{enumerate}
\item L'\'equation \'equivaut \`a $2z^2 - 3z + 2 = i(z^2 - 3z + 2)$ ou $2z^2 - 3z + 2 = -i(z^2 - 3z + 2)$, soit
      $(2 - i)z^2 - 3(1 - i)z + 2 - 2i = 0$ ou $(2 + i)z^2 - 3(1 + i)z + 2 + 2i = 0$.
      Pour la premi\`ere : $\Delta = 9(1-i)^2 - 4(2-i)(2-2i) = -18i - 8 + 24i = -8 + 6i = (1 + 3i)^2$, d'o\`u
      $z = \frac{3 - 3i \pm (1 + 3i)}{2(2 - i)}$, soit $\frac{4}{2(2-i)} = \frac{4}{5} + \frac{2}{5}i$ et $\frac{2 - 6i}{2(2-i)} = 1 - i$.
      La seconde \'equation a les coefficients conjugu\'es, donc les solutions conjugu\'ees.
      $S = \left\{1 - i ; 1 + i ; \frac{4}{5} + \frac{2}{5}i ; \frac{4}{5} - \frac{2}{5}i\right\}$.
\item L'\'equation \'equivaut \`a $z^2 + 4z + 1 + i(3z + 5) = 0$ ou $z^2 + 4z + 1 - i(3z + 5) = 0$.
      Les solutions de la premi\`ere sont celles de 3 ; la seconde, \`a coefficients conjugu\'es, a les solutions conjugu\'ees.
      $S = \{-1 - i ; -3 - 2i ; -1 + i ; -3 + 2i\}$.
\item $z^2 + (4 + 3i)z + 1 + 5i = 0$ : $\Delta = (4 + 3i)^2 - 4 - 20i = 3 + 4i = (2 + i)^2$, d'o\`u $z = \frac{-4 - 3i \pm (2 + i)}{2}$ :
      $S = \{-1 - i ; -3 - 2i\}$.
\end{enumerate}

\paragraph{Exercice 25.}
\begin{enumerate}
\item Avec $Z = z^2$ : $Z^2 - Z + 1 = 0$, $Z = \frac{1 \pm i\sqrt{3}}{2} = \expo{\pm i\frac{\pi}{3}}$.
      Les racines carr\'ees sont $\pm\expo{i\frac{\pi}{6}}$ et $\pm\expo{-i\frac{\pi}{6}}$ :
      $S = \left\{\frac{\sqrt{3}}{2} + \frac{1}{2}i ; -\frac{\sqrt{3}}{2} - \frac{1}{2}i ; \frac{\sqrt{3}}{2} - \frac{1}{2}i ; -\frac{\sqrt{3}}{2} + \frac{1}{2}i\right\}$.
\item En groupant les racines conjugu\'ees : $\left(z - \expo{i\frac{\pi}{6}}\right)\left(z - \expo{-i\frac{\pi}{6}}\right) = z^2 - \sqrt{3}z + 1$ et
      $\left(z + \expo{i\frac{\pi}{6}}\right)\left(z + \expo{-i\frac{\pi}{6}}\right) = z^2 + \sqrt{3}z + 1$.
      Donc $P(z) = (z^2 - \sqrt{3}z + 1)(z^2 + \sqrt{3}z + 1)$, ce que l'on v\'erifie aussi par $(z^2 + 1)^2 - 3z^2 = z^4 - z^2 + 1$.
\end{enumerate}

\paragraph{Exercice 26.}
Avec $Z = z^3$ : $Z^2 - (1 - i)Z - i = 0$, $\Delta = (1-i)^2 + 4i = 2i = (1 + i)^2$, donc $Z = \frac{1 - i \pm (1 + i)}{2}$, soit $Z = 1$ ou $Z = -i$.
\begin{itemize}
\item $z^3 = 1$ : $z \in \left\{1 ; -\frac{1}{2} + \frac{\sqrt{3}}{2}i ; -\frac{1}{2} - \frac{\sqrt{3}}{2}i\right\}$.
\item $z^3 = -i = \expo{-i\frac{\pi}{2}}$ : $z = \expo{i\left(-\frac{\pi}{6} + \frac{2k\pi}{3}\right)}$, soit $z \in \left\{\frac{\sqrt{3}}{2} - \frac{1}{2}i ; i ; -\frac{\sqrt{3}}{2} - \frac{1}{2}i\right\}$.
\end{itemize}
Les six images sont sur le cercle unit\'e ; elles forment deux triangles \'equilat\'eraux inscrits dans ce cercle.

\paragraph{Exercice 27.}
\begin{enumerate}
\item $x^5 = -1 = \expo{i\pi}$ : les solutions sont $\expo{i\frac{(2k+1)\pi}{5}}$, $0 \leq k \leq 4$ : toutes de module $1$, d'arguments
      $\frac{\pi}{5}$, $\frac{3\pi}{5}$, $\pi$, $\frac{7\pi}{5}$, $\frac{9\pi}{5}$.
\item $\frac{7\pi}{5} \equiv -\frac{3\pi}{5}$ et $\frac{9\pi}{5} \equiv -\frac{\pi}{5}$ : les racines autres que $-1$ sont deux \`a deux conjugu\'ees et
      $\left(x - \expo{i\varphi}\right)\left(x - \expo{-i\varphi}\right) = x^2 - 2x\cos\varphi + 1$. Comme $x^5 + 1$ est unitaire de degr\'e $5$ et a ces cinq racines,
      $x^5 + 1 = (x + 1)\left(x^2 - 2x\cos\frac{\pi}{5} + 1\right)\left(x^2 - 2x\cos\frac{3\pi}{5} + 1\right)$.
\item On a aussi $x^5 + 1 = (x+1)(x^4 - x^3 + x^2 - x + 1)$. En d\'eveloppant le produit des deux trin\^omes et en identifiant, avec
      $c_1 = \cos\frac{\pi}{5}$, $c_3 = \cos\frac{3\pi}{5}$ : $-2(c_1 + c_3) = -1$ et $2 + 4c_1c_3 = 1$, soit $c_1 + c_3 = \frac{1}{2}$ et $c_1c_3 = -\frac{1}{4}$.
      $c_1$ et $c_3$ sont les racines de $X^2 - \frac{1}{2}X - \frac{1}{4} = 0$, soit $\frac{1 \pm \sqrt{5}}{4}$. Comme $c_1 > 0$ :
      $\cos\frac{\pi}{5} = \boxed{\frac{1 + \sqrt{5}}{4}}$ et $\cos\frac{3\pi}{5} = \boxed{\frac{1 - \sqrt{5}}{4}}$.
\end{enumerate}

\paragraph{Exercice 28.}
\begin{enumerate}
\item $z_1 = \frac{1 + i}{2}$, $|z_1| = \frac{\sqrt{2}}{2} < 1$. Si $|z_n| < 1$, alors $|z_{n+1}| = \frac{|z_n + i|}{2} \leq \frac{|z_n| + 1}{2} < 1$.
      Par r\'ecurrence, $|z_n| < 1$ pour tout $n \geq 1$.
\item (a) $u_{n+1} = z_{n+1} - i = \frac{z_n + i}{2} - i = \frac{z_n - i}{2} = \frac{1}{2}u_n$.
      (b) $(u_n)$ est g\'eom\'etrique de raison $\frac{1}{2}$ et $u_0 = 1 - i$, donc $u_n = \frac{1 - i}{2^n}$ et $z_n = \frac{1}{2^n} + i\left(1 - \frac{1}{2^n}\right)$.
      Ainsi $x_n = \frac{1}{2^n}$ : suite g\'eom\'etrique de raison $\frac{1}{2}$, qui tend vers $0$ ; $y_n = 1 - \frac{1}{2^n} \to 1$ ;
      $|z_n| = \sqrt{x_n^2 + y_n^2} \to \sqrt{0 + 1} = 1$.
\item $|z_n - i| = |u_n| = \frac{\sqrt{2}}{2^n}$, suite d\'ecroissante. $\frac{\sqrt{2}}{2^n} < 10^{-5} \iff 2^n > \sqrt{2}\cdot 10^5 \approx 141\,421$.
      Comme $2^{17} = 131\,072$ et $2^{18} = 262\,144$, on obtient $\boxed{n_0 = 18}$.
\end{enumerate}

\paragraph{Exercice 29.}
\begin{enumerate}
\item La suite $d_n = z_n - z_{n-1}$ ($n \geq 1$) v\'erifie $d_n = i\,d_{n-1}$ et $d_1 = i$ : elle est g\'eom\'etrique de raison $i$, donc $z_n - z_{n-1} = i^n$.
\item $z_n = z_0 + \sum_{k=1}^{n} i^k = i\,\frac{1 - i^n}{1 - i} = \frac{i(1+i)}{2}(1 - i^n) = \frac{-1 + i}{2}(1 - i^n) = \frac{1 - i}{2}(i^n - 1)$.
\item $i^{n+4} = i^n$, donc $z_{n+4} = z_n$ : la suite est p\'eriodique de p\'eriode $4$.
\item $2017 = 4\times 504 + 1$, etc. : $z_{2017} = z_1 = i$, $z_{2018} = z_2 = -1 + i$, $z_{2019} = z_3 = -1$, $z_{2020} = z_0 = 0$.
      Les points $M_{2017}$, $M_{2018}$, $M_{2019}$, $M_{2020}$ sont les sommets d'un carr\'e de c\^ot\'e $1$.
\end{enumerate}

\paragraph{Exercice 30.}
Soit $q = \frac{\sqrt{2} + i\sqrt{6}}{4} = \frac{\sqrt{2}}{2}\left(\frac{1}{2} + \frac{\sqrt{3}}{2}i\right) = \frac{\sqrt{2}}{2}\expo{i\frac{\pi}{3}}$.
\begin{enumerate}
\item $z_n = 8q^n = 8\left(\frac{\sqrt{2}}{2}\right)^n\expo{i\frac{n\pi}{3}}$ :
      $z_0 = 8$ ; $z_1 = 4\sqrt{2}\expo{i\frac{\pi}{3}} = 2\sqrt{2} + 2\sqrt{6}i$ ; $z_2 = 4\expo{i\frac{2\pi}{3}} = -2 + 2\sqrt{3}i$ ;
      $z_3 = 2\sqrt{2}\expo{i\pi} = -2\sqrt{2}$ ; $z_4 = 2\expo{i\frac{4\pi}{3}} = -1 - \sqrt{3}i$ ; $z_5 = \sqrt{2}\expo{i\frac{5\pi}{3}} = \frac{\sqrt{2}}{2} - \frac{\sqrt{6}}{2}i$.
\item $z_n = 8\left(\frac{\sqrt{2}}{2}\right)^n\expo{i\frac{n\pi}{3}}$ ; $|z_n| = 8\left(\frac{\sqrt{2}}{2}\right)^n$ : suite g\'eom\'etrique de raison $\frac{\sqrt{2}}{2}$.
      Les arguments $\frac{n\pi}{3}$ forment une suite arithm\'etique de raison $\frac{\pi}{3}$.
\item $|z_n| = 8\cdot 2^{-\frac{n}{2}} < 10^{-2} \iff 2^{\frac{n}{2}} > 800$. Or $2^{9{,}5} \approx 724$ et $2^{10} = 1024$ : le plus petit entier est $\boxed{n = 20}$.
      $z_{20} = \frac{8}{2^{10}}\expo{i\frac{20\pi}{3}} = \frac{1}{128}\expo{i\frac{2\pi}{3}} = -\frac{1}{256} + \frac{\sqrt{3}}{256}i$.
      Comme $|z_n| \to 0$, $z_n \to 0$ : les points $M_n$ se rapprochent de $O$ en tournant autour de lui (spirale).
\item $OM_{n+1} = |z_{n+1}| = \frac{\sqrt{2}}{2}\,OM_n$ et $\left(\overrightarrow{OM_n},\overrightarrow{OM_{n+1}}\right) = \arg\frac{z_{n+1}}{z_n} = \arg q = \frac{\pi}{3}$.
      $M_{n+1}$ est l'image de $M_n$ par la similitude de centre $O$, de rapport $\frac{\sqrt{2}}{2}$ et d'angle $\frac{\pi}{3}$ :
      on trace la demi-droite issue de $O$ faisant l'angle $\frac{\pi}{3}$ avec $\overrightarrow{OM_n}$ et on y place le point \`a la distance $\frac{\sqrt{2}}{2}OM_n$ de $O$.
\end{enumerate}

\paragraph{Exercice 31.}
\begin{enumerate}
\item Figure : $A(0,-1)$, $B(3,0)$, $C(2,3)$, $D(-1,2)$.
\item (a) $\left|\frac{z_C - z_A}{z_D - z_B}\right| = \frac{AC}{BD}$ et un argument de ce quotient mesure $\left(\overrightarrow{BD},\overrightarrow{AC}\right)$.
      (b) $\frac{2 + 4i}{-4 + 2i} = \frac{(2+4i)(-4-2i)}{20} = \frac{-20i}{20} = -i$.
      (c) Module $1$ : $AC = BD$ ; argument $-\frac{\pi}{2}$ : $(AC) \perp (BD)$. Les diagonales ont m\^eme longueur et sont perpendiculaires.
\item (a) Le milieu de $[AC]$ a pour affixe $\frac{-i + 2 + 3i}{2} = 1 + i$ et celui de $[BD]$ a pour affixe $\frac{3 - 1 + 2i}{2} = 1 + i$ :
      $ABCD$ est un parall\'elogramme dont les diagonales sont de m\^eme longueur et perpendiculaires : c'est un \textbf{carr\'e}.
      (b) $AB = |3 + i| = \sqrt{10}$, donc $s_0 = 10$ unit\'es d'aire, soit $10 \times 4 = 40$~cm$^2$.
\end{enumerate}

\paragraph{Exercice 32.}
\begin{enumerate}
\item Pour $x$ r\'eel : $\mathfrak{Re}\,P(x) = (x - 6)(x - 3)(x + 3)$ et $\mathfrak{Im}\,P(x) = -9(x - 3)(x - 2)$ ; seule valeur commune : $z_0 = 3$.
      $P(z) = (z - 3)\left(z^2 - (3 + 9i)z - 18 + 18i\right)$.
\item $\Delta = (3 + 9i)^2 + 72 - 72i = -18i = (3 - 3i)^2$ ; racines $\frac{3 + 9i \pm (3 - 3i)}{2}$, soit $3 + 3i$ et $6i$.
      $z_0 = 3$ (module $3$, argument $0$) ; $z_1 = 3 + 3i = 3\sqrt{2}\expo{i\frac{\pi}{4}}$ ; $z_2 = 6i = 6\expo{i\frac{\pi}{2}}$.
\item $\frac{z_1}{z_0} = 1 + i$ et $\frac{z_2}{z_1} = \frac{6i}{3 + 3i} = \frac{2i}{1 + i} = 1 + i$ : ce sont les premiers termes de la suite g\'eom\'etrique de raison $1 + i$,
      $z_n = 3(1 + i)^n = 3(\sqrt{2})^n\expo{i\frac{n\pi}{4}}$.
\item $|z_n| = 3(\sqrt{2})^n$ : suite g\'eom\'etrique de raison $\sqrt{2}$ ; les arguments $\frac{n\pi}{4}$ forment une suite arithm\'etique de raison $\frac{\pi}{4}$.
\item $z_n$ r\'eel $\iff \frac{n\pi}{4} \equiv 0\ [\pi] \iff n \equiv 0\ [4]$ ; r\'eel n\'egatif $\iff \frac{n\pi}{4} \equiv \pi\ [2\pi] \iff n \equiv 4\ [8]$ ;
      imaginaire pur $\iff \frac{n\pi}{4} \equiv \frac{\pi}{2}\ [\pi] \iff n \equiv 2\ [4]$.
\end{enumerate}

\paragraph{Exercice 33.}
\begin{enumerate}
\item Pour $x$ r\'eel : $\mathfrak{Re}\,P(x) = (x + 2)(x^2 - 2x - 18)$, $\mathfrak{Im}\,P(x) = 3(x + 2)(3x - 2)$ : $z_1 = -2$.
\item $Q(z) = z^2 + (-2 + 9i)z - 18 - 6i$.
\item $Q(iy) = (-y^2 - 9y - 18) + i(-2y - 6)$ ; la partie imaginaire s'annule pour $y = -3$, et alors $-9 + 27 - 18 = 0$ : $z_2 = -3i$.
\item La somme des racines de $Q$ vaut $2 - 9i$, donc $z_3 = 2 - 9i + 3i = 2 - 6i$. $S = \{-2 ; -3i ; 2 - 6i\}$.
\item $\frac{z_3 - z_1}{z_2 - z_1} = \frac{4 - 6i}{2 - 3i} = 2$ est r\'eel : $\overrightarrow{AC} = 2\overrightarrow{AB}$, les points $A$, $B$, $C$ sont align\'es ($B$ milieu de $[AC]$).
\end{enumerate}

\paragraph{Exercice 34.}
\begin{enumerate}
\item Pour $x$ r\'eel : $\mathfrak{Re}\,P(x) = (2x + 1)(2x^2 - 3x + 4)$ et $\mathfrak{Im}\,P(x) = \frac{1}{2}(2x + 1)^2$ : $\alpha = -\frac{1}{2}$.
\item $P(z) = \left(z + \frac{1}{2}\right)\left(4z^2 - (6 - 2i)z + 8 + i\right)$ ; $\Delta = (6 - 2i)^2 - 16(8 + i) = -96 - 40i = (2 - 10i)^2$,
      racines $\frac{6 - 2i \pm (2 - 10i)}{8}$ : $z_1 = 1 - \frac{3}{2}i$ et $z_2 = \frac{1}{2} + i$.
      $S = \left\{-\frac{1}{2} ; 1 - \frac{3}{2}i ; \frac{1}{2} + i\right\}$.
\item Par soustraction : $\alpha - z_1 = a(z_2 - \alpha)$, soit $-\frac{3}{2} + \frac{3}{2}i = a(1 + i)$, donc $a = \frac{3}{2}\cdot\frac{-1 + i}{1 + i} = \frac{3}{2}i$ ;
      puis $b = z_1 - a\alpha = 1 - \frac{3}{2}i + \frac{3}{4}i = 1 - \frac{3}{4}i$.
\item $\omega = \frac{b}{1 - a} = \frac{4 - 3i}{4 - 6i} = \frac{(4 - 3i)(4 + 6i)}{52} = \frac{34 + 12i}{52} = \boxed{\frac{17}{26} + \frac{3}{13}i}$.
      ($\Omega$ est le centre de la similitude $z \mapsto az + b$, de rapport $\frac{3}{2}$ et d'angle $\frac{\pi}{2}$, qui envoie $C$ sur $A$ et $A$ sur $B$.)
\item Figure : $A\left(-\frac{1}{2}, 0\right)$, $B\left(1, -\frac{3}{2}\right)$, $C\left(\frac{1}{2}, 1\right)$, $\Omega \approx (0{,}65 ; 0{,}23)$.
\end{enumerate}

\paragraph{Exercice 35.}
\begin{enumerate}
\item $a = 0$, $b = 1$, $d = i$, $c = 1 + i$, $e = \frac{i}{2}$.
      $\overrightarrow{EF}$ est l'image de $\overrightarrow{ED}$ (affixe $\frac{i}{2}$) par la rotation d'angle $\frac{\pi}{2}$ : affixe $i\cdot\frac{i}{2} = -\frac{1}{2}$.
      Donc $f = -\frac{1}{2} + \frac{1}{2}i$ et $g = d + (f - e) = -\frac{1}{2} + i$.
\item On cherche $s : z \mapsto \alpha z + \beta$ avec $\alpha d + \beta = f$ et $\alpha b + \beta = d$.
      (a) $\alpha = \frac{f - d}{d - b} = \frac{-\frac{1}{2} - \frac{1}{2}i}{-1 + i} = \frac{i}{2}$ : rapport $k = |\alpha| = \frac{1}{2}$, angle $\theta = \frac{\pi}{2}$.
      (b) $\beta = d - \alpha = \frac{i}{2}$ : $s : z' = \frac{i}{2}z + \frac{i}{2}$.
      (c) $\omega = \frac{\beta}{1 - \alpha} = \frac{i/2}{1 - i/2} = \frac{i}{2 - i} = \boxed{-\frac{1}{5} + \frac{2}{5}i}$.
\end{enumerate}

\paragraph{Exercice 36.}
\begin{enumerate}
\setcounter{enumi}{1}
\item $\frac{z_C - z_A}{z_B - z_A} = \frac{4 + 4i}{2 - 2i} = 2i$ : $AC = 2AB$ et $\left(\overrightarrow{AB},\overrightarrow{AC}\right) = \frac{\pi}{2}$ :
      le triangle $ABC$ est rectangle en $A$ (et $AC = 2AB$).
\item (a) $f : z \mapsto \alpha z + \beta$ avec $2i\alpha + \beta = -1 + i$ et $2\alpha + \beta = 2i$ ; par diff\'erence $\alpha(2i - 2) = -1 - i$,
      d'o\`u $\alpha = \frac{i}{2}$ et $\beta = i$ : $f : z' = \frac{i}{2}z + i$.
      (b) Angle $\frac{\pi}{2}$, rapport $\frac{1}{2}$, centre $\omega = \frac{i}{1 - i/2} = \frac{2i}{2 - i} = -\frac{2}{5} + \frac{4}{5}i$.
      (c) $f(A) = D$, $f(B) = A$ et $f(C) = \frac{i}{2}(4 + 6i) + i = -3 + 3i = z_E$ : l'image du triangle $ABC$ est $DAE$.
      (d) Une similitude conserve les angles et multiplie les longueurs par $\frac{1}{2}$ : $DAE$ est rectangle en $D$, avec $DE = 2DA$.
\item (a) $f$ transforme la droite $(BC)$ en la droite $(AE)$ et le cercle $(\mathscr{C}_1)$ de diam\`etre $[AB]$ en le cercle de diam\`etre $[f(A)f(B)] = [DA]$, c'est-\`a-dire $(\mathscr{C}_2)$.
      Les points communs \`a $(\mathscr{C}_1)$ et $(BC)$ sont $B$ et $M$ ; leurs images sont les points communs \`a $(\mathscr{C}_2)$ et $(AE)$, soit $f(B) = A$ et $f(M)$.
      Donc $f(M) = N$. (Calcul : $M$ est le projet\'e orthogonal de $A$ sur $(BC)$, $z_M = \frac{12}{5} + \frac{6}{5}i$, et $z_N = f(z_M) = -\frac{3}{5} + \frac{11}{5}i$.)
      (b) $\Omega N = \frac{1}{2}\Omega M$ et $\left(\overrightarrow{\Omega M},\overrightarrow{\Omega N}\right) = \frac{\pi}{2}$ : le triangle $\Omega MN$ est rectangle en $\Omega$, avec $\Omega N = \frac{1}{2}\Omega M$.
\end{enumerate}

\paragraph{Exercice 37.}
\begin{enumerate}
\item $P(\lambda + \lambda i) = -2\lambda(\lambda + 1)(2\lambda - 3) + 2i(\lambda + 1)(2\lambda^2 - 5\lambda + 5)$ ; $2\lambda^2 - 5\lambda + 5$ n'a pas de racine r\'eelle,
      donc $\boxed{\lambda = -1}$ : $-1 - i$ est racine. (On obtient ensuite $P(z) = (z + 1 + i)\left(2z^2 - (5 + 3i)z + 5 + 5i\right)$, dont les autres racines sont
      $1 + 2i$ et $\frac{3}{2} - \frac{1}{2}i$ : ce sont les affixes de $A$, $B$, $C$.)
\item (a) $s : z' - z_A = \alpha(z - z_A)$ avec $\alpha = \frac{z_C - z_A}{z_B - z_A} = \frac{\frac{5}{2} + \frac{1}{2}i}{2 + 3i} = \frac{(5 + i)(2 - 3i)}{26} = \frac{1 - i}{2}$.
      Donc $z' = \frac{1 - i}{2}(z + 1 + i) - 1 - i$, soit $\boxed{z' = \frac{1 - i}{2}z - i}$.
      (b) Centre $A$, rapport $\left|\frac{1 - i}{2}\right| = \frac{\sqrt{2}}{2}$, angle $-\frac{\pi}{4}$.
\item (a) $|Z| = \frac{MB}{MA}$ et $\Arg{Z}$ est une mesure de l'angle $\left(\overrightarrow{MA},\overrightarrow{MB}\right)$ ($M \neq A$, $M \neq B$).
      (b) $|Z| = 1 \iff MA = MB$ : $(\mathscr{D})$ est la m\'ediatrice de $[AB]$.
      (c) $s$ conserve les rapports de distances et $s(A) = A$, $s(B) = C$ : $M$ est \'equidistant de $A$ et $B$ si et seulement si $s(M)$ est \'equidistant de $A$ et $C$.
      L'image de $(\mathscr{D})$ est la m\'ediatrice de $[AC]$.
\end{enumerate}

\paragraph{Exercice 38.}
\begin{enumerate}
\item $(2i)^2 = -4$, $(2i)^3 = -8i$, donc
      $P(2i) = -8i - (4+4i)(-4) + (2+8i)(2i) - 12i = -8i + 16 + 16i + 4i - 16 - 12i = \boxed{0}$.
\item En d\'eveloppant $(z - 2i)(az^2 + bz + c) = az^3 + (b - 2ia)z^2 + (c - 2ib)z - 2ic$ et en identifiant :
      $a = 1$, $b - 2i = -4 - 4i$ donc $b = -4 - 2i$, $-2ic = -12i$ donc $c = 6$ (et l'on v\'erifie $c - 2ib = 6 + 8i - 4 = 2 + 8i$).
\item $z^2 - (4 + 2i)z + 6 = 0$ : $\Delta = (4 + 2i)^2 - 24 = -12 + 16i = (2 + 4i)^2$, racines $\frac{4 + 2i \pm (2 + 4i)}{2}$, soit $3 + 3i$ et $1 - i$.
      $S = \{2i ; 1 - i ; 3 + 3i\}$.
\item $\frac{z_C - z_A}{z_B - z_A} = \frac{3 + i}{1 - 3i} = \frac{(3 + i)(1 + 3i)}{10} = \frac{10i}{10} = i$.
      Donc $AC = AB$ et $\left(\overrightarrow{AB},\overrightarrow{AC}\right) = \frac{\pi}{2}$ : le triangle $ABC$ est \textbf{rectangle isoc\`ele en $A$}.
\end{enumerate}

\paragraph{Exercice 39.}
\begin{enumerate}
\item Pour $x$ r\'eel, $P(x) = (x^3 - 2x^2 - 11x + 120) + i(-x^2 + 2x + 35)$.
      $-x^2 + 2x + 35 = -(x - 7)(x + 5)$ s'annule pour $x = 7$ et $x = -5$.
      Or $7^3 - 2\cdot 49 - 77 + 120 = 288 \neq 0$ et $(-5)^3 - 2\cdot 25 + 55 + 120 = 0$.
      L'unique solution r\'eelle est $\boxed{z_0 = -5}$.
\item $(z + 5)(z^2 + az + b) = z^3 + (a + 5)z^2 + (b + 5a)z + 5b$ ; en identifiant : $a = -7 - i$, $b = \frac{120 + 35i}{5} = 24 + 7i$
      (et $b + 5a = 24 + 7i - 35 - 5i = -11 + 2i$ convient).
\item $z^2 - (7 + i)z + 24 + 7i = 0$ : $\Delta = (7 + i)^2 - 96 - 28i = -48 - 14i$. On cherche $\delta = x + iy$ :
      $x^2 + y^2 = \sqrt{48^2 + 14^2} = 50$, $x^2 - y^2 = -48$, $xy = -7$ ; d'o\`u $\delta = 1 - 7i$.
      Racines $\frac{7 + i \pm (1 - 7i)}{2}$ : $4 - 3i$ et $3 + 4i$. $S = \{-5 ; 4 - 3i ; 3 + 4i\}$.
\item $|-5| = |4 - 3i| = |3 + 4i| = 5$ : les trois images sont sur le cercle de centre $O$ et de rayon $5$.
      Avec $B(3 + 4i)$ et $C(4 - 3i)$ : $\frac{z_C}{z_B} = \frac{(4 - 3i)(3 - 4i)}{25} = \frac{-25i}{25} = -i$,
      donc $OB = OC$ et $\left(\overrightarrow{OB},\overrightarrow{OC}\right) = -\frac{\pi}{2}$ : $OBC$ est rectangle isoc\`ele en $O$.
\end{enumerate}

\paragraph{Exercice 40.}
\begin{enumerate}
\item $(2+i)^2 = 3 + 4i$, $(2+i)^3 = 2 + 11i$, donc $P(2+i) = 2 + 11i - 3 - 4i + 10 + 5i - 9 - 12i = 0$.
\item En identifiant $(z - 2 - i)(z^2 + az + b) = z^3 + (a - 2 - i)z^2 + (b - (2+i)a)z - (2+i)b$ :
      $a = 1 + i$ et $b = \frac{9 + 12i}{2 + i} = \frac{(9 + 12i)(2 - i)}{5} = 6 + 3i$ (on v\'erifie $b - (2+i)(1+i) = 6 + 3i - 1 - 3i = 5$).
\item $z^2 + (1+i)z + 6 + 3i = 0$ : $\Delta = 2i - 24 - 12i = -24 - 10i = (1 - 5i)^2$, racines $\frac{-1 - i \pm (1 - 5i)}{2}$ : $-3i$ et $-1 + 2i$.
      $S = \{2 + i ; -1 + 2i ; -3i\}$, avec $z_B = -1 + 2i$ et $z_C = -3i$.
\item $i z_A = 2i + i^2 = -1 + 2i = z_B$ : $B$ est l'image de $A$ par la rotation de centre $O$ et d'angle $\frac{\pi}{2}$ ($OA = OB$ et $(OA) \perp (OB)$).
\item $ABCD$ parall\'elogramme $\iff \overrightarrow{AB} = \overrightarrow{DC} \iff z_D = z_A + z_C - z_B = 2 + i - 3i + 1 - 2i = \boxed{3 - 4i}$.
      Son centre est le milieu de $[AC]$ : $\frac{2 + i - 3i}{2} = \boxed{1 - i}$.
\end{enumerate}

\paragraph{Exercice 41.}
\begin{enumerate}
\item $z^2 - 2z + 5 = 0$ : $\Delta = -16 = (4i)^2$, $S' = \{1 - 2i ; 1 + 2i\}$.
\item $(z^2 - 2z + 5)(z^2 + az + b) = z^4 + (a - 2)z^3 + (b - 2a + 5)z^2 + (5a - 2b)z + 5b$.
      En identifiant : $a - 2 = -2 - 3i$, donc $a = -3i$ ; $5b = -10$, donc $b = -2$ ;
      on v\'erifie $b - 2a + 5 = 3 + 6i$ et $5a - 2b = 4 - 15i$.
\item $(E) \iff z^2 - 2z + 5 = 0$ ou $z^2 - 3iz - 2 = 0$. Pour la seconde, $\Delta = -9 + 8 = -1 = i^2$, racines $\frac{3i \pm i}{2}$ : $2i$ et $i$.
      $S = \{1 - 2i ; 1 + 2i ; i ; 2i\}$.
\item La somme vaut $2 + 3i$, oppos\'e du coefficient $-(2 + 3i)$ de $z^3$ (comme pour le second degr\'e, o\`u la somme des racines vaut $-\frac{b}{a}$).
\end{enumerate}

\paragraph{Exercice 42.}
\begin{enumerate}
\item $(E') : Z^2 + (2 - 16i)Z - 63 + 16i = 0$. $\Delta = (2 - 16i)^2 + 4(63 - 16i) = 4 - 64i - 256 + 252 - 64i = -128i$,
      et $(8 - 8i)^2 = 64 - 128i - 64 = -128i$. Donc $Z = \frac{-2 + 16i \pm (8 - 8i)}{2}$ : $Z_1 = 3 + 4i$ et $Z_2 = -5 + 12i$.
\item $\delta = x + iy$ avec $\delta^2 = 3 + 4i$ : $x^2 + y^2 = 5$, $x^2 - y^2 = 3$, $xy = 2$ : $\delta = \pm(2 + i)$.
      $\delta^2 = -5 + 12i$ : $x^2 + y^2 = 13$, $x^2 - y^2 = -5$, $xy = 6$ : $\delta = \pm(2 + 3i)$.
\item $S = \{2 + i ; -2 - i ; 2 + 3i ; -2 - 3i\}$.
\item Les solutions sont deux \`a deux oppos\'ees : les segments joignant $2 + i$ \`a $-2 - i$ et $2 + 3i$ \`a $-2 - 3i$ ont le m\^eme milieu $O$.
      Le quadrilat\`ere de sommets successifs $2 + i$, $2 + 3i$, $-2 - i$, $-2 - 3i$ a ses diagonales de m\^eme milieu : c'est un parall\'elogramme de centre $O$.
\end{enumerate}

\paragraph{Exercice 43.}
\begin{enumerate}
\item $f$ est d\'efinie sur $\mathbb{C} \setminus \{1\}$.
\item (a) $a = f(2) = \frac{2 - 2i}{-1} = -2 + 2i = 2\sqrt{2}\left(\cos\frac{3\pi}{4} + i\sin\frac{3\pi}{4}\right)$.
      (b) $a^k = (2\sqrt{2})^k\expo{i\frac{3k\pi}{4}}$ est r\'eel $\iff \frac{3k\pi}{4} \equiv 0\ [\pi] \iff 4 \mid 3k \iff 4 \mid k$ :
      $a^k$ est r\'eel pour $k$ multiple de $4$.
\item (a) $z'(1 - z) = 2 - iz \iff z(i - z') = 2 - z'$. Comme $z' = i$ conduirait \`a $0 = 2 - i$, on a $z' \neq i$ et $z = \frac{z' - 2}{z' - i}$.
      (b) $OM = |z| = \frac{|z' - 2|}{|z' - i|} = \frac{M'A}{M'B}$ avec $A(2)$ et $B(i)$.
      (c) Si $|z| = 1$ alors $M'A = M'B$ : $M'$ appartient \`a la m\'ediatrice $(\mathscr{D})$ du segment $[AB]$.
      (d) $(\mathscr{C})$ est le cercle de centre $O$ et de rayon $1$ priv\'e du point d'affixe $1$ ; $(\mathscr{D})$ est la m\'ediatrice de $[AB]$,
      passant par le milieu $\left(1, \frac{1}{2}\right)$ et perpendiculaire \`a $(AB)$ (droite d'\'equation $y = 2x - \frac{3}{2}$).
\item (a) $\frac{\sqrt{3} + i}{2} = \expo{i\frac{\pi}{6}}$ est de module $1$ et diff\'erent de $1$ : $S$ est une rotation d'angle $\frac{\pi}{6}$.
      (b) $z = \expo{i\frac{\pi}{6}}z + 1 + 4i \iff z = \frac{1 + 4i}{1 - \expo{i\frac{\pi}{6}}}$ : un seul point invariant $\Omega$.
      Avec $1 - \expo{i\frac{\pi}{6}} = \frac{2 - \sqrt{3}}{2} - \frac{1}{2}i$, le calcul donne
      $\omega = -\frac{7}{2} - 2\sqrt{3} + i\left(3 + \frac{\sqrt{3}}{2}\right)$.
      (c) $S$ est la rotation de centre $\Omega$ et d'angle $\frac{\pi}{6}$.
\end{enumerate}

\end{corrige}

